\documentclass[11pt,a4paper,oneside]{report}
\usepackage[margin=2.5cm]{geometry}
\usepackage{amsmath,amssymb,amsthm}
\usepackage{booktabs,longtable,array}
\usepackage{enumitem}
\usepackage{float}
\usepackage{xcolor}
\usepackage{caption}
\usepackage{tikz}
\usetikzlibrary{positioning,arrows.meta}
\usepackage{pgfplots}
\pgfplotsset{compat=1.18}
\usepackage[hidelinks]{hyperref}
\newtheorem{theorem}{Theorem}[chapter]
\newtheorem{proposition}[theorem]{Proposition}
\newtheorem{lemma}[theorem]{Lemma}
\newtheorem{corollary}[theorem]{Corollary}
\theoremstyle{definition}
\newtheorem{definition}[theorem]{Definition}
\newtheorem{remark}[theorem]{Remark}
\newtheorem{heuristic}[theorem]{Heuristic}
\newtheorem{observation}[theorem]{Observation}
\newtheorem{example}[theorem]{Example}
\newcommand{\Z}{\mathbb{Z}}
\newcommand{\F}{\mathbb{F}}
\newcommand{\Cm}{\mathfrak{Cm}}
\newcommand{\RA}{\mathrm{RA}}
\newcommand{\RRA}{\mathrm{RRA}}
\newcommand{\At}{\mathrm{At}}
\newcommand{\score}{\operatorname{score}}
\newcommand{\sgn}{\operatorname{sign}}
\newcommand{\dom}{\operatorname{Dom}}
\newcommand{\vtwo}{v_2}
\newcommand{\Fix}{\operatorname{Fix}}
\newcommand{\code}[1]{\texttt{#1}}
\newcommand{\ALO}{\operatorname{ALO}}
\newcommand{\AMO}{\operatorname{AMO}}
\newcommand{\EO}{\operatorname{EO}}
\newcommand{\st}[1]{\textsf{#1}}
\renewcommand{\chaptername}{Paper}
\setcounter{tocdepth}{1}
\setlength{\emergencystretch}{3em}
\newcommand{\paperauthor}{\begin{center}\large Charles Gretton\footnote{Corresponding author: \href{mailto:charles.gretton@anu.edu.au}{charles.gretton@anu.edu.au}.}\\[2pt]\normalsize Australian National University\end{center}\vspace{0.5em}}
\newcommand{\paperabstract}[1]{\begin{quote}\small\noindent\textsc{Abstract.} #1\end{quote}\vspace{0.5em}}
\makeatletter
\newenvironment{paperrefs}{\section*{References}\list{\@biblabel{\@arabic\c@enumiv}}{\settowidth\labelwidth{\@biblabel{99}}\leftmargin\labelwidth\advance\leftmargin\labelsep\usecounter{enumiv}\let\p@enumiv\@empty\renewcommand\theenumiv{\@arabic\c@enumiv}}\sloppy\small}{\endlist}
\makeatother

\begin{document}
\begin{titlepage}
\centering
\vspace*{3cm}
{\LARGE Costas Arrays of Order 32\par}
\vspace{0.6cm}
{\Large Virtual Proceedings of an Investigation\par}
\vspace{2cm}
{\large Charles Gretton\textsuperscript{$\ast$}\par}
{\normalsize Australian National University\par}
\vspace{1cm}
{\normalsize October 2026\par}
\vfill
{\small\textsuperscript{$\ast$}Corresponding author:
\href{mailto:charles.gretton@anu.edu.au}{charles.gretton@anu.edu.au}\par}
\end{titlepage}

\chapter*{Licence and errata}
\addcontentsline{toc}{chapter}{Licence and errata}

\paragraph{Licence.} \textcopyright\ 2026 Charles Gretton. This work is licensed under the Creative Commons Attribution 4.0 International licence (CC BY 4.0).
To view a copy of the licence, visit \url{https://creativecommons.org/licenses/by/4.0/}. The software released
with it in the Costas-32 repository is licensed under the Apache License, Version 2.0
(\url{https://www.apache.org/licenses/LICENSE-2.0}). Material by others that is cited or reproduced here,
including published data and catalogues, remains under its original terms.

\paragraph{Errata.} We welcome reports of errors, omissions and overlooked prior work. Please send them to the
corresponding author, Charles Gretton, at \href{mailto:charles.gretton@anu.edu.au}{charles.gretton@anu.edu.au}.
Every report that leads to a correction will be acknowledged in the errata and in any revised edition,
cited as a private communication; where the point was already made in a publication that appeared before
this report's publication date, that publication will be cited as well.

\chapter*{AI usage acknowledgment}
\addcontentsline{toc}{chapter}{AI usage acknowledgment}

Large language models were used extensively throughout this investigation and in preparing this
volume. Claude Opus 5.5 and GPT 5.6 SOL were used for the report and for the proceedings as a
whole. Their use covered developing and checking mathematical arguments; designing, writing,
testing and running the software; planning, monitoring and analysing the computational campaigns;
searching and assessing the literature; and drafting, reviewing and revising the text, including
assembling these proceedings from the project's reports.

The author directed the work, decided what to run and what to claim, and checked its major
claims. He has not made a complete inventory of the claims in this volume, nor checked each one
individually; doing so would take far longer than the investigation itself, so errors may remain,
and reports of them are welcome (see the errata above). What the volume provides instead is the
means to check: each computational claim rests on released software and recorded outputs, and the
report states how its claims were checked. The use of models also shaped the methods: much of the
emphasis on fail-closed accounting, independent verification of every answer and hypotheses stated
before the runs, described in the note on software engineering, exists so that results produced with heavy machine assistance can be audited without
trusting the process that produced them.

The author calls a release of this kind a \emph{virtual proceedings}: a self-contained manuscript
published together with the code package that produced its results --- the volume and its source,
the software needed to rerun every experiment, and the derived evidence behind every reported
number --- as one open repository rather than through a conventional venue. The name reflects both
how such work is now produced, largely in dialogue with AI tools under the author's direction, and
how it is meant to be used: read, rerun and audited by others.

\chapter*{Preface}
\addcontentsline{toc}{chapter}{Preface}

Orders 32 and 33 are the smallest orders for which the existence of a Costas array is unknown. This
volume collects the results of a systematic investigation of order 32 as ten self-contained
papers, one for each strand of the work. Each paper is readable on its own at the level of its main claim. Shared notation, recurring definitions and release conventions are collected once, in the front matter, and each paper states its dependencies on the others.

The investigation does not resolve existence at order 32. What it does is close a series of
structured routes by which an order-32 array could arise, and measure or explain why others fail.
The table summarises the ten papers and the standing of their results.

\begin{center}
\small
\begin{tabular}{@{}p{0.04\textwidth}>{\raggedright\arraybackslash}p{0.42\textwidth}>{\raggedright\arraybackslash}p{0.44\textwidth}@{}}
\toprule
& result & standing \\
\midrule
1 & No known construction reaches order 32; no circular Costas array of order 32 exists. &
Mostly from the literature; single periodicity remains open pending the Golomb--Taylor conjecture. \\
2 & No order-32 array is a two-point augmentation of the 664 generated algebraic order-30 cores;
the closest miss has a complete genealogy. & Computational; relative to the generated order-30
catalogue. \\
3 & No order-32 array contains an order-29 or an order-28 Costas deletion minor. &
Computational; quantifies over complete published censuses; neither result is known to imply the other. \\
4 & The order-29 census, reconstructed and independently verified. & Verification, partly
kernel-checked in Lean; completeness rests on the published enumeration. \\
5 & A telescoping-sum propagation rule, proved sound. & Passes a frozen benchmark; slows exact search
in practice. \\
6 & A first-moment model forecasting about $10^{-4}$ sporadic classes at order 32. & Conditional
forecast; an equally good fit diverges. \\
7 & At least 112 of the 240 diagonal second differences of any order-32 permutation are even. &
Proved; does not bear on existence. \\
8 & No transpose-symmetric Costas array of order 32 exists. & Known since 1993; determined
independently here twice, by SAT and by exhaustive direct search. \\
9 & An exhaustive search of all order-32 permutations, by a straightforward method, would take about $1.5\times10^6$ core-years as first written, about $1.9\times10^5$ once engineered. & Measured, with validated estimators; the faster 2011 method would need an estimated $10^4$ core-years. \\
10 & Late dead prefixes are refuted by a pair-marginal linear relaxation; nothing tested refutes earlier prefixes more cheaply than search. & Sampled and independently replicated; caps any rule acting at the transition at about nine-fold at order 32. \\
\bottomrule
\end{tabular}
\end{center}

Two cautions apply throughout. First, most results here are computational: they rest on the
correctness of the software, which is released, tested and in places independently re-implemented,
but not formally certified except where a paper says so. Each paper ends with a section stating
the standing of its claims. Second, the result of Paper~8 turned out to have been published in 1993
and confirmed in 2009, which the project discovered only after its own campaign was complete. The
other strands have not yet been checked against the literature as thoroughly, and their novelty
should be read with that in mind.

Some papers point to the project's consolidated report for material outside their scope, such as
structural approaches that were tried and refuted.

\paragraph{Data and code availability.} The consolidated report, all software, Lean sources,
catalogues, logs and result ledgers referred to in this volume are released together as one software
archive (release \texttt{costas32-1.3}), whose manifest records the hash of every file. A persistent
identifier for the archive is to be assigned on deposit.

\paragraph{Dependencies and standing.} The table records, for each paper, the kind of result, what it
depends on, and what it takes from outside the project.

\begin{center}
\small
\begin{tabular}{@{}p{0.04\textwidth}>{\raggedright\arraybackslash}p{0.21\textwidth}>{\raggedright\arraybackslash}p{0.23\textwidth}>{\raggedright\arraybackslash}p{0.23\textwidth}>{\raggedright\arraybackslash}p{0.19\textwidth}@{}}
\toprule
& result type & depends on & external premise & standing \\
\midrule
1 & background and analysis & shared notation & published theorems on constructions and periodicity & cited or derived \\
2 & catalogue-relative exclusion & catalogue generator & none; order 30 is unenumerated & computational, uncertified \\
3 & census-complete exclusions & rung engine; Paper~4 & completeness of the order-28 and order-29 censuses & computational, uncertified \\
4 & census verification & published representatives & completeness of the order-29 census & partly kernel-checked \\
5 & inference rule and evaluation & none & none & proved rule; measured performance \\
6 & conditional model & Paper~1; published counts & the model form & heuristic \\
7 & theorem and bounds & none & none & proved; computational \\
8 & independent confirmation & regional encoding, solvers, audit & none; the 1993 and 2009 counts are for comparison & computational, uncertified \\
9 & cost measurement & Papers~1 and~8 & the reported cost of the 2011 enumeration & estimated, validated \\
10 & obstruction search & Papers~2 and~9 & the complete catalogues of orders 12--15 & sampled, replicated \\
\bottomrule
\end{tabular}
\end{center}

What remains open is the residual problem, stated more fully in the closing chapter: order-32 arrays that are not algebraic, not periodic,
not geometrically symmetric, contain no Costas deletion minor of order 28 or 29, and are not
augmentations of algebraic order-30 cores. No structural classification in this volume controls
that class, and exhaustive enumeration of all order-32 permutations remains far out of reach.


\chapter*{A note on software engineering}
\addcontentsline{toc}{chapter}{A note on software engineering}

The investigation was run as a software project as much as a mathematical one, and several terms
from that practice recur in the papers. Two of them, \emph{rungs} and \emph{gates}, carry technical
weight and are introduced here.

\paragraph{Rungs.} Several searches are organised as a ladder of finite exhaustive searches indexed
by a parameter. Each \emph{rung} is one complete run, with its own declared coverage and its own
verdict. The main instance is the augmentation ladder of Papers~2 and~3. Rung $k$ searches every
permutation of order $n$ obtained from a Costas core of order $n-k$ by inserting $k$ columns, $k$
rows and $k$ dots. Rung~1 runs from order 30 to order 31; rungs 2, 3 and 4 target order 32 from
cores of orders 30, 29 and 28. Rungs are run in
increasing cost, and each is self-contained: a negative result at one rung is a theorem about that
rung alone. Rungs need not imply one another; Paper~3 explains why neither of rungs 3 and 4 is known to imply the other.

\paragraph{Gates.} A \emph{gate} is an acceptance criterion written down and frozen before the run
it judges, in the manner of a pre-registered experiment or a release gate in continuous
integration. Freezing it prevents the criterion from being tuned to the result. Paper~5 evaluates a
propagation rule in two stages: Stage~1 on a benchmark of partial assignments, and Stage~2 inside an
exact search. It uses two kinds of gate. A \emph{scientific gate} asks whether the rule deduces
more than the baseline without ever rejecting a completable assignment. An \emph{operational gate}
asks whether it makes search faster. A result can pass one and fail the other, and Paper~5's does. (Paper~8 uses ``Stage~1'' independently, for the first stage of its own campaign.)

\paragraph{Related practice.} Other recurring terms are as follows.
\begin{itemize}
\item \emph{Fail-closed accounting.} A run counts as complete only when its declared coverage is known to cover the whole search space, by a proof such as Paper~8's partition and cover theorems or by construction, as with Paper~3's one work unit per dihedral orbit, and every unit of that coverage reports an explicit terminal status. Anything missing, failed or unparsed counts as not done.
\item \emph{Census match.} Before a census-complete search runs, its catalogue is checked against
the published census figures and its hash recorded.
\item \emph{Work units and independent merging.} Large searches are split into deterministic units
whose outputs are merged and checked by a separate tool.
\item \emph{Independent re-scoring.} Every reported near-solution is re-scored by a separately
written scorer.
\item \emph{Release gate.} A repository self-test that must pass before any archive is released.
\end{itemize}

\chapter*{Notation and glossary}
\addcontentsline{toc}{chapter}{Notation and glossary}

Papers 1--7 write a permutation as a sequence of heights $h$, with the dot of column $i$ in row
$h_i$; Paper~8 writes it as $\pi$. All indices are zero-based.

\section*{Notation}
\begin{longtable}{@{}l >{\raggedright\arraybackslash}p{0.62\textwidth}@{}}
\toprule
Symbol & Meaning \\
\midrule
\endhead

$n$ & Order of a Costas array; the board is $n\times n$. Usually $n=32$. \\
$h$ & A permutation $h:\{0,\dots,n-1\}\to\{0,\dots,n-1\}$, written zero-based.
      $h_i$ is the \emph{height} (row) of the dot in column $i$. \\
$\Delta_{m,i}$ & $h_{i+m}-h_i$, the vertical difference at horizontal
      separation $m$ starting at column $i$. \\
$m$, $d$ & Horizontal separation (distance) between two columns. $d$ is used in
      solver contexts, $m$ in partition-counting contexts; they denote the same
      quantity. \\
$e$ & A signed vertical difference; $e\in\mathcal V$. \\
$g$ & An absolute vertical distance, $g=|e|$. \\
$\mathcal V$ & The signed palette $\{-(n-1),\dots,-1,1,\dots,n-1\}$,
      $|\mathcal V| = 2n-2$; for $n=32$, $|\mathcal V|=62$. \\
$L_m$ & $n-m$, the number of column pairs at separation $m$. \\
$\score(h)$ & Costas defect: the number of repeated $(d,\Delta)$ pairs,
      $\sum_{d,e}\max(0,\chi_{d,e}-1)$, over $1\le d\le n-1$ and nonzero $e$ with $|e|\le n-1$ (differences of a permutation are never zero). $\score(h)=0$ iff $h$ is Costas. \\
$\chi_{m,e}$ & $\#\{i: \Delta_{m,i}=e\}$, the occupancy of signed value $e$
      at separation $m$. \\
$y_{m,g}$ & $\chi_{m,g}+\chi_{m,-g}\in\{0,1,2\}$, magnitude occupancy. \\
$T_m$ & $\sum_i \Delta_{m,i}$, the telescoping row sum at separation $m$. \\
$N_q(r)$ & $\#\{e\in\mathcal V : e\equiv r \pmod q\}$, the signed residue
      capacity. \\
$\vtwo(x)$ & The 2-adic valuation of $x$; $\vtwo(0)=\infty$. \\
$D_i$ & Domain of the height variable $h_i$ during propagation. \\
$E_{i,j}$, $D_{ij}$ & Domain of the edge (difference) variable $c_{ij}=h_j-h_i$. \\
$k$ & Number of columns and rows deleted in an augmentation rung; the core has
      order $n-k$. \\
$K$ & A Costas core array of order $n-k$. \\
$C$, $R$ & The deleted column set and deleted row set, each of size $k$. \\
$\delta,\epsilon$ & Compressed displacement of a core pair: column distance and
      row difference \emph{within the core}. \\
$D_4$ & The dihedral group of order 8 acting on the square; the Costas
      symmetry group. \\
$W_1(p,g,c)$ & Welch construction, order $p-1$: $h_i = g^{i+c}\bmod p - 1$. \\
$W_2,W_3$ & $W_1$ with one or two corner dots deleted, orders $p-2$, $p-3$. \\
$G_2(q,\alpha,\beta)$ & Golomb construction over $\F_q$, order $q-2$:
      dot at $(i,j)$ iff $\alpha^i+\beta^j=1$. \\
$G_3,G_4,G_5$ & $G_2$ with corner dots deleted, orders $q-3,q-4,q-5$. \\
$C_{31}$ & The unique order-31 Costas array reachable by one-dot augmentation
      of the algebraic order-30 cores (one $D_4$ orbit). \\
$N_{32}$ & The known order-32 permutation of defect 2; the corner extension of
      $C_{31}$. \\
$\pi$ & In Paper~8, a permutation, written $\pi$ rather than $h$; transpose symmetry is
      $\pi=\pi^{-1}$. \\
$\Fix(\pi)$, $f$ & The fixed points of an involution $\pi$, and their number. \\
$F$, $R_F$ & A fixed-point set, and the region of Costas involutions whose fixed-point set is
      exactly $F$ (Paper~8). \\
$x_{ab}$, $h_{ij}$, $p_{s,x,\delta}$ & The edge, image and displacement-indicator variables of
      Paper~8's encoding of a region. \\
$Z(\pi)$ & In Paper~8, the number of large-large pairs: pairs $\{x,y\}$ with $|x-y|\ge16$ and
      $|\pi(x)-\pi(y)|\ge16$. \\
$k$ (Paper~8) & The number of split coordinates in a native split. In Papers~2 and~3, $k$ is the
      rung parameter instead. \\
$M(D,k)$ & In Paper~8, the number of messages of a native split with $k$ split coordinates and $D$
      free partners. \\
\bottomrule
\end{longtable}

\section*{Glossary}
\begin{description}[style=nextline,leftmargin=0pt,itemsep=2pt]
\item[Costas array] A permutation array of order $n$ in which all $\binom{n}{2}$ vectors connecting
pairs of dots are distinct. Equivalently, for every horizontal separation $m$, the vertical
differences $h_{i+m}-h_i$ are pairwise distinct.
\item[Displacement] The vector $(j-i,\,h_j-h_i)$ between the dots of columns $i<j$.
\item[Defect, score] The number of repeated displacement vectors, $\score(h)$. A Costas array has
defect 0; the known order-32 near-solution $N_{32}$ has defect 2.
\item[Dihedral orbit, equivalence class] The set of arrays obtained from one array under the eight
symmetries of the square, $D_4$. Censuses count arrays both individually and by class.
\item[Census] A complete list of the Costas arrays of one order, obtained by exhaustive enumeration.
Complete censuses exist through order 29.
\item[Census-complete, catalogue-relative] A result is census-complete if it quantifies over a
complete census, and catalogue-relative if it quantifies only over a generated catalogue, such as
the 664 order-30 arrays of Paper~2, at an order that has not been enumerated.
\item[Welch and Golomb constructions] The two principal algebraic constructions of Costas arrays,
over prime fields and finite fields respectively; see the notation table.
\item[Circular and singly periodic sequences] A Costas sequence $a_0,\dots,a_{n-1}$ with values in
$\{1,\dots,n\}$ is \emph{circular} if, for every lag $k$, the circular differences
$a_{(i+k)\bmod n}-a_i$ are distinct modulo $n+1$: its dots, read in $\Z_n\times\Z_{n+1}$, have
distinct displacements. It is \emph{singly periodic} if every cyclic shift of its terms is again
Costas. Circularity lies between single periodicity and Moreno's shifting property. Circular
sequences exist only when $n+1$ is prime (Golomb and Moreno, 1996) and are exactly the Welch
sequences (Muratovi\'c-Ribi\'c, Pott, Thomson and Wang, 2015); every circular sequence is an
ordinary Costas array, but not conversely. Arrays required to stay Costas when extended periodically
in both directions are a further notion, studied separately. See Paper~1.
\item[Corner deletion, corner addition] Removing or adding a dot at a corner of the board, together
with its row and column. The compression is then a translation, so these are the deletions that
preserve the Costas property in general.
\item[Sporadic array] A Costas array not accounted for by a recognised generated or emergent
construction. The category shrinks whenever a new construction is found. Order 28 is the first
order above 5 with no sporadic arrays, and order 29 has none either.
\item[Emergent array] An array produced by a semi-empirical construction rather than an algebraic
family, such as the 16 order-29 arrays produced by the Rickard--Welch operation.
\item[Augmentation rung] The exhaustive search over all permutations of order $n$ formed by placing
a Costas core of order $n-k$ in a compressed sub-board and adding $k$ dots; see the note on software
engineering.
\item[Deletion minor] The permutation array left after deleting $k$ columns, $k$ rows and the $k$
dots they share from a larger permutation array, and closing up the gaps: the inverse of a rung-$k$
augmentation.
\item[Core-first pruning] In a rung search, testing the placed core's own displacements before
considering the added dots. Exact for defect 0, since a repeated displacement among a subset of the
dots is a repeated displacement of the whole.
\item[Involution, fixed point, transposition] A permutation equal to its inverse; a coordinate $c$
with $\pi(c)=c$, whose dot lies on the main diagonal; a pair $(a\,b)$ with $\pi(a)=b$ and $\pi(b)=a$.
\item[Geometric and transpose symmetry] An array is geometrically symmetric if a non-identity motion
of the square fixes it, and transpose-symmetric if the main-diagonal reflection does. For orders of
at least 3 the two notions agree up to the choice of diagonal (Paper~8). The literature calls
transpose-symmetric arrays \emph{diagonal}.
\item[Anti-reflective and consecutive arrays] Symmetry types studied in the literature that are not
motions of the square. They are mentioned for comparison only.
\item[Levels of assurance] A \emph{validated} computation agrees with independent implementations,
regression tests or known answers; an \emph{audited} computation has its coverage and execution
records independently checked; a \emph{machine-checked witness} is an object, such as a Costas
permutation or a Lean theorem, checked independently; a \emph{proof-certified} unsatisfiable verdict
has a checked DRAT or LRAT proof; an \emph{uncertified} one rests on the solver's exhaustion.
\item[Region] In Paper~8, the set of Costas involutions with a given fixed-point set.
\item[Native split, message, unsplit branch] Paper~8's distributed decomposition fixes the partners
of $k$ chosen coordinates; each assignment of partners is a message, and a message in which two
chosen coordinates are paired with each other is an unsplit branch.
\item[Seed-refuted] In Paper~8, a pairing that repeats a displacement with the fixed dots alone.
\item[Case] In Paper~8, a region with one transposition fixed.
\item[Gate] An acceptance criterion frozen before the run it judges; see the note on software
engineering.
\item[CDCL, DRAT and LRAT] Conflict-driven clause learning, the algorithm of modern SAT solvers; and
machine-checkable proof formats for unsatisfiability. Without such a proof, an unsatisfiable verdict
is a computational observation rather than a certificate.
\item[Verdicts] Paper~8 records \st{SAT-VALIDATED} (a model that passed independent checking),
\st{UNSAT-UNCERTIFIED} (normal solver exhaustion, no checked proof), \st{TIMEOUT} and \st{ERROR}. A
region closed by an \st{UNSAT-UNCERTIFIED} verdict is \emph{search-closed}.
\end{description}


\tableofcontents

\chapter{The Landscape at Order 32: Constructions, Periodicity and the Cost of Enumeration}\label{p1-paper}
\paperauthor
\paperabstract{Orders 32 and 33 are the smallest orders at which the existence of a Costas array is unknown.
This paper sets out the landscape at order 32. No known algebraic construction reaches order 32,
and we determine which construction families reach which orders near it. No circular Costas array
of order 32 exists, because a circular Costas sequence exists only when $n+1$ is prime and 33 is
composite; whether a singly periodic array can exist remains tied to the Golomb--Taylor conjecture.
We also record the measured growth law of exact enumeration, a constraint that any non-existence
proof must satisfy, and a three-branch decomposition of the existence question that organises the
other papers in this volume.}

\section{Introduction}\label{p1-sec:paperintro}
Costas arrays are known at every order up to 31, mostly through two algebraic constructions
due to Welch and to Golomb~\cite{p1-golombtaylor}, together with arrays found by exhaustive enumeration, which is complete through order 29~\cite{p1-drakakis27,p1-drakakis28,p1-drakakis29}; the counts are tabulated in~\cite{p1-oeis}. At orders 32 and 33 neither route has produced an array. This paper explains why,
in terms of what each route can reach, and fixes the vocabulary and decomposition used in the rest
of this volume. Most of its content is drawn from the literature; the contribution is its
organisation, the measurements of enumeration cost, and the analysis of which construction
families reach which orders.

\paragraph{Dependencies.} This paper uses only the shared notation; the other papers draw on its account of the constructions.

\section{Preliminaries}\label{p1-sec:notation}
A \emph{Costas array} of order $n$ is an $n\times n$ permutation array in which the displacement
vectors between pairs of dots are all distinct. Writing the array as a permutation $h$ of
$\{0,\dots,n-1\}$, with the dot of column $i$ in row $h_i$, the condition is that for each
horizontal separation $m$ the vertical differences $h_{i+m}-h_i$ are distinct. Costas arrays are
known at every order up to 31; orders 32 and 33 are the smallest for which existence is open. The
dihedral group $D_4$ of the square maps Costas arrays to Costas arrays, and arrays are counted
either individually or by $D_4$ orbit (equivalence class). Further notation and terms are collected
in the notation table and glossary at the front of this volume.


\section{Status of the existence question}

Costas arrays are known to exist for all orders $n \le 31$. Orders 32 and 33 are
the smallest for which none is known. Exhaustive enumeration has been carried
out for all orders up to 29; orders 30 and 31 have never been exhaustively
enumerated, although arrays are known at both.

Reported enumeration costs (cited) are: order 26, $29.52$ single-processor
years; order 27, approximately 30 CPU-years across 120 machines over three
months; order 28, about 70 CPU-years; order 29, $366.55$ CPU-years. The observed
growth of the search is slightly worse than a factor of 5 per increment of $n$.

\subsection{Recent work, 2022--2026}
\label{p1-sec:recent}
The enumeration frontier has not moved: Gulec and Abolghasemi, proposing a framework of ``universal
Costas matrices'' for reconstruction-based search, note that all Costas arrays are known only up to
order 29~\cite{p1-gulec2026ucm}. Two structural results bear on searches at higher orders. Warnke, Correll
and Swanson proved that the proportion of permutations of order $n$ that are Costas arrays decays
exponentially, $C(n)/n!\le e^{-cn}$~\cite{p1-warnke2023density}. Correll and Swanson showed that, unless
undiscovered sporadic arrays exist at high orders, the number of positive first differences of an
order-$n$ Costas array lies within 2.5 of $(n-1)/2$, and that this property alone yields only an
$O(n^{-1/2})$ density among permutations~\cite{p1-correll2023difference}.

At order 32 itself there are two near-solutions and two restricted searches. An ant-colony search reported
an order-32 permutation with 17 violations of the Costas property after more than 10{,}000
CPU-days~\cite{p1-vulakh2022mdraco}. In August 2026 Rosi\'c publicly documented a permutation with only two
repeated displacement vectors, the array $N_{32}$ of Paper~2, with a proof checked in DRAT format that no
Costas array lies within Hamming distance 12 of it and a program enumeration of distances 13 to 16 that
found none; the eight such balls around its symmetry orbit exclude about $3.3\times10^{17}$ permutations,
some $10^{-18}$ of all, and are expressly not a non-existence result~\cite{p1-rosic2026gist}. The proof
package is held outside the published record and was not rechecked here. In a preprint, Dumitru and
Budala restrict the search at a new order to prefixes whose seven-dot difference signatures occur at
smaller orders: at order 28 such a vocabulary built from order 27 recovers 67 of the 89 known classes;
restricted sweeps at orders 30 and 31 return exactly the known arrays; and at order 32 the restricted
class is empty, after 102 GPU-hours. They stress that an order-32 array outside the class would simply be
sporadic, so this is not evidence against existence~\cite{p1-dumitru2026localpatterns}. In a second preprint
they report an exhaustive census of transpose-symmetric arrays of orders 37 to 42, finding 4, 0, 16, 2, 12
and 4 of them, and no sporadic symmetric array at any order from 24 to 42~\cite{p1-dumitru2026symmetriccensus}.

\section{Why orders 32 and 33 specifically}

The algebraic constructions produce arrays of orders $p-1$, $p-2$, $p-3$ for $p$
prime and $q-2$, $q-3$, and conditionally $q-4$ and $q-5$, for $q$ a prime
power. The orders reachable near 32 are therefore determined by the prime powers
in the window:

\begin{center}
\begin{tabular}{@{}lccccccccc@{}}
\toprule
order & 28 & 29 & 30 & \textbf{31} & \textbf{32} & \textbf{33} & 34 & 35 & 36 \\
\midrule
direct construction & $29{-}1$ & $31{-}2$ & $31{-}1$ & --- & --- & --- &
  $37{-}3$ & $37{-}2$ & $37{-}1$ \\
dots from nearest seed & 0 & 0 & 0 & 1 & 2 & 2 & 0 & 0 & 0 \\
\bottomrule
\end{tabular}
\end{center}

Since $33 = 3\cdot 11$, $34 = 2\cdot 17$, $35 = 5\cdot 7$ and $36 = 6^2$ are all
non-prime-powers, the only prime power in the window above 31 is 37, which
reaches down to 34 unconditionally and to 33 and 32 only through the conditional
$q-4$ and $q-5$ constructions, whose conditions fail. Orders 31, 32 and 33 are
therefore the only orders in $[28,36]$ with no direct construction, at distances
1, 2 and 2 from the nearest seed. Order 31 succeeds at distance 1; orders 32 and
33 do not at distance 2.

\begin{proposition}[Orders of the unmodified standard constructions]
\label{p1-prop:polydual}
For the standard Welch and Golomb parameterisations to induce an order-$N$
permutation without deleting field elements or selecting a proper output window,
one requires
\[
q-1=N \qquad\text{or}\qquad q-2=N,
\]
respectively. Equivalently, these are precisely the cases in which the standard
parameter domain has the required cardinality automatically.

For $N=32$ the required field orders are therefore $q=33$ and $q=34$. Neither is
a prime power, so neither finite field exists. Consequently there is no automatically loss-free Welch or Golomb parameterisation at order 32.
\end{proposition}

\noindent The proposition concerns only the unmodified constructions. It says nothing about deletion
variants, augmentation, subfield restrictions, other polynomial constructions or emergent
transformations, which this paper and Papers~2 and~3 treat separately.

The next available field order is $q=37$. A construction over this field must
therefore impose additional structure restricting its values to an appropriate
width-32 window. Under the heuristic model in which the 32 outputs are
independently uniform over the 36 available nonzero field values, the probability
of filling one prescribed window injectively is
\[
P_{\mathrm{window}} \;=\; \frac{32!}{36^{32}} \;=\; 4.1542\times10^{-15}.
\]
This factor is \emph{not} a non-existence theorem for structured polynomial
families. It quantifies why uninformed sampling over the next available field is
ineffective.

The automatic loss-free condition is the polynomial dual of the prime-power gap
visible in the construction table above. The associated window factor is a
separate random-output heuristic explaining why unstructured sampling beyond that
gap is ineffective. The computation is recorded in the project's consolidated report.

\begin{remark}
Nothing about $32 = 2^5$ plays any role. The arithmetic that distinguishes order
32 is the distribution of prime powers in $\{33,\dots,37\}$. This is worth
stating explicitly because the binary structure of 32 motivated several of the
approaches refuted in the project's consolidated report, and in every case the
binary structure turned out to be irrelevant.
\end{remark}

\section{The constraint any non-existence proof must satisfy}

Costas arrays exist at orders 31 and 36. If none exists at 32, the counting
function is non-monotonic in a way tied to prime-power gaps. Consequently:

\begin{observation}[Order-sensitivity requirement]
\label{p1-obs:uniformity}
A successful non-existence argument must contain some feature capable of
distinguishing order 32 from nearby orders at which Costas arrays exist, notably
31 and 36. Bounds or relaxations whose hypotheses and numerical strength vary
smoothly and monotonically across this range are unlikely to suffice unless they
include an arithmetic, boundary, or discrete threshold effect.
\end{observation}

This is a research heuristic, not an impossibility theorem. A theorem may be
uniform in $n$ while its numerical or arithmetic conclusion varies with $n$: a
uniform bound $C(n)\le F(n)$ could in principle satisfy $F(31)\ge1$, $F(32)<1$,
$F(36)\ge1$, and Bruck--Ryser is exactly a uniform theorem whose arithmetic
hypothesis distinguishes orders. The existence of arrays at 31 and 36 therefore
does not rule out a general theorem implying non-existence at 32. What the
measurements below establish is that each \emph{particular} relaxation tested
returns the same verdict at 31, 32 and 36.

This observation provides a diagnostic for the particular relaxations examined in the project's consolidated report, not an impossibility theorem. In each case
tested, the relaxation lacked sufficient order-sensitive strength and returned
essentially the same verdict at orders 31, 32 and 36. A stronger density theorem
or arithmetic bound capable of distinguishing these orders remains possible.

\section{The three-branch decomposition}

Any proof of non-existence must dispose of three classes.

\begin{enumerate}[label=(\roman*)]
\item \textbf{Periodic arrays.} Golomb and Moreno proved in 1996 that an order-$n$ circular Costas sequence exists only when $n+1$ is prime~\cite{p1-golombmoreno}. Since 33 is composite,
\emph{no circular Costas array of order 32 exists}; this branch is closed. The full Golomb--Moreno conjecture, that every circular Costas array is Welch, was proved in 2015 by
Muratovi\'c-Ribi\'c, Pott, Thomson and Wang, in its equivalent form that every Costas polynomial over a
prime field is a monomial; since circularity forces $n+1$ to be prime, this covers every circular sequence~\cite{p1-mpttw2015,p1-torresrubio2023}. The remaining open statement is the Golomb--Taylor conjecture of
1984, that a Costas sequence is singly periodic if and only if it is Welch;
Moreno proved in 1992~\cite{p1-moreno92} that the stronger \emph{shifting property} characterises
the Welch construction, leaving the gap between single periodicity and the
shifting property open.

\item \textbf{Constructed and augmented arrays.} This is where the present
investigation makes progress; see Papers~2 and~3.

\item \textbf{Sporadic arrays.} Everything else. This branch is defined by
absence of structure and is therefore immune to algebraic attack by
construction. Paper~3 shows that branch (ii) is thinner than
expected, which makes branch (iii) larger, not smaller.
\end{enumerate}

\begin{remark}[Circular arrays are a strict subclass]
\label{p1-rem:circular}
A circular Costas sequence becomes an ordinary Costas array once its entries are read as integers:
two equal integer displacements would also be equal modulo $n$ horizontally and modulo $n+1$
vertically, which circularity forbids. The converse fails. The order-3 Costas array $(1,3,2)$ is
not circular, since its lag-1 circular differences $3-1$, $2-3$ and $1-2$ are $2$, $3$ and $3$
modulo 4. An exhaustive check at orders 3 to 6 finds circular sequences only at orders 4 and 6, 8
and 12 of them, which is exactly the number $\varphi(p-1)(p-1)$ of exponential Welch arrays of
those orders, as the two theorems require. Excluding circular arrays at order 32 therefore closes a
subclass of the problem, not the problem.
\end{remark}

\begin{remark}
The decomposition is a partition only because branch (iii) is the complement.
The category ``sporadic'' is epistemic. During this investigation a
catalogue-generation bug produced an array that appeared unreachable by
construction until an additional closure operation was added; nothing about the
array changed. The same caution applies to published sporadic counts, although
those rest on the enumeration authors' own classification and not on any
generator built here.
\end{remark}

%=====================================================================
\section{Computational baseline}
\label{p1-sec:baseline}

\subsection{Exact solver engineering}

The starting point was a depth-first search on vertex potentials that allocated
per-node vectors and sorted candidate values at every node, achieving
$5.84\times 10^5$ nodes/s at order 32.

For $n\le 32$ the set of used differences at a fixed horizontal distance lies in
$[-(n-1),n-1]$, that is $2n-1 \le 63$ values, and therefore fits in a single
64-bit word. Rewriting with one \texttt{uint64\_t} mask per distance, no
per-node allocation and no sort gives the same search tree and the following
measured speedup:

\begin{center}
\begin{tabular}{@{}lrr@{}}
\toprule
task & original & bitmask rewrite \\
\midrule
$n=13$ exhaustive (18{,}290{,}044 nodes, identical) & $8.55$ s & $1.26$ s \\
$n=32$, 20 s pilot & $5.84\times10^5$ n/s & $4.34\times10^6$ n/s \\
\bottomrule
\end{tabular}
\end{center}

\subsection{Growth law}

Exhaustive enumeration with the vertical-complement symmetry $p_0\le
\lfloor(n-1)/2\rfloor$ gives the following measured tree sizes. Solution counts
were validated against the published census in every case.

\begin{center}
\begin{tabular}{@{}rrrr@{}}
\toprule
$n$ & nodes & ratio & time (1 core) \\
\midrule
11 & 789{,}512 & --- & 0.04 s\\
12 & 3{,}488{,}681 & 4.42 & 0.21 s\\
13 & 18{,}290{,}044 & 5.24 & 1.26 s\\
14 & 83{,}176{,}085 & 4.55 & 19.6 s\\
15 & 435{,}090{,}281 & 5.23 & 82.9 s\\
16 & 2{,}001{,}473{,}413 & 4.60 & 237.9 s\\
\bottomrule
\end{tabular}
\end{center}

The geometric mean ratio is $\mathbf{4.80}$ per order, consistent with the
published observation of slightly worse than $5\times$ per order. Extrapolating
from the measured order-29 cost of $366.55$ CPU-years gives roughly
$5\times10^4$ CPU-years on 2010-era hardware for order 32, or of order $10^3$--$10^4$ modern core-years. This projection concerns unrestricted enumeration of all
order-32 permutations. It should not be confused with the involution-only search of Paper~8, which
the symmetry restriction makes many orders of magnitude cheaper, projected at 30--45 single-core hours. The projection also assumes the efficiency of the 2011 enumeration: Paper~9
measures a straightforward column-by-column search at about $1.5\times10^{20}$ nodes at order 32, over
$10^6$ core-years, and finds the 2011 method at least 36 times more efficient.

\subsection{Order 17 as a calibration point}

Order 17 exhaustive enumeration was measured by splitting on the nine
symmetry-reduced top-level branches $p_0 = 0,\dots,8$:

\begin{center}
\begin{tabular}{@{}lrrr@{}}
\toprule
branch & nodes & time & rate \\
\midrule
$p_0=0$ & $1.099\times10^9$ & 95.8 s & $1.15\times10^7$/s \\
$p_0=1$ & $1.138\times10^9$ & 94.8 s & $1.20\times10^7$/s \\
$p_0=2$ & $1.144\times10^9$ & 97.0 s & $1.18\times10^7$/s \\
\bottomrule
\end{tabular}
\end{center}

Nine branches at that rate is approximately $10^{10}$ nodes and
\textbf{about 14 minutes on one core}. The project's consolidated report uses this figure to calibrate distributed SAT approaches.

\subsection{Fail-closed discipline}\label{p1-sec:discipline}

Two defects in the original solver are recorded because they are easy to
reproduce in any successor. First, it had no completeness flag, so a
time-limited subtree was indistinguishable from an exhausted one --- fatal for
any non-existence claim. Second, the vertical-complement symmetry was silently
skipped whenever a prefix was supplied, since the guard tested $k=0$ and $k$
never reaches 0 under a prefix; a prefix manifest must therefore be
symmetry-reduced itself. Both are corrected in the released solver, which emits
one of \texttt{EXHAUSTED}, \texttt{WITNESS\_FOUND}, \texttt{TIME\_LIMIT},
\texttt{NODE\_LIMIT}, \texttt{INVALID\_PREFIX} or \texttt{INTERNAL\_ERROR}, and
whose merge tool refuses an exhaustive verdict unless every unit reports
\texttt{EXHAUSTED}, all units agree on parameters and catalogue hash, and the
coverage is exactly $0..N{-}1$ with no gap or duplicate.

%=====================================================================
\section{Order reachability settles the taxonomy without the names}
\label{p1-sec:reachability}

The literature's $G_3$/$G_4$/Taylor/Lempel designations are not used consistently. Resolving them into exact source-defined equations before specialising each to $\F_{32}$ would be worth having, but the question that resolution exists to answer
can be settled first, and independently of every naming decision, by asking which
\emph{orders} corner surgery reaches rather than what each family is called.

Every $W_1(p)$ for prime $p\le43$ and every $G_2(q)$ for supported prime power
$q\le43$ was generated over all primitive parameter pairs, and each seed set was
closed under corner deletion and corner addition:

\begin{center}
\begin{tabular}{@{}lrrl@{}}
\toprule
family & native order & seeds & orders reachable by corner surgery \\
\midrule
$G_2(29)$ & 27 & 144 & 26, 27, 28, 29 \\
$G_2(31)$ & 29 & 64 & 28, 29, 30 \\
$G_2(32)$ & 30 & 180 & 28, 29, 30 \\
$G_2(37)$ & 35 & 144 & \textbf{34, 35} \\
$W_1(29)$ & 28 & 12 & 26, 27, 28, 29 \\
$W_1(31)$ & 30 & 8 & 29, 30 \\
$W_1(37)$ & 36 & 12 & \textbf{36} \\
\bottomrule
\end{tabular}
\end{center}

\begin{observation}[The corner-surgery order gap]
\label{p1-obs:ordergap}
Corner surgery from these seeds reaches exactly
$\{26,\dots,30\}\cup\{34,35,36\}$. \textbf{Orders 31, 32 and 33 are reached by
nothing.} In particular:
\begin{itemize}
\item $G_2(37)$ reaches 34 by one corner deletion and stops. No second deletion
  survives, so whatever equations the $G_4$ and $G_5$ labels denote, the families
  they name are \emph{empty at $q=37$}. Since 37 is the only prime power that
  could reach 33 or 32 by deletion, the conditional constructions cannot produce
  either order.
\item $W_1(37)$ admits no corner deletion at all, so order 35 is not reached that
  way either.
\end{itemize}
\end{observation}

The naming ambiguity therefore does not affect the order-32 reachability conclusion, although it remains material to the completeness and historical accuracy of the construction catalogue. Whatever exact
equations $G_3$, $G_4$, $G_5$, Taylor and Lempel denote, each is a corner
deletion of a $G_2$ or $W_1$ seed, and no such deletion has order 31, 32 or 33.
The taxonomy remains worth completing for correctness of the record and for
neighbouring-order work, but it cannot supply an order-32 witness.

\begin{remark}[Order 31 is not a corner-surgery descendant]
The gap is informative about order 31 as well. Order-31 Costas arrays exist, and
rung $k=1$ exhibits eight of them in one dihedral orbit --- but they do not appear
in the table above, because rung $k=1$ performs \emph{general} one-point
augmentation into a compressed sub-board, not corner addition. So the eight
order-31 arrays are precisely non-corner augmentations, which is consistent with the heredity measurement of Paper~3: corner operations preserve structure and are rare,
and order 31 needs the operation that does not. Applying the same general
augmentation twice is rung $k=2$, exhausted over all 664 generated order-30 cores
with zero order-32 arrays (Paper~2).
\end{remark}

%=====================================================================

\section{Status of the claims}\label{p1-sec:status}
The exclusion of circular arrays at order 32 is due to Golomb and Moreno~\cite{p1-golombmoreno}, and
the theorem that every circular Costas array is Welch to Muratovi\'c-Ribi\'c, Pott, Thomson and
Wang~\cite{p1-mpttw2015}, whose proof completed the Golomb--Moreno conjecture~\cite{p1-torresrubio2023};
Moreno's shifting-property characterisation is also published~\cite{p1-moreno92}; single periodicity remains open
pending the Golomb--Taylor conjecture. The growth-law figures are measurements made with this
project's exact solvers, and the projections drawn from them are extrapolations. The reachability
analysis is a deductive argument over the construction families as stated.

\paragraph{Relation to prior work.} Known material, newly organised; the reachability analysis and the growth measurements are this project's.
\begin{paperrefs}
\bibitem{p1-golombtaylor} S.~W.~Golomb and H.~Taylor, \emph{Constructions and
properties of Costas arrays}, Proc.\ IEEE \textbf{72} (1984).
\bibitem{p1-drakakis27} K.~Drakakis, S.~Rickard, J.~K.~Beard, R.~Caballero,
F.~Iorio, G.~O'Brien and J.~Walsh, \emph{Results of the enumeration of Costas
arrays of order 27}, IEEE Trans.\ Inform.\ Theory \textbf{54} (2008).
\bibitem{p1-drakakis28} K.~Drakakis, F.~Iorio and S.~Rickard, \emph{The enumeration of Costas arrays of order 28 and its consequences}, Advances in Mathematics of Communications, 2011.
\bibitem{p1-drakakis29} K.~Drakakis, F.~Iorio, S.~Rickard and J.~Walsh,
\emph{Results of the enumeration of Costas arrays of order 29}, IEEE Trans.\
Inform.\ Theory \textbf{57} (2011).
\bibitem{p1-oeis} OEIS A008404 (number of Costas arrays of order $n$) and A001441
(inequivalent Costas arrays under the dihedral group).
\bibitem{p1-gulec2026ucm} F.~Gulec and V.~Abolghasemi, \emph{Universal Costas matrices: towards a general framework for Costas array construction}, IEEE International Conference on Communications (ICC), Glasgow, 2026; arXiv:2602.03407.
\bibitem{p1-warnke2023density} L.~Warnke, B.~Correll~Jr. and C.~N.~Swanson, \emph{The density of Costas arrays decays exponentially}, IEEE Transactions on Information Theory 69(1) (2023), 575--581.
\bibitem{p1-correll2023difference} B.~Correll~Jr. and C.~N.~Swanson, \emph{Difference-based structural properties of Costas arrays}, Designs, Codes and Cryptography 91(3) (2023), 779--794.
\bibitem{p1-vulakh2022mdraco} D.~Vulakh and R.~Finkel, \emph{Parallel $m$-dimensional relative ant colony optimization (mDRACO) for the Costas-array problem}, Soft Computing 26(12) (2022), 5765--5772.
\bibitem{p1-rosic2026gist} V.~Rosi\'c, \emph{Order-32 Costas array: reproducible bounded partial result}, GitHub gist, created 19 August 2026, \url{https://gist.github.com/vukrosic/afa4096cd8b16fe17dbf413f63422b72}. Unpublished; the DRAT proof package is held outside the gist.
\bibitem{p1-dumitru2026localpatterns} B.~Dumitru and C.~Budala, \emph{Local patterns and the cost of Costas array search}, preprint hal-05731061v1, 2026, \url{https://hal.science/hal-05731061}. Not peer-reviewed.
\bibitem{p1-dumitru2026symmetriccensus} B.~Dumitru and C.~Budala, \emph{An exhaustive census of main-diagonal symmetric Costas arrays of orders 37--42}, preprint arXiv:2608.28690, 2026. Not peer-reviewed.
\bibitem{p1-golombmoreno} S.~W.~Golomb and O.~Moreno, \emph{On periodicity properties of Costas arrays and a conjecture on permutation polynomials}, IEEE Transactions on Information Theory 42(6) (1996), 2252--2253.
\bibitem{p1-mpttw2015} A.~Muratovi\'c-Ribi\'c, A.~Pott, D.~Thomson and Q.~Wang, \emph{On the characterization of a semi-multiplicative analogue of planar functions over finite fields}, in Topics in Finite Fields, Contemporary Mathematics 632, American Mathematical Society, 2015, 317--325.
\bibitem{p1-torresrubio2023} J.~Torres and I.~Rubio, \emph{Circular Costas maps: a multidimensional analog of circular Costas sequences}, Cryptography and Communications 15(5) (2023), 941--958, \href{https://doi.org/10.1007/s12095-023-00654-2}{\texttt{doi:10.1007/s12095-023-00654-2}}.
\bibitem{p1-moreno92} O.~Moreno, \emph{Survey of results on signal patterns for
locating one or multiple targets}, and the shifting-property characterisation
(1992).
\end{paperrefs}

\clearpage
\chapter{Two-Point Augmentations of Algebraic Order-30 Cores}\label{p2-paper}
\paperauthor
\paperabstract{An order-32 Costas array might arise by adding two dots to an order-30 Costas array. Order 30
has never been enumerated, but its algebraically derived arrays can be generated: closing the Welch
and Golomb constructions under corner deletion and corner addition yields 664 order-30 arrays in 85
dihedral orbits, and, as a check, reproduces the complete published order-28 census exactly. An
exhaustive search of all $326{,}709{,}248$ two-point augmentations of these cores finds no Costas
array. The minimum defect is 2, attained by a single dihedral orbit, that of the known near-solution
$N_{32}$. One-dot augmentation of the same cores yields exactly one orbit of order-31 Costas arrays,
and $N_{32}$ is its corner extension, which gives the near-solution a complete genealogy.}

\section{Introduction}\label{p2-sec:paperintro}
Costas arrays and their algebraic constructions are surveyed in~\cite{p2-golombtaylor}. A natural way for an order-32 Costas array to arise is by augmentation: inserting extra rows,
columns and dots into a smaller Costas array. This paper studies the rungs of that ladder that start
from order 30, the largest order with algebraic arrays near 32. The question is whether any order-32
Costas array is a two-point augmentation of an algebraic order-30 core, and, since none is, what the
closest misses look like.

\paragraph{Dependencies.} This paper uses the shared notation and the catalogue generator described here. Its exclusion does not depend on the other papers.

\section{Preliminaries}\label{p2-sec:notation}
A \emph{Costas array} of order $n$ is an $n\times n$ permutation array in which the displacement
vectors between pairs of dots are all distinct. Writing the array as a permutation $h$ of
$\{0,\dots,n-1\}$, with the dot of column $i$ in row $h_i$, the condition is that for each
horizontal separation $m$ the vertical differences $h_{i+m}-h_i$ are distinct. Costas arrays are
known at every order up to 31; orders 32 and 33 are the smallest for which existence is open. The
dihedral group $D_4$ of the square maps Costas arrays to Costas arrays, and arrays are counted
either individually or by $D_4$ orbit (equivalence class). Further notation and terms are collected
in the notation table and glossary at the front of this volume.

\section{Formulation}

\begin{definition}[Augmentation rung]
Fix $k$ and a Costas core $K$ of order $n-k$. Choose a set $C$ of $k$ columns
and a set $R$ of $k$ rows of the $n\times n$ board. Place the core dots at the
surviving columns and rows in increasing order, and place $k$ further dots as a
bijection $C\to R$. The rung is the exhaustive search over all
$(K,C,R,\mu)$.
\end{definition}

Two remarks on the formulation. First, when $C$ and $R$ both consist of
end-blocks the compression is a translation and the core's displacements are
inherited unchanged; this \emph{contiguous} sub-case is the actual construction
mechanism and is reported separately. Second, core-first pruning --- rejecting a
placement as soon as the placed core carries a repeated displacement --- is
exact for defect 0, because the core dots are a subset of the $n$ dots. It is
\emph{not} exact for defects above 0, so the reported histograms above score 0
count only core-intact configurations.

\section{Catalogue generation and census verification}
\label{p2-sec:catalogue}

The algebraic possibilities at each order near 32 are determined by the prime
powers available. For order 30 they are $W_1(31)$ and $G_2(32)$; no other family
applies, since $W_2$ needs $p=32$, $W_3$ needs $p=33$, $G_3$ needs $q=33$, $G_4$
needs $q=34$ and $G_5$ needs $q=35$, none of which is prime or a prime power.

A first generator seeded $W_1(29)$, $W_1(31)$, $G_2(31)$ and $G_2(32)$ and
closed \emph{downward} under corner deletion. It produced 704 order-28 arrays in
88 dihedral orbits against a published census of 712 in 89, and 660 order-30
arrays.

\begin{observation}[The missing orbit]
\label{p2-obs:missing}
The deficit was exactly one dihedral orbit and was caused by the closure being
downward-only. $G_2(29)$ has order $q-2 = 27$ and can reach order 28 only
\emph{upwards}. Adding $G_2(29)$ as a seed and corner \emph{addition} to the
closure recovers the missing orbit, whose canonical representative is
\[
(0,16,21,12,18,22,23,26,10,7,24,14,27,6,17,5,25,11,20,2,4,19,15,13,8,1,9,3).
\]
Eight configurations land in this orbit, including $G_2(29,\alpha=3,\beta=3)$
with a corner dot at $(0,27)$ and $G_2(29,\alpha=3,\beta=10)$ with one at
$(0,0)$.
\end{observation}

With the corrected closure the generator produces:

\begin{center}
\begin{tabular}{@{}lrrl@{}}
\toprule
order & generated & published census & status \\
\midrule
27 & 192 / 27 orbits & 204 / 29 & 2 orbits short \\
\textbf{28} & \textbf{712 / 89 orbits} & \textbf{712 / 89} & \textbf{complete} \\
29 & 148 / 21 orbits & 164 / 23 & 16 arrays short \\
30 & 664 / 85 orbits & never enumerated & --- \\
\bottomrule
\end{tabular}
\end{center}

\begin{proposition}[Completeness of the order-28 catalogue]
The generated order-28 catalogue is the complete census.
\end{proposition}
\begin{proof}
The published enumeration is exhaustive with 712 arrays. The generated
catalogue contains 712 distinct permutations of $\{0,\dots,27\}$, each verified
to have Costas defect 0, and is closed under $D_4$ with 89 orbits. Two sets of
equal finite cardinality, one contained in the other, are equal.
\end{proof}

\begin{remark}
Order 29 landing on exactly $164-16=148$ is an internal consistency check: the
literature records precisely 16 order-29 arrays not accounted for by the
Golomb and Welch constructions, and the corrected generator finds every array
that is. Order 27 falling two orbits short is consistent with the $T_4(31,1)$
family and the single known sporadic array, neither of which the closure
reaches.
\end{remark}

\section{The rungs $k=1$ and $k=2$}
\label{p2-sec:rungs}

All rungs were run with core-first pruning, independent re-scoring of every
reported low-defect configuration, and a fail-closed coverage manifest.

\begin{center}
\begin{tabular}{@{}rlrrrr@{}}
\toprule
$k$ & core order & cores & measured placements & Costas found & minimum defect \\
\midrule
1 & 30 (algebraic) & 664 & $6.3\times10^5$ & \textbf{8} & 0 \\
2 & 30 (algebraic) & 664 & $2.091\times10^{7}$ & \textbf{0} & \textbf{2} \\
\bottomrule
\end{tabular}
\end{center}

\paragraph{Rung $k=1$ (order 30 $\to$ 31).}
Exactly 8 order-31 Costas arrays are found, forming exactly one dihedral orbit.
Hence there is, up to the square symmetries, \emph{exactly one} order-31 Costas
array reachable by one-dot augmentation of the algebraic order-30 cores. Defect
histogram: $0{:}8$, $1{:}0$, $2{:}24$, $3{:}16$, $4{:}124$, $5{:}112$,
$6{:}596$.

\paragraph{Rung $k=2$ (order 30 $\to$ 32).}
The engine processed $20{,}911{,}360=85\binom{32}{2}^2$ core placements in 33 seconds on one core: one representative of each of the catalogue's 85
dihedral orbits, in every choice of two inserted columns and two inserted rows. Each placement admits $2!$ bijections of the
added dots, and by $D_4$-equivariance (the remark on canonical mode below) the search covers all $664\binom{32}{2}^2\,2!=326{,}709{,}248$ a priori configurations. It
finds \textbf{zero} Costas arrays. (Earlier statements of this count gave $324{,}741{,}120$, the value for a superseded 660-array catalogue.)
Defect histogram: $0{:}0$, $1{:}0$, $2{:}8$, $3{:}0$, $4{:}24$, $5{:}16$,
$6{:}40$, $7{:}56$, $8{:}104$, $9{:}168$, $10{:}584$.

The eight defect-2 configurations were dumped and re-verified with an
independent scorer: all are valid permutations, all have exact defect 2, and
they are \emph{precisely} the dihedral orbit of the known candidate $N_{32}$.
Nothing else in the space attains 2, and defects 1 and 3 are empty. The
contiguous sub-case comprises 11{,}880 configurations with the same optimum.

\begin{theorem}[Two-point augmentation]
\label{p2-thm:rungs}
No order-32 Costas array arises as a two-point augmentation of any of the 664 generated order-30
Costas cores. Among all $326{,}709{,}248$ two-point augmentations the minimum defect is 2, attained
uniquely up to the dihedral group.

The conclusion is catalogue-relative: order 30 has never been enumerated, so its 664 cores are the
generated set.
\end{theorem}

\begin{remark}[Canonical mode is exact]
Restricting cores to one representative per dihedral orbit is exact for
existence and for the minimum defect, because the configuration space is
$D_4$-equivariant and the defect is $D_4$-invariant, so every configuration maps
to one with a canonical core and the same defect. It is \emph{not} exact for
counting solutions.
\end{remark}

\section{Genealogy of the score-2 object}
\label{p2-sec:genealogy}

The rung searches identify the parent structure of $N_{32}$ directly. The
order-31 array found at $k=1$ is
{\footnotesize\[
C_{31} = (30,2,8,1,14,12,7,4,20,0,18,9,13,21,10,11,26,25,6,28,24,3,15,5,19,22,27,29,16,23,17),
\]}
and the order-32 optimum found at $k=2$ is
{\footnotesize\[
N_{32}' = (31,30,2,8,1,14,12,7,4,20,0,18,9,13,21,10,11,26,25,6,28,24,3,15,5,19,22,27,29,16,23,17),
\]}
which is $C_{31}$ with a single dot adjoined at the corner. $N_{32}'$ is a
dihedral image of the known candidate $N_{32}$.

\paragraph{Attribution.} $N_{32}$ was publicly documented by Rosi\'c on 19 August
2026~\cite{p2-rosic2026gist}, constructed as a corner extension of an exceptional order-31 array (``W0'') in
Beard's database of known Costas arrays. That order-31 array is a dihedral image of $C_{31}$, so the two
genealogies coincide: our searches reached the same order-31 array by one-dot augmentation of the
algebraic order-30 cores, and the same order-32 permutation by one further corner dot. We reached it
independently and claim no priority for it. Rosi\'c also reports a proof, checked in DRAT format, that
no order-32 Costas array lies within Hamming distance 12 of the candidate, and an exhaustive program
enumeration of distances 13 to 16 that found none; the 3.79\,GB proof package is held outside the
published record and was not rechecked here. For comparison, the best order-32 permutation reported by an
ant-colony search had 17 violations of the Costas property~\cite{p2-vulakh2022mdraco}.

\begin{observation}[Genealogy]
\label{p2-obs:genealogy}
\[
\begin{aligned}
&\text{Welch-30 (algebraic)}\;\xrightarrow{\ \text{free adjunction}\ }\;C_{31}\text{ (unique derived order-31 orbit)}\\
&\qquad\xrightarrow{\ \text{minimum-cost adjunction}\ }\;N_{32}\ (\text{defect }2).
\end{aligned}
\]
The first adjunction is free; the second costs 2 and nothing costs less.
\end{observation}

This is a better description of $N_{32}$ than ``Welch-30 plus two dots'': it
identifies an exact intermediate Costas object and localises the failure
entirely to the final adjunction.

\section{The toll is a minimum, not a pointwise quantity}
\label{p2-sec:toll}

The bottoms of the two defect spectra coincide after a shift of 2
($8,0,24,16$ appears at defects $0,1,2,3$ for $k=1$ and at $2,3,4,5$ for
$k=2$). This suggested a pointwise parent--child relation. It was tested and is
false.

For each two-dot child $Q$ with defect at most 10, let
$\tau(Q)=\score(Q)-\min\{\score(Q-x),\score(Q-y)\}$ over its two one-dot
parents. Measured distribution:

\begin{center}
\begin{tabular}{@{}lrrrrrrrrrr@{}}
\toprule
$\tau$ & 0 & 1 & 2 & 3 & 4 & 5 & 6 & 7 & 8 & negative \\
\midrule
count & 76 & 48 & 152 & 56 & 296 & 108 & 192 & 24 & 16 & 32 \\
\bottomrule
\end{tabular}
\end{center}

Thirty-two children score \emph{below} both parents, so the defect is not even
monotone along the ladder. The histogram coincidence is a coincidence of
multiplicities.

\begin{observation}
The correct statement is: \emph{the second-rung class has a uniform minimum toll
of 2, attained uniquely up to $D_4$}. It is not the case that each second
adjunction costs exactly 2.
\end{observation}

%=====================================================================

\section{Status of the claims}\label{p2-sec:status}
The order-28 census used as a check is that of~\cite{p2-drakakis28}. The results are computational and uncertified. They are relative to the generated order-30
catalogue, since order 30 has never been enumerated: an order-30 Costas array outside the catalogue
would not be covered. The rung computations were re-run from a single binary, which reproduced the
recorded outcomes exactly, and every reported low-defect configuration was re-scored by an
independent scorer. We have not found these results in the literature, but a systematic search has
not been completed.

\paragraph{Relation to prior work.} An apparently new computational result. Its novelty has not been established by a systematic literature search. The defect-2 permutation $N_{32}$ was publicly documented by another author in August 2026; we reached it independently and claim no priority for it.
\begin{paperrefs}
\bibitem{p2-golombtaylor} S.~W.~Golomb and H.~Taylor, \emph{Constructions and
properties of Costas arrays}, Proc.\ IEEE \textbf{72} (1984).
\bibitem{p2-rosic2026gist} V.~Rosi\'c, \emph{Order-32 Costas array: reproducible bounded partial result}, GitHub gist, created 19 August 2026, \url{https://gist.github.com/vukrosic/afa4096cd8b16fe17dbf413f63422b72}. Unpublished; the DRAT proof package is held outside the gist.
\bibitem{p2-vulakh2022mdraco} D.~Vulakh and R.~Finkel, \emph{Parallel $m$-dimensional relative ant colony optimization (mDRACO) for the Costas-array problem}, Soft Computing 26(12) (2022), 5765--5772.
\bibitem{p2-drakakis28} K.~Drakakis, F.~Iorio and S.~Rickard, \emph{The enumeration of Costas arrays of order 28 and its consequences}, Advances in Mathematics of Communications, 2011.
\end{paperrefs}

\clearpage
\chapter{Deletion-Minor Exclusions at Order 32 from the Complete Order-29 and Order-28 Censuses}\label{p3-paper}
\paperauthor
\paperabstract{We show computationally that no order-32 Costas array contains an order-29 Costas array as a
three-row, three-column deletion minor, and that none contains an order-28 Costas array as a
four-row, four-column deletion minor. Both statements quantify over complete published censuses: all
164 order-29 arrays in 23 classes, and all 712 order-28 arrays in 89 classes. The order-28 search
covers $115{,}087{,}822{,}400$ core placements in 89 work units, all exhausted and independently
merged; the order-29 search covers $565{,}836{,}800$. Neither result is known to imply the other: deriving either would need a hereditary-minor property, and over the complete
order-28 census a single corner deletion preserves the Costas property only $14.6\%$ of the time.}

\section{Introduction}\label{p3-sec:paperintro}
If an order-32 Costas array exists, deleting some rows and columns from it leaves a smaller
permutation array, which may or may not be Costas. This paper asks whether any order-32 Costas array
can contain an order-29 or an order-28 Costas array in this way. Because orders 28 and 29 have been enumerated completely~\cite{p3-drakakis28,p3-drakakis29}, the answers quantify over every Costas array of those orders, not just over
algebraic ones.

\paragraph{Dependencies.} This paper uses the order-29 census reconstructed in Paper~4 and the order-28 catalogue verified here. Its theorems do not otherwise depend on the other papers.

\section{Preliminaries}\label{p3-sec:notation}
A \emph{Costas array} of order $n$ is an $n\times n$ permutation array in which the displacement
vectors between pairs of dots are all distinct. Writing the array as a permutation $h$ of
$\{0,\dots,n-1\}$, with the dot of column $i$ in row $h_i$, the condition is that for each
horizontal separation $m$ the vertical differences $h_{i+m}-h_i$ are distinct. Costas arrays are
known at every order up to 31; orders 32 and 33 are the smallest for which existence is open. The
dihedral group $D_4$ of the square maps Costas arrays to Costas arrays, and arrays are counted
either individually or by $D_4$ orbit (equivalence class). Further notation and terms are collected
in the notation table and glossary at the front of this volume.

\section{Formulation}

\begin{definition}[Augmentation rung]
Fix $k$ and a Costas core $K$ of order $n-k$. Choose a set $C$ of $k$ columns
and a set $R$ of $k$ rows of the $n\times n$ board. Place the core dots at the
surviving columns and rows in increasing order, and place $k$ further dots as a
bijection $C\to R$. The rung is the exhaustive search over all
$(K,C,R,\mu)$.
\end{definition}

\begin{proposition}[Augmentation and deletion minors]
\label{p3-prop:augminor}
An order-$n$ permutation $P$ contains a Costas array $K$ of order $n-k$ as a $k$-row, $k$-column
deletion minor if and only if $P$ lies in the rung-$k$ augmentation space generated from $K$.
\end{proposition}
\begin{proof}
Suppose deleting columns $C$ and rows $R$, with $|C|=|R|=k$, removes exactly the $k$ dots of $P$
in $C\times R$ and leaves $K$ once the gaps are closed. The surviving dots then sit in the surviving
columns and rows at the order-preserving positions of $K$'s dots, and $P$'s dots in the deleted
columns define a bijection $\mu:C\to R$. So $(K,C,R,\mu)$ is a rung-$k$ configuration whose array
is $P$. Conversely, a configuration $(K,C,R,\mu)$ yields a permutation from which deleting $C$, $R$
and the dots of $\mu$ recovers $K$.
\end{proof}

\noindent Hence a rung-$k$ search over a complete census of order $n-k$ finds every order-$n$ Costas
array that has a Costas deletion minor of that order.

Two remarks on the formulation. First, when $C$ and $R$ both consist of
end-blocks the compression is a translation and the core's displacements are
inherited unchanged; this \emph{contiguous} sub-case is the actual construction
mechanism and is reported separately. Second, core-first pruning --- rejecting a
placement as soon as the placed core carries a repeated displacement --- is
exact for defect 0, because the core dots are a subset of the $n$ dots. It is
\emph{not} exact for defects above 0, so the reported histograms above score 0
count only core-intact configurations.

\section{Catalogue generation and census verification}
\label{p3-sec:catalogue}

The algebraic possibilities at each order near 32 are determined by the prime
powers available. For order 30 they are $W_1(31)$ and $G_2(32)$; no other family
applies, since $W_2$ needs $p=32$, $W_3$ needs $p=33$, $G_3$ needs $q=33$, $G_4$
needs $q=34$ and $G_5$ needs $q=35$, none of which is prime or a prime power.

A first generator seeded $W_1(29)$, $W_1(31)$, $G_2(31)$ and $G_2(32)$ and
closed \emph{downward} under corner deletion. It produced 704 order-28 arrays in
88 dihedral orbits against a published census of 712 in 89, and 660 order-30
arrays.

\begin{observation}[The missing orbit]
\label{p3-obs:missing}
The deficit was exactly one dihedral orbit and was caused by the closure being
downward-only. $G_2(29)$ has order $q-2 = 27$ and can reach order 28 only
\emph{upwards}. Adding $G_2(29)$ as a seed and corner \emph{addition} to the
closure recovers the missing orbit, whose canonical representative is
\[
(0,16,21,12,18,22,23,26,10,7,24,14,27,6,17,5,25,11,20,2,4,19,15,13,8,1,9,3).
\]
Eight configurations land in this orbit, including $G_2(29,\alpha=3,\beta=3)$
with a corner dot at $(0,27)$ and $G_2(29,\alpha=3,\beta=10)$ with one at
$(0,0)$.
\end{observation}

With the corrected closure the generator produces:

\begin{center}
\begin{tabular}{@{}lrrl@{}}
\toprule
order & generated & published census & status \\
\midrule
27 & 192 / 27 orbits & 204 / 29 & 2 orbits short \\
\textbf{28} & \textbf{712 / 89 orbits} & \textbf{712 / 89} & \textbf{complete} \\
29 & 148 / 21 orbits & 164 / 23 & 16 arrays short \\
30 & 664 / 85 orbits & never enumerated & --- \\
\bottomrule
\end{tabular}
\end{center}

\begin{proposition}[Completeness of the order-28 catalogue]
The generated order-28 catalogue is the complete census.
\end{proposition}
\begin{proof}
The published enumeration is exhaustive with 712 arrays. The generated
catalogue contains 712 distinct permutations of $\{0,\dots,27\}$, each verified
to have Costas defect 0, and is closed under $D_4$ with 89 orbits. Two sets of
equal finite cardinality, one contained in the other, are equal.
\end{proof}

\begin{remark}
Order 29 landing on exactly $164-16=148$ is an internal consistency check: the
literature records precisely 16 order-29 arrays not accounted for by the
Golomb and Welch constructions, and the corrected generator finds every array
that is. Order 27 falling two orbits short is consistent with the $T_4(31,1)$
family and the single known sporadic array, neither of which the closure
reaches.
\end{remark}

\section{The rungs $k=3$ and $k=4$}
\label{p3-sec:rungs}

All rungs were run with core-first pruning, independent re-scoring of every
reported low-defect configuration, and a fail-closed coverage manifest.

\begin{center}
\begin{tabular}{@{}rlrrrr@{}}
\toprule
$k$ & core order & cores & measured placements & Costas found & minimum defect \\
\midrule
3 & 29 (\emph{census-complete}) & 164 arrays & $5.658\times10^{8}$ & \textbf{0} & 4 \\
4 & 28 (\emph{census-complete}) & 712 arrays & $1.151\times10^{11}$ & \textbf{0} & --- \\
\bottomrule
\end{tabular}
\end{center}

\paragraph{Rung $k=4$ (order 28 $\to$ 32): completed.}
Run in canonical mode as 89 deterministic work units, one dihedral orbit each,
against the census-complete catalogue (\texttt{CENSUS ... MATCH}, catalogue hash
\texttt{b64cc86a0058f1b3}). Independently merged with the fail-closed tool:

\begin{center}
\begin{tabular}{@{}lr@{}}
\toprule
merge verdict & \texttt{EXHAUSTED\_AND\_VERIFIED} \\
units & 89 of 89, all \texttt{EXHAUSTED} \\
core placements & $115{,}087{,}822{,}400$ \\
survivors of the core-first test & $2{,}401$ \\
configurations evaluated & $57{,}624$ \\
\textbf{Costas arrays found} & \textbf{0} \\
aggregate CPU & $2{,}389.8$ s; longest unit $41.1$ s \\
\bottomrule
\end{tabular}
\end{center}

\noindent The engineering rate was $3.79\times10^7$ core placements per second,
achieved by bucketing the core-pair test by column displacement and replacing the
per-placement clear with a generation stamp. Because order 28 has no sporadic
arrays, ``census-complete'' and ``all algebraic'' coincide here, so this is the
first rung whose negative statement quantifies over \emph{every} core of its
order rather than an algebraic subset.

\begin{theorem}[Census-complete rung $k=4$]
\label{p3-thm:rung4census}
No order-32 Costas array contains an order-28 Costas array as a four-row, four-column deletion
minor. The statement quantifies over all 712 order-28 Costas arrays.
\end{theorem}

\begin{remark}[Measured placements]
\label{p3-rem:colconv}
The placements column records core placements measured in the runs, not a priori counts of the
configuration space. Re-running the rungs from a single binary reproduced the recorded outcomes
exactly: zero Costas arrays, with minimum defect 4 at $k=3$.
\end{remark}


\paragraph{Rung $k=3$ (order 29 $\to$ 32).}
The order-29 catalogue is the complete published census, reconstructed and verified as
described in the paper on the order-29 census in these proceedings. In brief: The complete order-29 census was reconstructed from Table 2 of the enumeration
paper, which lists the 23 lexicographically minimal equivalence-class
representatives with their construction method. Every representative was
independently validated as an order-29 Costas permutation and expanded under
$D_4$, reproducing all published figures exactly: \textbf{164 arrays, 23
equivalence classes, 5 symmetric}, with per-method counts
$G_0\,4$, $G_2\,64$, $G_3\,24$, $W_0\,24$, $W_2\,32$, $RW_0\,16$.

Running the rung over the full census: \texttt{CENSUS ... MATCH}, \texttt{STATUS
EXHAUSTED}, $565{,}836{,}800$ placements, 420 survivors, 2{,}520 configurations,
\textbf{zero Costas arrays}, minimum defect 4 attained by 3 configurations,
12.7\,s. Fail-closed merge re-scored all 3 with zero failures:
\texttt{EXHAUSTIVE, NO COSTAS ARRAY FOUND}.

\begin{theorem}[Census-complete rung $k=3$]
\label{p3-thm:rung3census}
No order-32 Costas array contains an order-29 Costas array as a three-row,
three-column deletion minor.
\end{theorem}


Neither this theorem nor Theorem~\ref{p3-thm:rung4census} is known to imply the other
(Remark~\ref{p3-rem:independent}): deriving either would need a hereditary-minor property that is not established and is not supported by the
order-28 measurements.

\begin{remark}[Neither census-complete rung is known to imply the other]
\label{p3-rem:independent}
Deriving Theorem~\ref{p3-thm:rung4census} from Theorem~\ref{p3-thm:rung3census} would require every order-29
Costas array to contain an order-28 Costas deletion minor; deriving the latter from
the former would require the converse extension property. Neither implication is
established. Theorem~\ref{p3-thm:heredity} measures that a single corner deletion
preserves the Costas property only $14.6\%$ of the time at order 28, showing the
required deletion property is far from automatic and therefore providing no basis
for deriving either rung theorem from the other. No counterexample is exhibited;
the point is that neither present proof implies the other, since they close different extension channels --- extension from an order-29 core and extension from an
order-28 core.
\end{remark}


\begin{remark}[Canonical mode is exact]
Restricting cores to one representative per dihedral orbit is exact for
existence and for the minimum defect, because the configuration space is
$D_4$-equivariant and the defect is $D_4$-invariant, so every configuration maps
to one with a canonical core and the same defect. It is \emph{not} exact for
counting solutions.
\end{remark}

\section{Heredity of the Costas property}
\label{p3-sec:heredity}

The value of deepening the ladder depends on how often a Costas array contains a
smaller one. With a complete census at order 28 this is directly measurable.
Deleting $k$ columns forces deletion of the $k$ rows they occupy; the remaining
dots are compressed and tested.

\begin{center}
\begin{tabular}{@{}rrrrr@{}}
\toprule
$k$ & sub-order & deletions per array & arrays with a Costas sub-array & \\
\midrule
1 & 27 & 28 & 104 / 712 & \textbf{14.6\%} \\
2 & 26 & 378 & 24 / 712 & \textbf{3.4\%} \\
3 & 25 & 3{,}276 & 0 / 712 & \textbf{0.0\%} \\
4 & 24 & 20{,}475 & 0 / 120 sampled & 0.0\% \\
5 & 23 & 98{,}280 & 0 / 120 sampled & 0.0\% \\
\bottomrule
\end{tabular}
\end{center}

\begin{theorem}[Heredity is corner-driven]
\label{p3-thm:heredity}
Over the complete order-28 census, exactly 104 of 712 arrays possess a corner
dot ($14.6\%$, against a random-permutation corner rate of $4/28 = 14.3\%$). Of
the 112 order-27 Costas sub-arrays found, \textbf{104 arise from corner
deletion} and only 8 are non-corner coincidences.
\end{theorem}

\noindent The table counts arrays with at least one Costas sub-array; the theorem counts
sub-arrays. A corner deletion of a Costas array is always Costas, so the 104 arrays with a corner
dot are exactly the 104 arrays of the table, and the 8 non-corner sub-arrays come from among them.

The explanation is immediate: corner deletion induces a translation, which
preserves all displacements, whereas any interior deletion induces a compression,
which does not.

\begin{heuristic}[Implication for rung prioritisation]
\label{p3-cor:ladder-value}
If the deletion-minor frequencies at order 32 resemble those measured in the
order-28 census, then an order-32 Costas array would be much more likely to
possess an order-31 or order-30 Costas minor than an order-28 one. This motivates
prioritising small $k$. The $k=1$ and $k=2$ searches are exhaustive relative to
the generated order-30 catalogue; $k=3$ and $k=4$ are census-complete at their
respective core orders. The measurement lowers the expected value of $k\ge5$ but
does not prove those channels empty.
\end{heuristic}

This is the single most consequential measurement in this investigation. It caps the
reach of branch (ii) of the three-branch decomposition, and correspondingly
enlarges branch (iii).


%=====================================================================

\section{Status of the claims}\label{p3-sec:status}
The results are computational and uncertified. Their universal quantifiers rest on the
completeness of the published enumerations of orders 28 and 29, an external premise. The generated
order-28 catalogue is shown equal to the published census; the order-29 catalogue is the published
census, reconstructed and verified as described in the paper on the order-29 census, where its
representatives are also kernel-checked in Lean. The order-28 search was run as 89 deterministic
work units and merged by a separate fail-closed tool. We have not found these results in the
literature, but a systematic search has not been completed.

\paragraph{Relation to prior work.} Apparently new computational results. Their novelty has not been established by a systematic literature search.
\begin{paperrefs}
\bibitem{p3-drakakis28} K.~Drakakis, F.~Iorio and S.~Rickard, \emph{The enumeration of Costas arrays of order 28 and its consequences}, Advances in Mathematics of Communications, 2011.
\bibitem{p3-drakakis29} K.~Drakakis, F.~Iorio, S.~Rickard and J.~Walsh,
\emph{Results of the enumeration of Costas arrays of order 29}, IEEE Trans.\
Inform.\ Theory \textbf{57} (2011).
\end{paperrefs}

\clearpage
\chapter{Reconstructing and Verifying the Order-29 Costas Census}\label{p4-paper}
\paperauthor
\paperabstract{The published enumeration of order-29 Costas arrays lists 23 lexicographically minimal
equivalence-class representatives. We reconstruct the full census from them: every representative
validates as an order-29 Costas permutation, and expansion under $D_4$ reproduces all published
figures exactly, namely 164 arrays in 23 classes, 5 of them symmetric, with matching per-method
counts. An independent algebraic generator produces 148 of the 164 arrays; the 16 it misses are
exactly the two classes produced by the Rickard--Welch operation, which corner surgery does not
reach. A Lean 4 development kernel-checks every representative and all eight of its $D_4$ images,
without \texttt{native\_decide}. The certificate covers verification, that every listed array is
Costas, not completeness, which remains the published enumeration's.}

\section{Introduction}\label{p4-sec:paperintro}
Several results in this volume quantify over every Costas array of order 29, and so depend on
having that census in hand and correct. The published enumeration~\cite{p4-drakakis29} gives only class representatives.
This paper reconstructs the full census, checks it against every published figure, cross-checks it
with an independent generator, and builds a machine-checked certificate for the part of the claim
that can be certified.

\paragraph{Dependencies.} This paper uses only the published class representatives. Paper~3 depends on it.

\section{Preliminaries}\label{p4-sec:notation}
A \emph{Costas array} of order $n$ is an $n\times n$ permutation array in which the displacement
vectors between pairs of dots are all distinct. Writing the array as a permutation $h$ of
$\{0,\dots,n-1\}$, with the dot of column $i$ in row $h_i$, the condition is that for each
horizontal separation $m$ the vertical differences $h_{i+m}-h_i$ are distinct. Costas arrays are
known at every order up to 31; orders 32 and 33 are the smallest for which existence is open. The
dihedral group $D_4$ of the square maps Costas arrays to Costas arrays, and arrays are counted
either individually or by $D_4$ orbit (equivalence class). Further notation and terms are collected
in the notation table and glossary at the front of this volume.

\section{Reconstruction from the published representatives}
\label{p4-sec:recon}
The complete order-29 census was reconstructed from Table 2 of the enumeration
paper, which lists the 23 lexicographically minimal equivalence-class
representatives with their construction method. Every representative was
independently validated as an order-29 Costas permutation and expanded under
$D_4$, reproducing all published figures exactly: \textbf{164 arrays, 23
equivalence classes, 5 symmetric}, with per-method counts
$G_0\,4$, $G_2\,64$, $G_3\,24$, $W_0\,24$, $W_2\,32$, $RW_0\,16$.


\paragraph{The generator gap, identified.}
The corrected algebraic generator produces 148 of the 164 arrays, and the
shortfall is now explained exactly: the 16 missing arrays form the 2 $RW_0$
equivalence classes produced by the semi-empirical Rickard--Welch method, which
appends a blank row to an exponential $W_1$ array, cyclically shifts, and adjoins
a column. That operation is not corner surgery, and the implemented
corner-deletion/addition closure does not reach its two classes. The gap is
therefore explained by a missing \emph{operation class} rather than by a defect
in the implemented closure; it does not exclude some unrelated genealogy for
those arrays through corner operations from a different seed. The generated set is a verified subset of the
census. This also confirms the order-29 taxonomy: no array of order 29 is
sporadic, but 2 classes lie outside $G\cup PE$ as the generator implements it.


\section{A machine-checked certificate in Lean}
\label{p4-sec:census29lean}

The census underwrites the universal quantifier in the order-29 deletion-minor theorem of Paper~3, and certifying it is the cheapest real certificate available to this investigation. Every constituent theorem has been kernel-checked in an individual compilation, in Lean 4.15.0
with no Mathlib, no \texttt{sorry} and no \texttt{native\_decide}. The released aggregate target
has not yet been built in one invocation, because the available 4\,GB environment was
insufficient.

\paragraph{What kernel-checks.} All 23 published equivalence-class
representatives, each as an individual theorem
\texttt{checkSeed29 census29SeedNN = true} discharged \texttt{by decide}. Each
verifies that representative \emph{and all eight of its $D_4$ images} as order-29
Costas permutations: length, range, permutation uniqueness and fixed-lag
difference uniqueness. Total elaboration is about 230\,s, roughly 10\,s per seed.
The two orbit-level theorems also compile individually: the expansion has
cardinality exactly 164 (10\,s), and every member of it is a Costas permutation
(120\,s).

\paragraph{Three findings about the mechanism.} First, per-seed chunking is what
makes kernel decision feasible at all. A single aggregate over the 23
representatives re-runs every seed inside one reduction and exceeds the
elaborator's budget; it is also redundant, since the per-seed theorems are the
certificate of record. Second, the binding constraint is \emph{heartbeats}, not
recursion depth: the orbit theorems fail at the default 200{,}000-heartbeat limit
regardless of \texttt{maxRecDepth}, and pass with
\texttt{set\_option maxHeartbeats 4000000}. Third, kernel reduction is
memory-hungry --- the complete module exhausts a 4\,GB container even when every
constituent theorem compiles alone, so the development is split into a seed
module and an orbit module. The orbit module needs a larger machine than was
available here; each of its theorems has compiled individually, but the module has not been built in one invocation.

\paragraph{The boundary, which matters more than the build.} This certifies the
supplied representatives and their 164-element orbit expansion. It does
\emph{not} prove that the catalogue is exhaustive. Exhaustiveness of the order-29
census is an external premise inherited from the published enumeration, and it is
precisely the premise on which that theorem's universal quantifier
rests. So the certificate closes the verification half --- every array claimed to
be Costas is one --- while the completeness half remains a citation. The presence of a Lean certificate should not be read as implying more than it delivers.

\paragraph{An earlier attempt, and why it failed.} A first version placed the
whole computation --- generation, $D_4$ closure, deduplication and Costas
checking over all 164 arrays --- under a single \texttt{by decide}. It returns
\texttt{true} instantly under \texttt{\#eval} and fails in the kernel, because
\texttt{decide} does not use the interpreter: it asks the kernel to reduce the
computation by \texttt{whnf}, with no compiled code and no fast paths through
\texttt{List}. Roughly $10^6$ comparisons is well past what that will do.
\texttt{native\_decide} would finish at once and is exactly the kernel-bypassing
route the choice of \texttt{decide} was meant to avoid. Agreement under
\texttt{\#eval} is the positive-arm test, not the certificate.



%=====================================================================

\section{Status of the claims}\label{p4-sec:status}
The reconstruction and its agreement with the published figures are computational. The Lean
development kernel-checks all 23 representatives with their $D_4$ images; every constituent theorem, including the orbit-level ones, has been kernel-checked in an individual compilation, but the released aggregate target has not been built in one invocation, because the available 4\,GB environment was insufficient. Exhaustiveness of the census is not certified: it is a premise
inherited from the published enumeration.

\paragraph{Relation to prior work.} A verification of a published census; the reconstruction and the Lean certificate are this project's.
\begin{paperrefs}
\bibitem{p4-drakakis29} K.~Drakakis, F.~Iorio, S.~Rickard and J.~Walsh,
\emph{Results of the enumeration of Costas arrays of order 29}, IEEE Trans.\
Inform.\ Theory \textbf{57} (2011).
\end{paperrefs}

\clearpage
\chapter{Telescoping-Sum Propagation for Costas Search: A Sound Rule, a Frozen Benchmark and an Operational Failure}\label{p5-paper}
\paperauthor
\paperabstract{For any sequence, the sum of its differences at lag $m$ telescopes to the sum of its last $m$
terms minus the sum of its first $m$. In a Costas array the differences at a fixed lag are distinct,
which turns the identity into an inference rule for partial assignments. We prove the rule sound and
evaluate it under a protocol frozen in advance. On the complete benchmark of Costas-valid
six-leading-column prefixes at order 12, it detects $117{,}880$ infeasible prefixes against
$77{,}074$ for the baseline, for about $5\%$ extra measured work, and rejects none of the $7{,}632$
completable prefixes. In an exact search at orders 11--13 it leaves the solution sets byte-identical
and reduces the search tree to $68$--$74\%$ of the baseline's, passing the scientific gate, but runs
$12$--$18\%$ slower, failing the operational gate.}

\section{Introduction}\label{p5-sec:paperintro}
Exhaustive Costas search, the method behind the complete enumerations through order 29~\cite{p5-drakakis27,p5-drakakis28,p5-drakakis29}, lives or dies by propagation: how much it can deduce from a partial
array before branching. This paper develops one new rule, based on a telescoping identity, proves it
sound, and tests it in two stages under criteria fixed before the results were seen: first on a
complete benchmark population, then inside an exact search. It also records two implementation bugs
that produced an earlier wrong verdict, because how they were caught is itself a methodological
point.

\paragraph{Dependencies.} This paper does not depend on the others.

\section{Preliminaries}\label{p5-sec:notation}
A \emph{Costas array} of order $n$ is an $n\times n$ permutation array in which the displacement
vectors between pairs of dots are all distinct. Writing the array as a permutation $h$ of
$\{0,\dots,n-1\}$, with the dot of column $i$ in row $h_i$, the condition is that for each
horizontal separation $m$ the vertical differences $h_{i+m}-h_i$ are distinct. Costas arrays are
known at every order up to 31; orders 32 and 33 are the smallest for which existence is open. The
dihedral group $D_4$ of the square maps Costas arrays to Costas arrays, and arrays are counted
either individually or by $D_4$ orbit (equivalence class). Further notation and terms are collected
in the notation table and glossary at the front of this volume.


The telescoping identity $\sum_i(h_{i+m}-h_i)=(\text{sum of last }m)-(\text{sum
of first }m)$ couples the endpoint sums of a partial assignment to its
lag-$m$ difference multiset. This section records the inference rule it yields,
a proof that the rule is sound, and its evaluation under a prospectively frozen,
cost-controlled protocol on complete finite populations. It is the one
prospectively gated positive result in this investigation.

\section{Benchmark construction, and a problem with the original gate}
Regenerating the order-12 infeasibility benchmark from scratch --- enumerate all
Costas-valid assignments of the first $s$ columns, then test by exhaustive search
whether each extends to a full order-12 Costas array:

\begin{center}
\begin{tabular}{@{}rrrr@{}}
\toprule
$s$ & Costas-valid prefixes & infeasible & rate \\
\midrule
4 & 10{,}112 & 5{,}120 & 50.63\% \\
5 & 64{,}000 & 56{,}898 & 88.90\% \\
6 & 304{,}288 & \textbf{296{,}656} & 97.49\% \\
7 & 987{,}664 & 979{,}894 & 99.21\% \\
\bottomrule
\end{tabular}
\end{center}

\noindent Among the $304{,}288$ Costas-valid six-leading-column prefixes at
order 12, $296{,}656$ --- about $97.49\%$ --- admit no completion. Local Costas
validity at depth six is therefore a weak predictor of global extendibility. The
benchmark is precisely the complete population of infeasible Costas-valid
six-\emph{leading-column} prefixes at order 12; results on it do not generalise
to arbitrary six-cell partial assignments.

\begin{observation}[The predeclared Stage~1 gate is underpowered]
\label{p5-obs:gate}
The recorded benchmark of ``296 infeasible configurations'' is a \emph{sample},
three orders of magnitude smaller than the $296{,}656$ available at $s=6$. With
$m=296$ draws and an observed union detection rate of $96/296=32.4\%$, the
binomial standard error is
\[
\sqrt{296\cdot0.324\cdot0.676}\;\approx\;8.1,
\]
so the gate at 106 detections sits about $1.25$ standard errors above the
baseline. A propagator with \emph{identical} power would clear it with probability $0.12$, the exact
binomial tail $\Pr(X\ge106)$ for $X\sim\mathrm{Bin}(296,96/296)$: roughly one time in eight. The gate as frozen cannot distinguish a real improvement from
sampling noise.
\end{observation}

The repair is to run against the complete finite population and preserve the
relative effect size encoded by the original gate. The revised gate therefore
requires
\[
|D_T|\;\ge\;\bigl\lceil 1.104167\,|D_B|\bigr\rceil,
\]
together with zero false positives and a cost ratio of at most two. This is an
exact finite-population comparison, not a significance test: on the complete
population the detection rate is the exact quantity $r_B=|D_B|/296{,}656$, and a
standard error would re-enter only if these prefixes were treated as a sample
from a wider family of orders, depths or geometries. A caveat on scope: the
reconstruction here fixes whole leading columns, and if the original benchmark
fixed arbitrary cells the population will differ. The power calculation does not
depend on which definition is used; a 296-element sample is underpowered either
way.


\section{The frozen protocol and the result}
\emph{Audit trail.} An earlier run reported a null failure. It was wrong, through
two implementation defects recorded here as findings rather than footnotes. (i)
The singleton-shaving routine rebuilt each domain from the unary filter on entry,
overwriting the tentative singleton, so shaving never fired; the conclusion that
``shaving contributes nothing'' was an artefact, and the weakened baseline made
the $T_m$ margin vanish. (ii) Relations between two fixed prefix positions were
filtered against the used-difference mask, which contains that pair's own
difference, so every such relation emptied and RA rejected every input. The
second was caught immediately by the complete positive arm. The methodological
conclusion is general: \emph{testing only known-infeasible configurations can
measure detection power but cannot establish even empirical soundness.}

\begin{lemma}[Soundness of the $T_m$ inference rule]
\label{p5-lem:tmsound}
Fix a partial assignment with sound position-value domains $D_i$, and a lag $m$.
Let $k$ be the number of pairs $(i,i+m)$ not wholly inside the prefix, $U$ the
union of the corresponding difference domains, $\det$ the sum of the determined
differences at lag $m$, and $R$ the remaining value set. If $U$ has fewer than $k$ values, reject: the remaining
differences at lag $m$ would have to be $k$ distinct values drawn from $U$. Otherwise put
\[
[\ell,h]=\Bigl[\det+\!\!\sum_{k \text{ smallest of } U}\!\!,\ \
\det+\!\!\sum_{k \text{ largest of } U}\Bigr],
\]
and let $[\tau_-,\tau_+]$ be the corresponding interval for
$T_m=(\text{sum of last }m)-(\text{sum of first }m)$, obtained from the extremal
$f_R$- and $l_R$-subset sums of $R$. If $[\ell,h]\cap[\tau_-,\tau_+]=\varnothing$
then the partial assignment has no Costas completion.
\end{lemma}
\begin{proof}
Let $\sigma$ be any Costas completion. Telescoping gives the identity
$\sum_{i=0}^{n-1-m}\bigl(\sigma_{i+m}-\sigma_i\bigr)
=\sum_{j=m}^{n-1}\sigma_j-\sum_{i=0}^{n-1-m}\sigma_i
=\bigl(\textstyle\sum\text{last }m\bigr)-\bigl(\sum\text{first }m\bigr)=T_m(\sigma)$,
valid for every sequence. Since $\sigma$ is Costas, its $n-m$ differences at lag
$m$ are pairwise distinct, so the $k$ undetermined ones are $k$ distinct integers,
each lying in its domain and hence in $U$ by soundness of the domains. Any $k$
distinct elements of a finite set have sum between that of its $k$ smallest and
its $k$ largest, so the left side lies in $[\ell,h]$. The unknown parts of the two
endpoint sums are subset sums of $R$ of fixed sizes, each between the extremal
subset sums of that size, so $T_m(\sigma)\in[\tau_-,\tau_+]$. As the two sides are
equal, the intervals meet. Contrapositive. Both interval constructions are
outer approximations: they may include sums incompatible with the individual
variable domains or with joint row usage, but they exclude no sum realised by a
completion, so disjointness is a valid certificate and can never be manufactured.
\end{proof}

\noindent The rule is sound but incomplete. The interval $[\ell,h]$ ignores which values can go to
which pairs, so the $k$ extremal values need not be simultaneously assignable, and the endpoint
interval ignores the joint feasibility of the row placements. In outline, the test at one lag is:
\begin{enumerate}[nosep]
\item collect the $k$ unresolved pairs at lag $m$, the union $U$ of their difference domains, the sum
  of the determined differences, and the remaining value set $R$;
\item if $|U|<k$, reject;
\item form $[\ell,h]$ from the $k$ smallest and $k$ largest values of $U$, and $[\tau_-,\tau_+]$ from
  the extremal subset sums of $R$ for the first and last $m$ positions;
\item reject if the two intervals are disjoint; otherwise pass.
\end{enumerate}

\subsection*{Benchmark}\noindent $\mathcal B_{n,s}$, every Costas-valid
$s$-leading-column prefix at order $n$, labelled by exhaustive completion. At
$(12,6)$: $304{,}288$ valid, $296{,}656$ infeasible, $7{,}632$ completable,
independently reproduced by a second completion procedure.
\textbf{Baselines.} $D_{\mathrm{RA}}$, relation-algebra closure by path
consistency on difference-relation masks under sumset composition; $D_{L_3}$,
position--value domains with unary filtering, arc consistency, Hall matching and
singleton shaving; $D_B=D_{\mathrm{RA}}\cup D_{L_3}$.
\textbf{Candidate.} $D_T=D_B$ plus the rule of Lemma~\ref{p5-lem:tmsound}, so
$D_B\subseteq D_T$ and the gate reduces to marginal coverage.
\textbf{Gate, frozen before the run:} $|D_T|\ge\lceil1.104167|D_B|\rceil$, zero
false positives, cost ratio $\le2$ on the \emph{declared deterministic propagation-work
counter}: \texttt{workRA} increments per relation initialisation, per sumset
shift and per composition; \texttt{workL3} per unary-filter test, per
arc-consistency support test and per matching call; \texttt{workTM} per domain
scan entry and per interval test. It is a frozen instrumentation metric, not a
literal instruction count or a hardware-independent measure of total cost.

\begin{theorem}[Stage~1, computational]
\label{p5-thm:f3}
Among the $296{,}656$ infeasible Costas-valid six-leading-column prefixes of order
12, the frozen baseline $D_B$ detects $77{,}074$. Adding the $T_m$-coupled rule
detects $117{,}880$, including $40{,}806$ additional infeasible prefixes, at
$1.050$ times the baseline declared propagation-work count. Neither system
rejects any of the $7{,}632$ completable prefixes.
\end{theorem}

The frozen threshold was $\lceil1.104167\times77{,}074\rceil=85{,}103$; the
candidate exceeds it by $32{,}777$ detections, an improvement of
$40{,}806/77{,}074\approx52.94\%$ against a required $10.4167\%$, for about $5\%$
additional work. \textbf{The gate passes on all three criteria.}

\medskip\noindent External validation on further complete populations (the formal
gate was declared at $(12,6)$; these rows are exploratory generalisation):

\begin{center}
\begin{tabular}{@{}rrrrrrl@{}}
\toprule
$n$ & $s$ & $|D_B|$ & required & $|D_T|$ & cost & verdict \\
\midrule
10 & 5 & 3{,}882 & 4{,}287 & 7{,}388 & 1.060 & pass \\
10 & 6 & 47{,}230 & 52{,}150 & 54{,}104 & 1.017 & pass \\
11 & 6 & 88{,}652 & 97{,}887 & 119{,}734 & 1.023 & pass \\
11 & 7 & 325{,}922 & 359{,}873 & 341{,}520 & 1.008 & gate inapplicable$^\ast$ \\
12 & 7 & 875{,}236 & 966{,}407 & 969{,}828 & 1.006 & pass \\
\bottomrule
\end{tabular}
\end{center}

\noindent$^\ast$\emph{Ceiling saturation}, not failure: with
$|D_B|=325{,}922$ and $|\mathcal B|=342{,}010$ the maximum possible relative
improvement is $4.94\%$, so a $10.4167\%$ gate is arithmetically impossible. A
multiplicative total-coverage gate is inapplicable whenever the baseline rate
exceeds $1/1.104167\approx90.6\%$. For such regimes the appropriate secondary
criterion is miss reduction,
$\rho_{\mathrm{miss}}=|D_T\setminus D_B|/|\mathcal B\setminus D_B|$; here
$15{,}598/16{,}088\approx96.95\%$ of the baseline's remaining misses are
eliminated. This metric is recommended for future high-baseline benchmarks and is
\emph{not} used retroactively --- $(12,6)$ passed its own frozen gate.

\begin{observation}[Mechanistic reading]
\label{p5-obs:f3}
The $T_m$ rule supplies substantial marginal propagation beyond the implemented
RA and $L_3$ procedures across multiple complete leading-prefix benchmarks: at the
declared benchmark it improves detection by about $52.94\%$ over baseline for
about $5.0\%$ additional deterministic work. The $40{,}806$ marginal detections
are certificates the implemented relation-algebra closure and $L_3$ singleton
shaving do not produce, so the contribution cannot be explained as a different
execution order for those two procedures. On the tested leading-prefix
benchmarks, RA contributes no detections beyond $L_3$
($D_{\mathrm{RA}}\setminus D_{L_3}=\varnothing$ at every point), so the working
stack there is $L_3+T_m$; this is benchmark-specific and does not generalise.
\end{observation}

\noindent\textbf{What Stage~1 establishes:} a correctly implemented new propagation
mechanism, sound by Lemma~\ref{p5-lem:tmsound} and exhaustively validated on the
complete positive arm; substantial marginal detection beyond the frozen baseline;
reproducibility across several complete finite populations; a favourable
coverage-to-cost ratio. \textbf{What it does not:} improved total runtime in a
full order-32 search; the same marginal gain at order 32; that RA is always
subsumed by $L_3$; or that no other domain-consistency method derives the same
contradictions --- that would need a formal expressiveness theorem.



\section{Stage~2: does the rule reduce an exact search?}
\label{p5-sec:f3prime}

Stage~1 measured detection on static prefixes. Whether those deductions reduce the
tree of a running solver is a different question, and it was asked under a
separately frozen protocol. One executable selects the propagator by flag; every
non-propagation choice --- branching order, value order, symmetry handling,
incremental displacement check, thread count --- is identical across arms.
Relation algebra is omitted from the candidate, since it contributed no
detections beyond $L_3$ on every tested benchmark. The invocation schedule
(propagate when at most 8 positions are free) was \emph{declared a priori and
never tuned}, so no order is a development order and all three are held out.

\begin{center}
\begin{tabular}{@{}rrrrrrrrr@{}}
\toprule
$n$ & solutions & $N_I$ & $N_{L_3}$ & $N_{L_3+T_m}$ & $R_{\text{tree}}$ & $W$ ratio & $t_{L_3}$ & $t_{L_3+T_m}$ \\
\midrule
11 & 4{,}368 & 1{,}437{,}514 & 379{,}290 & 259{,}414 & $0.316$ & $1.055$ & 5.0\,s & 5.9\,s \\
12 & 7{,}852 & 6{,}977{,}361 & 1{,}522{,}289 & 1{,}119{,}809 & $0.264$ & $1.047$ & 31.8\,s & 36.8\,s \\
13 & 12{,}828 & 33{,}831{,}238 & 5{,}785{,}132 & 4{,}175{,}962 & $0.278$ & $1.037$ & 174.4\,s & 197.1\,s \\
\bottomrule
\end{tabular}
\end{center}

\noindent Solution counts match the published censuses. Set identity is
established \emph{exactly}, not by hashing: each mode emits every solution in
canonical form, and the sorted files are compared byte for byte. At orders 11,
12 and 13 all three modes produce identical files
(\texttt{sha256} prefixes \texttt{3b346690b48c2527}, \texttt{2baed7d03cc6f9de},
\texttt{b6cf5cfc1580bff2} for 4{,}368, 7{,}852 and 12{,}828 solutions). The two
order-independent 64-bit hashes are retained as a fast regression guard;
aggregate hashes are not collision-free certificates of set equality, so the
claim rests on the exact comparison. The fixed column-by-column DFS also
generates each full assignment at most once, so duplicates are impossible by
construction.

\begin{observation}[Stage~2 verdict: structural pass, operational failure]
\label{p5-obs:f3prime}
Against the gate frozen before the run: \textbf{(a)} identical solution sets ---
pass; \textbf{(b)} $N_{L_3+T_m}\le0.90N_{L_3}$ --- pass, measured $0.684$,
$0.736$, $0.722$; \textbf{(c)} $W_{L_3+T_m}\le1.25W_{L_3}$ --- pass, measured
$1.055$, $1.047$, $1.037$. The \emph{scientific} gate passes at every order.
\textbf{(d)} $t_{L_3+T_m}<t_{L_3}$ on a majority --- \textbf{fails}, 0 of 3; the
candidate is $12$--$18\%$ slower in wall time. The \emph{operational} gate fails.
$T_m$ removes a quarter to a third of the remaining search tree for about $4\%$
additional deterministic work, and the current implementation does not convert
that into elapsed time.
\end{observation}

The larger finding is about the baseline, not the candidate. At order 13,
$L_3+T_m$ explores only $12.3\%$ of the incremental-only tree
($4{,}175{,}962$ against $33{,}831{,}238$ nodes) and yet takes $197$\,s against
$3.1$\,s. Across the three orders incremental-only explores $5.5$ to $8.1$ times
more nodes and runs $53$--$63$ times faster, because the exact bitmask displacement
check costs almost nothing per node while shaving costs a great deal. At these
orders heavyweight propagation is an operational loss \emph{in either
configuration}, and $T_m$ inherits that. The node-reduction trend is stable
rather than growing with order ($0.316$, $0.264$, $0.278$), so nothing here
supports projecting a crossover; establishing one would require either a far
cheaper $L_3$, a sparser invocation schedule, or orders well beyond those
enumerable here.

The honest reading: $T_m$ is a correct and genuinely useful deduction rule that
pays for itself in \emph{deductions per unit of declared work} and does not pay
for itself in seconds against this baseline.


\section{The representation boundary}
\label{p5-sec:repboundary}

Lemma~\ref{p5-lem:tmsound} is proved for any value of the formal contract. It does
not follow that \texttt{f3\_complete.cpp} builds one. That gap is now an object
in the development rather than a paragraph in this investigation.

\paragraph{Proved.} The bitmask/list bridge. The executable stores position
domains as \texttt{unsigned D[u]}, bit $k$ meaning ``the $k$-th remaining value is
still available at position $u$''; the Lean development uses lists of indices.
\texttt{domainList} and \texttt{mem\_domainList} identify the two exactly, and
\texttt{mem\_domainList\_of\_submask} shows that every filtering step --- each is
\texttt{D[u] \&= \~{}(1u <\!< k)}, producing a submask --- only shrinks a domain.
Pool membership proved against the \emph{filtered} domains therefore reads back
as membership in the original.

\paragraph{Not proved.} \texttt{CppObligations} is a structure whose eight fields
are the complete residue, each documented against the construct it refers to:

\begin{center}
\begin{tabular}{@{}ll@{}}
\toprule
obligation & claim about the C++ \\
\midrule
\texttt{domains\_sound} & \texttt{l3Detect} clears a bit only for a value no completion can take \\
\texttt{pairs\_enumerated} & the lag loop visits each unresolved pair once; \texttt{a < s} classifies it \\
\texttt{pools\_match} & the two \texttt{push\_back} branches generate the two pool functions \\
\texttt{dedup\_matches} & \texttt{sort}/\texttt{unique} has the membership of \texttt{dedupInt} \\
\texttt{prefix\_sums\_extremal} & \texttt{pre\_lo}, \texttt{pre\_hi} are the extremal $k$-subset sums \\
\texttt{endpoints\_match} & \texttt{slo}, \texttt{shi}, \texttt{tlo}, \texttt{thi} are the interval endpoints \\
\texttt{cardinalities\_correct} & \texttt{fk}, \texttt{fR}, \texttt{lk}, \texttt{lR} are the true window sizes \\
\texttt{no\_overflow} & no accumulation overflows; all sums are $O(n^2)$ \\
\bottomrule
\end{tabular}
\end{center}

\noindent An opaque \texttt{Prop} field, however, is an unfalsifiable
placeholder: it can be neither proved nor refuted, and a differential harness can
only ever refute. \texttt{Reference.lean} therefore supplies a Lean reference
implementation of the computation \texttt{tmDetect} performs and discharges the
shallow obligations as theorems about it:

\begin{center}\footnotesize
\begin{tabular}{@{}ll@{}}
\toprule
former obligation & now \\
\midrule
\texttt{pools\_match} & \texttt{ref\_fixedFree\_eq}, \texttt{ref\_freeFree\_eq} (by \texttt{rfl}, no axioms) \\
\texttt{dedup\_matches} & \texttt{ref\_dedup\_preserves}, \texttt{ref\_dedup\_nodup} \\
\texttt{prefix\_sums\_extremal} & \texttt{ref\_lower\_eq} (by \texttt{rfl}, no axioms) \\
\texttt{endpoints\_match} & \texttt{ref\_disjoint\_iff} \\
\texttt{no\_overflow} & \texttt{ref\_sum\_bounded}: sum bounded by $\text{length}\times n$, so $O(n^2)$ \\
\texttt{pairs\_enumerated}, \texttt{cardinalities\_correct} & absorbed below \\
\bottomrule
\end{tabular}
\end{center}

\noindent What survives is \texttt{CppResidue}, with \textbf{two} fields rather
than eight. \texttt{domains\_sound} --- that \texttt{l3Detect} clears a bit only
for a value no completion can take --- is the deep one, and is not a
representation question at all: it asserts that arc consistency and singleton
shaving are correct filters. \texttt{cpp\_matches\_reference} is a single
refinement claim about straight-line code, and is precisely what a differential
harness \emph{can} test on the complete order-12 population.

\begin{remark}
The reduction matters more than the count. Eight opaque propositions were an
unbounded liability; two typed contracts, one of them testable and the other a
named mathematical property of a filtering algorithm, are a work item. Replacing
placeholders with a reference implementation is what made the difference ---
five obligations became theorems because there was finally something to prove
them about.
\end{remark}



\section{Status of the claims}\label{p5-sec:status}
Soundness of the rule is proved. The benchmark and search results are computational, with
solution-set identity established byte for byte. The rule passed its scientific gate and failed its
operational gate; nothing here supports projecting a crossover at higher orders.

\paragraph{Relation to prior work.} A new inference rule as far as we know; its novelty has not been established by a systematic literature search.
\begin{paperrefs}
\bibitem{p5-drakakis27} K.~Drakakis, S.~Rickard, J.~K.~Beard, R.~Caballero,
F.~Iorio, G.~O'Brien and J.~Walsh, \emph{Results of the enumeration of Costas
arrays of order 27}, IEEE Trans.\ Inform.\ Theory \textbf{54} (2008).
\bibitem{p5-drakakis28} K.~Drakakis, F.~Iorio and S.~Rickard, \emph{The enumeration of Costas arrays of order 28 and its consequences}, Advances in Mathematics of Communications, 2011.
\bibitem{p5-drakakis29} K.~Drakakis, F.~Iorio, S.~Rickard and J.~Walsh,
\emph{Results of the enumeration of Costas arrays of order 29}, IEEE Trans.\
Inform.\ Theory \textbf{57} (2011).
\end{paperrefs}

\clearpage
\chapter{A Conditional First-Moment Heuristic for Unexplained Costas Classes}\label{p6-paper}
\paperauthor
\paperabstract{The published count of sporadic Costas arrays, those not explained by a fixed catalogue of constructions, first reaches zero at order 28 and remains zero at order 29. We derive a first-moment
model whose constraint cost, about $n^2/12$, is fixed before fitting. It matches the published
sporadic counts to about $2\%$, tracks the curve across four orders of magnitude, and places
extinction at order 27 against an observed 28. Extrapolated, it forecasts about $10^{-4}$ sporadic
equivalence classes at order 32, where the algebraic term is empty. The forecast is conditional: it
rests on superlinear growth of the constraint cost, which the data do not establish. A cubic fits the
same eight points better and diverges, and under linear growth at the terminal slope the probability of no further sporadic arrays, read heuristically as a Poisson probability, falls from $0.97$ to $0.54$.}

\section{Introduction}\label{p6-sec:paperintro}
Every attack on order 32 in this volume fails or closes only a structured subclass. A different
question is whether there is any reason to expect an order-32 array at all. This paper approaches it populationally, using the published counts~\cite{p6-drakakis27,p6-drakakis28,p6-drakakis29,p6-oeis}, and in the spirit of the result that the density of Costas arrays decays exponentially~\cite{p6-warnke}: because ``sporadic'' is defined relative to the construction catalogue used in the published classification, the quantity modelled is the number $S(n)$ of order-$n$ classes not explained by that fixed catalogue. The paper models how $S(n)$ declines with order, and asks what the
model predicts at order 32, and how much weight the prediction can bear. The constraint cost of about $n^2/12$ is a scaling heuristic obtained from an independence approximation; differences of a random permutation are not independent, and the paper returns to this.

\paragraph{Dependencies.} This paper uses the construction taxonomy of Paper~1 and the published counts. No other paper depends on it.

\section{Preliminaries}\label{p6-sec:notation}
A \emph{Costas array} of order $n$ is an $n\times n$ permutation array in which the displacement
vectors between pairs of dots are all distinct. Writing the array as a permutation $h$ of
$\{0,\dots,n-1\}$, with the dot of column $i$ in row $h_i$, the condition is that for each
horizontal separation $m$ the vertical differences $h_{i+m}-h_i$ are distinct. Costas arrays are
known at every order up to 31; orders 32 and 33 are the smallest for which existence is open. The
dihedral group $D_4$ of the square maps Costas arrays to Costas arrays, and arrays are counted
either individually or by $D_4$ orbit (equivalence class). Further notation and terms are collected
in the notation table and glossary at the front of this volume.


The investigation's negative results all share a shape: the methods fail because
order 32 is not structurally distinguished. This section argues that the
distinguishing fact is populational rather than structural, and that it is
quantitative.

\section{The four-class taxonomy}

Drakakis, Iorio and Rickard partition Costas arrays into four classes:
\emph{generated} (G), obtained by algebraic constructions whose applicability
depends only on the order; \emph{predictably emergent} (PE), arising from
transformations of generated arrays and guaranteed by additional conditions;
\emph{unpredictably emergent} (UE), arising from such transformations but
requiring direct verification; and \emph{sporadic} (S), everything else.

This sharpens language used loosely earlier in this investigation. The catalogue generator of Paper~2 produces $G\cup PE$ and part of $UE$; the
16 arrays of order 29 it cannot reach are \emph{emergent}, not sporadic --- they
were found previously by a semi-empirical technique. Order 29 has \emph{zero}
sporadic equivalence classes, and order 28 is the first order above 5 with none.

\section{The published sporadic curve}

Table 1 of the order-28 enumeration gives sporadic equivalence-class counts:

\begin{center}
\begin{tabular}{@{}lrrrrrrrrr@{}}
\toprule
order & 13 & 14 & 17 & 19 & 24 & 25 & 26 & 27 & 28 \\
\midrule
total ECs & 1616 & 2168 & 2294 & 1283 & 25 & 12 & 8 & 29 & 89 \\
sporadic ECs & 1602 & 2163 & 2270 & 1281 & 24 & 5 & 2 & \textbf{1} & \textbf{0} \\
\bottomrule
\end{tabular}
\end{center}

Sporadics dominate the middle orders, collapse through the twenties, and vanish
at 28. The extinction conjecture --- that beyond some $N$ only generated and
predictably emergent arrays survive --- is theirs, argued from the order-28
result, from unsuccessful searches for new emergent arrays, and from the
$O(n^3)$ growth in the number of constraints.

\section{A first-moment derivation}

The conjecture can be made quantitative. For lag $d$ the $n-d$ differences
$h_{i+d}-h_i$ must be distinct among roughly $2n-2$ available values, a birthday
cost of about $(n-d)^2/(4n)$ in the log. Summing over $d$,
\[
\log \Pr[\,p \text{ is Costas}\,] \;\approx\; -\frac{1}{4n}\sum_{k=1}^{n-1}k^2
\;\approx\; -\frac{n^2}{12},
\]
against the entropy $\log n! \approx n\log n - n$. The populations cross where
$n\log n - n = n^2/12$, that is $n \approx 12(\log n - 1)$; at $n=29$ the right
side is $28.4$.

\begin{observation}[Quadratic collision cost fits the observed sporadic curve]
\label{p6-obs:extinction}
Fitting $\log S(n) = \log(n!/8) - f(n)$ with $f$ quadratic to the sporadic EC
counts above gives
\[
f(n) = 0.08501\,n^2 + 0.173\,n - 3.657,
\]
whose leading coefficient agrees with the derived $1/12 = 0.08333$ to within
$2\%$.
\end{observation}

\begin{center}
\begin{tabular}{@{}lrrrrrrrr@{}}
\toprule
order & 13 & 14 & 17 & 19 & 24 & 25 & 26 & 27 \\
\midrule
sporadic ECs & 1602 & 2163 & 2270 & 1281 & 24 & 5 & 2 & 1 \\
model & 1830 & 2171 & 1941 & 1032 & 25.5 & 8.3 & 2.4 & 0.60 \\
\bottomrule
\end{tabular}
\end{center}

The fit holds across four orders of magnitude, is essentially exact at $n=14$
(2171 against 2163), and places the extinction point --- where the model first
falls below one equivalence class --- at $n=27$, one order before the observed
extinction at 28.

\section{Consequences for order 32}

Extrapolating: $0.025$ sporadic ECs at order 29, $0.0042$ at 30, $0.0006$ at 31,
and
\[
\boxed{\;\approx 10^{-4}\ \text{sporadic equivalence classes at order }32.\;}
\]

The contribution from the recognised direct Welch and Golomb construction
families is zero under the taxonomy and applicability conditions used here.
Orders 31 and 32 admit no direct construction: $W_1$ needs $n+1$ prime, $G_2$ needs $n+2$ a prime power, and for
$n=31,32$ none of $32,33,34,35,36$ supplies one in the required position
(Paper~1, Proposition~\ref{p1-prop:polydual}). Order 31 in particular has \emph{no} direct
construction --- $W_1$ would require $p=32$ prime --- yet it has arrays, and rung
$k=1$ identifies their origin exactly --- augmentation of order-30 cores, 8
arrays in one orbit. For order 32, several corresponding extension routes have
been exhausted: rung $k=2$ over all 664 \emph{generated} order-30 cores --- order
30 has never been enumerated, so this clause is catalogue-relative --- together
with the census-complete rungs $k=3$ and $k=4$ over every order-29 and every
order-28 Costas array respectively, the latter at $115{,}087{,}822{,}400$
placements. These close the specified augmentation and deletion-minor channels
but do not control an order-32 array possessing none of those minors.

So the expected order-32 population is $10^{-4}$ from the sporadic term plus zero
from the algebraic term. Conditional on the quadratic model, the projected population is so small that continued failure to
find an order-32 array is unsurprising. This is an interpretation of non-observation, not an
explanation independent of the model. On that reading, the structural methods examined in the project's consolidated report failed not because
they were too weak but because the fact they were searching for is populational.

\begin{remark}[A falsifiable prediction]
\label{p6-rem:falsifiable}
The model predicts that an order-30 enumeration will find \emph{only} generated
and emergent arrays --- $W_1(31)$ and $G_2(32)$ derivatives --- and zero
sporadic equivalence classes. A single sporadic array at order 30 would refute
the extinction model and reopen order 32. This raises the value of F4 from data
collection to hypothesis test.
\end{remark}

\begin{remark}[Honest error bar on the extrapolation]
\label{p6-rem:errorbar}
The point estimates above understate the uncertainty. The fit has residual
standard deviation $0.360$ in log space on 5 degrees of freedom, and the
extrapolation leverage at $n=32$ is $6.62$ against $0.46$--$0.50$ inside the data
range $13\le n\le27$. Proper 95\% prediction intervals are therefore wide:
$[6.1\times10^{-6},\,1.0\times10^{-3}]$ at $n=32$, and for the aggregate
$\sum_{n\ge29}\mathbb E[S(n)]$ a point estimate of $0.030$ with a band of roughly
$[0.007,\,0.14]$. If, additionally and heuristically, the class counts at successive orders are treated as
independent Poisson variables with these means, the probability that no sporadic array ever
appears again is $97\%$ as a point estimate but $87\%$--$99\%$ as an interval. Leave-one-out refits are stable (order-32 predictions from
$7.6\times10^{-5}$ to $1.6\times10^{-4}$), but leave-one-out tests interpolation
at leverage $0.5$ and says nothing about extrapolation at leverage $6.6$.
\end{remark}

\begin{remark}[Extinction versus completeness]
The natural theorem here is not $S_n=0$ for large $n$. ``Sporadic'' denotes what
present knowledge fails to explain, so that statement is contingent on the state
of the literature rather than on mathematics. The well-formed version is a
\emph{completeness} theorem: every Costas array of sufficiently large order
arises from an explicit list of constructions. That is the analogue of the
classification of finite simple groups, and no structural theory of the sporadic
class exists to support it.
\end{remark}

\section{Model-form adequacy, and where the forecast's weight actually sits}
\label{p6-sec:adequacy}

The interval in Remark~\ref{p6-rem:errorbar} accounts for estimation uncertainty
\emph{conditional on the quadratic model}. It says nothing about the model form.
With eight observations that distinction is decisive.

\paragraph{Stress test.} Fitting a cubic to the same data:

\begin{center}
\begin{tabular}{@{}lrrrrr@{}}
\toprule
model & RSS & resid.\ sd & $S(30)$ & $S(32)$ & $\sum_{n\ge29}$ \\
\midrule
quadratic & 0.649 & 0.360 & $4.2\times10^{-3}$ & $7.8\times10^{-5}$ & $0.030$ \\
cubic & \textbf{0.256} & \textbf{0.253} & $4.4\times10^{-2}$ & $8.8\times10^{-3}$ & diverges \\
\bottomrule
\end{tabular}
\end{center}

\noindent Information criteria do not settle the choice. With eight points, the small-sample
corrected AIC slightly favours the quadratic ($-8.1$ against $-6.2$), while BIC favours the cubic
($-19.2$ against $-13.9$). This is a model-adequacy stress test, not a competing forecast. The
cubic \emph{fits better} in sample and extrapolates to eventual explosion rather
than extinction; the two curves agree on the calibration range and diverge by
some 192 orders of magnitude outside it. The cubic is not rejected by goodness of
fit. It is rejected only relative to the proposed collision-cost mechanism, under
which $f(n)=-\log\Pr[\pi_n \text{ Costas}]$ should eventually increase and grow at
least quadratically.

\paragraph{That exclusion is a hypothesis, not a theorem.} It is tempting to
argue that adding a row adds constraints and removes none, so $f$ must be
increasing and convex. The argument fails: $\pi_n\sim\mathrm{Unif}(S_n)$ and
$\pi_{n+1}\sim\mathrm{Unif}(S_{n+1})$ are different probability spaces and are not
nested, deletion with compression does not preserve uniformity, and a superlinear
count of formal constraints does not imply convexity of their aggregate
logarithmic cost, since dependence among them also grows.

\paragraph{What the data say about the hypothesis.} Computing
$f(n)=\log n!-\log C(n)$ at every enumerated order:

\begin{center}
\begin{tabular}{@{}lr@{}}
\toprule
$\Delta f>0$ (monotone) & holds at 27 of 28 consecutive orders \\
$\Delta^2 f>0$ (convex) & holds at 20 of 27, failing at $n=4,8,10,22,25,26,27$ \\
\bottomrule
\end{tabular}
\end{center}

Monotonicity is well supported. \emph{Convexity is not}, and three of its seven
failures are the last three orders before the extrapolation begins. The
construction-free sporadic series shows the same thing: the slope of
$g(n)=\log(n!/8)-\log S(n)$ runs $+2.34,\,+2.76,\,+3.20,\,+3.88,\,+4.79$ and then
\emph{falls} to $+4.17$ and $+3.99$. Those terminal points carry counts of 5, 2
and 1, so Poisson noise of 45--100\% could account for it entirely; but the data
do not support convexity where it matters most.

\paragraph{Convexity is exactly what carries the conclusion.} Dropping it and
continuing $g$ linearly at its terminal slope $4.317$:

\begin{center}
\begin{tabular}{@{}lrrr@{}}
\toprule
model class for $g$ & decreasing phase & infinite horizon & $\Pr[\text{none}]$ on that phase \\
\midrule
superlinear (quadratic $f$) & $0.030$ & $0.030$ & $0.97$ \\
linear at terminal slope & $0.618$ through $n=73$ & \emph{diverges} & $0.54$ \\
concave (cubic fit) & model-dependent & \emph{diverges} & $\approx0$ \\
\bottomrule
\end{tabular}
\end{center}

\noindent The two descriptions must be kept apart. Under the linear model the
step ratio is $(n+1)e^{-s}$, which exceeds 1 once $n\ge74$ at $s=4.317$; the
figure $0.618$ is the expected count accumulated over the \emph{decreasing} phase
$28\le n\le73$ only, and $e^{-0.618}=0.539$ is the corresponding probability on
that finite window. Over the infinite horizon the series diverges, so a Poisson
reading gives probability zero of no further sporadic array. Extinction requires
$g$ to grow \emph{superlinearly}, and nothing in the enumeration establishes that:
under linear growth the sporadic population declines through
the thirties, bottoms out near $n=74$, and resumes growing.

\begin{observation}[Where the forecast's weight sits]
\label{p6-obs:weight}
The extinction forecast rests on two separable components: a structural
modelling assumption, that the effective logarithmic constraint cost is
eventually increasing and grows at least quadratically; and an empirical
calibration, which estimates the rate of that growth over the enumerated range.
The data support the calibration only after the model class has been chosen; they
do not establish the class. Within the class, $\Pr[\text{no further sporadic}]$ is
$0.97$ with an estimation band of $0.87$--$0.99$. Across plausible classes it
ranges from $0.97$ to below $0.54$.
\end{observation}

\paragraph{The constant is not pinned either.} Three levels should be
distinguished: the naive independent-uniform birthday calculation gives
$1/12=0.0833$; correcting for the approximately triangular difference law, whose
effective support is 47.2 slots at order 32 against 62 if uniform, raises it to
about $0.109$; the fitted value is $0.08501$. The gap between the second and third
shows that dependence or finite-order effects materially reduce the effective
collision cost relative to the refined independent model. It is not evidence for
any particular mechanism doing so. The derivation supports the quadratic scale and
order of magnitude far more convincingly than it identifies the constant.

\begin{remark}[The asymptotic theorem the argument needs]
\label{p6-rem:openproblem}
Extinction in the first-moment picture requires the entropy to be dominated:
$\mathbb E[S(n)]=(n!/8)e^{-f(n)}\to0$ summably, that is
\[
f(n)-\log(n!/8)\;\longrightarrow\;+\infty
\]
fast enough, with $f(n)\ge cn^2$ or $f(n)\ge(1+\varepsilon)n\log n$ as concrete
sufficient forms. Eventual monotonicity and eventual convexity are informative
intermediate results but neither implies extinction: a convex function may be
asymptotically linear, and linear $f$ is dominated by the entropy $n\log n$. A
theorem supplying entropy domination would be significant in its own right, is
the missing ingredient in every probabilistic statement above, and --- unlike the
order-32 existence question --- is uniform in $n$, so the uniformity caveat recorded in the project's consolidated report does not apply to it.
\end{remark}

The model is a first-moment heuristic, not a theorem. It says nothing about
whether a particular order-32 array exists, only what the expected number is
under a random model calibrated on nine measured orders. The quadratic scale and order of magnitude were motivated before fitting. The
fitted coefficient lies within $2\%$ of the naive independent-uniform value
$1/12$, but not of the triangular-support refinement near $0.109$; the agreement
should therefore not be read as identifying the asymptotic constant.



%=====================================================================

\section{Status of the claims}\label{p6-sec:status}
This is a conditional population forecast, not a proof. The model's form is chosen on stated
grounds and calibrated on published counts; its conclusion at order 32 depends on an asymptotic
growth assumption that the data do not establish, and an alternative fit of comparable quality
reverses it.

\paragraph{Relation to prior work.} A new heuristic model as far as we know; its novelty has not been established by a systematic literature search.
\begin{paperrefs}
\bibitem{p6-drakakis27} K.~Drakakis, S.~Rickard, J.~K.~Beard, R.~Caballero,
F.~Iorio, G.~O'Brien and J.~Walsh, \emph{Results of the enumeration of Costas
arrays of order 27}, IEEE Trans.\ Inform.\ Theory \textbf{54} (2008).
\bibitem{p6-drakakis28} K.~Drakakis, F.~Iorio and S.~Rickard, \emph{The enumeration of Costas arrays of order 28 and its consequences}, Advances in Mathematics of Communications, 2011.
\bibitem{p6-drakakis29} K.~Drakakis, F.~Iorio, S.~Rickard and J.~Walsh,
\emph{Results of the enumeration of Costas arrays of order 29}, IEEE Trans.\
Inform.\ Theory \textbf{57} (2011).
\bibitem{p6-oeis} OEIS A008404 (number of Costas arrays of order $n$) and A001441
(inequivalent Costas arrays under the dihedral group).
\bibitem{p6-warnke} L.~Warnke, B.~Correll and C.~Swanson, \emph{The density of
Costas arrays decays exponentially}, IEEE Trans.\ Inform.\ Theory \textbf{69}
(2023), 575--581.
\end{paperrefs}

\clearpage
\chapter{The Dyadic Defect Filtration and a Diagonal Parity Floor at Order 32}\label{p7-paper}
\paperauthor
\paperabstract{We study Costas defects through their 2-adic structure. For a permutation $p$ of order 32,
consider the 240 diagonal second differences $p_{r+2\ell}-2p_{r+\ell}+p_r$. We prove that at least
112 of them are even for every permutation, with equality exactly when exactly eight even positions
carry even values. This is the only unconditional universal bound the investigation produced. We also show that the zero-survivor case of the associated mod-4 minimisation is impossible, by reducing it to four complete graphs and exhausting signed cliques, although the minimum itself remains unknown, and we explain why neither result bears on existence. Finally, an explicit witness shows that unlabelled dyadic residue counts cannot characterise the Costas property at order 32, which closes the corresponding Fourier route, including the M10 transform.}

\section{Introduction}\label{p7-sec:paperintro}
For background on Costas arrays see~\cite{p7-golombtaylor}. Costas defects can be counted modulo powers of two, which organises them into a filtration
and suggests lower bounds that no Costas array could meet. This paper develops the filtration at
order 32, proves the one unconditional bound it yields, settles the associated mod-4 minimisation,
and explains precisely why these bounds do not reach the existence question.

\paragraph{Dependencies.} This paper does not depend on the others.

\section{Preliminaries}\label{p7-sec:notation}
A \emph{Costas array} of order $n$ is an $n\times n$ permutation array in which the displacement
vectors between pairs of dots are all distinct. Writing the array as a permutation $h$ of
$\{0,\dots,n-1\}$, with the dot of column $i$ in row $h_i$, the condition is that for each
horizontal separation $m$ the vertical differences $h_{i+m}-h_i$ are distinct. Costas arrays are
known at every order up to 31; orders 32 and 33 are the smallest for which existence is open. The
dihedral group $D_4$ of the square maps Costas arrays to Costas arrays, and arrays are counted
either individually or by $D_4$ orbit (equivalence class). Further notation and terms are collected
in the notation table and glossary at the front of this volume.


The mod-64 injectivity of the 2-adic tower, noted in the project's consolidated report, can be turned into a
filtration on the defect set rather than on the search. This section records what
that yields: one unconditional theorem, and a measurement showing why the
associated optimisation problem does not bear on the existence question.

\section{Exact counting and the filtration}

For a permutation $p$ of order 32 and integers $r\ge0$ and $\ell,h\ge1$ with $r+\ell+h\le31$, put
\[
R_{\ell,h,r}(p)=p_{r+\ell+h}-p_{r+\ell}-p_{r+h}+p_r .
\]
It vanishes exactly when the dot pairs in columns $(r,r+h)$ and $(r+\ell,r+\ell+h)$ have the same
displacement, so $p$ is Costas if and only if no $R_{\ell,h,r}(p)$ is zero. When $\ell=h$ it is the
diagonal second difference $p_{r+2\ell}-2p_{r+\ell}+p_r$.


Directed triples $(r,\ell,h)$ correspond to $0\le a<b<c\le31$ via $a=r$,
$b=r+\ell$, $c=r+\ell+h$, so there are $\binom{32}{3}=4960$ directed defect
conditions. The diagonal count is $\sum_{\ell=1}^{15}(32-2\ell)=240$, and after
the duality quotient $R_{\ell,h,r}=R_{h,\ell,r}$ there are
$240+\tfrac12(4960-240)=2600$ distinct defect expressions.

Since $-62\le R_{\ell,h,r}\le62$ at order 32, being a sum of two values in $[0,31]$ minus another two, $R=0$ in $\Z$ if and only if
$64\mid R$. Writing $Z_k(p;\Omega)=\#\{\omega\in\Omega : 2^k\mid R_\omega(p)\}$,
one has $Z_1\ge Z_2\ge\cdots\ge Z_6$ and
\[
p \text{ is Costas} \iff Z_6(p;\Omega_{32})=0 .
\]
With $E_k=\#\{\omega : 2^k\,\|\,R_\omega\}$ the escape recurrence is
$Z_{k+1}=Z_k-E_k$.

\section{An unconditional lower bound}

\begin{theorem}[Diagonal parity floor]
\label{p7-thm:parityfloor}
For the diagonal defects $D_{\ell,r}=p_{r+2\ell}-2p_{r+\ell}+p_r$ at order 32,
let $e_0=\#\{i \text{ even}: p_i \text{ even}\}$. Then
\[
Z_1^{\mathrm{diag}} \;=\; 2\binom{e_0}{2}+2\binom{16-e_0}{2}
\;=\; 2(e_0-8)^2+112 \;\ge\; 112,
\]
with equality exactly when $e_0=8$.
\end{theorem}
\begin{proof}
$2\mid D_{\ell,r}$ iff $p_r\equiv p_{r+2\ell}\pmod2$, and $r$ and $r+2\ell$ have
the same index parity, so the 240 diagonal pairs are exactly the same-parity
index pairs. Among the 16 even indices, $e_0$ carry even values; since there are
16 even values in total, the odd indices carry $16-e_0$ even values and $e_0$ odd
ones. Counting same-value-parity pairs within each index class gives the stated
expression.
\end{proof}

This is the only unconditional universal bound the investigation has produced.
Verified identically against 400 random permutations, with the floor 112
attained.

\section{Why the mod-4 optimisation does not bear on existence}

Let $q_i=p_i\bmod4$, let $Z_2^{\mathrm{diag}}(q)$ count diagonal triples with
$4\mid q_r+q_{r+2\ell}-2q_{r+\ell}$, and let $m_2$ be its minimum over balanced
residue words, of which there are $32!/(8!)^4 = 99{,}561{,}092{,}450{,}391{,}000$.
A stochastic search produced a balanced word with $Z_2^{\mathrm{diag}}=16$,
independently re-evaluated over all 240 triples, so $0\le m_2\le16$; the exact
value is open.

The statistic is invariant under the lift, since it depends only on
$q=p\bmod 4$. Calibrating it against arrays that actually exist:

\begin{center}
\begin{tabular}{@{}lrr@{}}
\toprule
population & $Z_2^{\mathrm{diag}}$ & fraction of triples \\
\midrule
order-32 minimisation witness & 16 of 240 & $6.7\%$ \\
order-28 Costas census (712 arrays) & 32--56, mean 48.6 of 182 & $26.7\%$ \\
order-28 random permutations (3000) & 28--71, mean 44.2 of 182 & $24.3\%$ \\
order-32 random permutations (2000) & min 39, mean 58.4 of 240 & $24.3\%$ \\
\bottomrule
\end{tabular}
\end{center}

\begin{observation}
\label{p7-obs:m2}
$Z_2^{\mathrm{diag}}$ does not separate Costas arrays from random permutations:
the complete order-28 census sits at mean 48.6 against a random mean of 44.2 on
the same 182 triples, slightly \emph{above} random and well inside the spread.
Real Costas arrays do not avoid level-2 survivors; they let them escape at higher
levels, as the full profiles show --- Welch-30 runs
$Z_1\ldots Z_6=[102,50,25,9,2,0]$ and the order-31 parent $C_{31}$ runs
$[111,57,20,7,1,0]$. Minimising $Z_2^{\mathrm{diag}}$ therefore moves away from
the region where Costas arrays live, and no order-28 Costas array arises from a
word with $Z_2^{\mathrm{diag}}$ anywhere near the minimum.
\end{observation}

Determining $m_2$ exactly is a well-posed finite problem. It does not bear on
whether an order-32 Costas array exists, because the quantity being minimised
carries no Costas signal.

\section{A clique reduction and the strict lower bound $m_2\ge1$}

The zero-survivor case can be excluded exactly, and the structure does it without a solver,
although the minimum itself remains unknown. Every same-parity index pair $(i,k)$ with $k-i=2\ell$ satisfies
$1\le\ell\le15$ automatically, so \emph{the 240 diagonal triples are exactly the
240 same-parity index pairs}, each with a determined midpoint $j=(i+k)/2$.
Survival requires $u_i=u_k$ and $v_i\oplus v_k=u_i\oplus u_j$, so for fixed $u$
the zero-survivor condition forces
\[
v_i\oplus v_k \;=\; 1\oplus u_i\oplus u_j
\]
on \emph{every} pair with $u_i=u_k$. The constraint graph is therefore four
\emph{complete} graphs, indexed by (index parity) $\times$ ($u$-value), with
midpoint-generated edge labels $\lambda_{ik}=1\oplus u_i\oplus u_{(i+k)/2}$.

\begin{remark}[A tempting argument that does not work]
It is tempting to conclude inconsistency from equation counting: a clique on $s$
vertices carries $\binom{s}{2}$ equations in $s$ unknowns. That is invalid. The
$\F_2$ incidence matrix of $K_s$ has rank $s-1$, and the remaining
$\binom{s-1}{2}$ equations are cycle conditions which may all hold. Taking
$\lambda_{ik}=w_i\oplus w_k$ for arbitrary $w$ makes every equation consistent at
any size: verified for $s=8,12,16$ (28, 66 and 120 equations respectively, all
consistent). The inconsistency below is a property of the \emph{midpoint}
labelling, not of the redundancy.
\end{remark}

\begin{proposition}[Computational]
\label{p7-prop:m2}
For every $u\in\F_2^{32}$ and every $v\in\F_2^{32}$, the mod-4 word $q=u+2v$ has
at least one diagonal defect divisible by 4. Hence $m_2\ge1$, and combined with
the verified score-16 witness, $1\le m_2\le16$.
\end{proposition}

An exhaustive parity union-find over $u$ (with $u_0=0$ by translation) visits 744 nodes, reaches index 17, and finds no consistent $v$. The search extends $u$ one index
at a time, keeps the induced parity constraints on $v$ in a union-find structure, and abandons a
prefix as soon as those constraints are contradictory; every word extending an abandoned prefix
fails for the same reason. That is how 744 visited nodes account for all $2^{17}$ words. Two remarks on scope. The
search imposes \emph{no} balance condition on $v$ and, as a separate run
confirms, imposing balance on $u$ changes nothing: with and without
$\sum_i u_i=16$ the search is byte-identical at 744 nodes and zero solutions. The
statement is therefore stronger than $m_2\ge1$ --- no $v$ exists at all, balanced
or not, for any $u$.

This is a \emph{computational} proposition, not a structural theorem. The
obstruction first becomes unavoidable on positions $0\ldots17$: for every
smaller prefix some $u$ is consistent, and at $m=18$ all $2^{17}=131{,}072$ words
fail. There is accordingly no short case analysis, and a structural proof --- that
the midpoint labelling always yields a frustrated triangle --- remains open and
would be worth more than the exact value of $m_2$.

The exact object is
\[
m_2 \;=\; \min_{\substack{u\in\F_2^{32}\\ \sum_i u_i=16}}
\operatorname{fr}_{\mathrm{bal}}(G_u),
\]
the balance-constrained frustration index of the four midpoint-labelled complete
signed graphs, where the minimisation over $v$ must respect
$\sum_{i:u_i=0}v_i=\sum_{i:u_i=1}v_i=8$. Frustration index on $\le16$ vertices is
$2^{16}$ colourings computed exactly, so the residual work is the outer search
over $u$.

\section{The survivor hypergraph and the escape system}
\label{p7-sec:escape}

Scalar counts discard which defects share which positions. Retaining it, choose
bit $b_k$ and note $R\mapsto R+2^kA_{k,\omega}$ with
$A_{k,\omega}=b_{r+\ell+h}-b_{r+\ell}-b_{r+h}+b_r$. A survivor escapes iff
$A_{k,\omega}\equiv R_\omega/2^k+1 \pmod 2$, so killing every survivor in one
step is an affine system over $\F_2$,
\[
H_k b_k = c_k,
\]
with rows indexed by $S_k$, columns by positions, $H_k[\omega,i]$ the parity of
the multiplicity of $i$ in the rectangle, and $c_k[\omega]=R_\omega/2^k+1$. Let
$\Lambda_k(p)$ count solutions that also respect the residue balance the lift
requires. Calibrated at order 28, 60 Costas arrays against 60 random
permutations over all 3{,}276 rectangle defects:

\begin{center}
\begin{tabular}{@{}rlrrrr@{}}
\toprule
$k$ & arm & mean $|S_k|$ & mean rank & inconsistent & $\Lambda_k>0$ \\
\midrule
1 & Costas & 1623.3 & 25.0 & 60/60 & 0 \\
1 & random & 1644.3 & 25.0 & 60/60 & 0 \\
2 & Costas & 756.6 & 25.0 & 60/60 & 0 \\
2 & random & 825.6 & 25.0 & 60/60 & 0 \\
3 & Costas & 351.3 & 25.0 & 60/60 & 0 \\
3 & random & 416.4 & 25.0 & 60/60 & 0 \\
\bottomrule
\end{tabular}
\end{center}

\begin{observation}
\label{p7-obs:escape}
The one-step all-escape system is inconsistent for \emph{every} permutation
tested in both arms, and $\Lambda_k\equiv0$. The statistic separates nothing
because it is vacuous: with over 1{,}600 survivors at level 1, a rank-25 system
in 28 unknowns is over-determined by two orders of magnitude, and genuine Costas
arrays do not satisfy it either. They do not kill their survivors in one step;
they need only $S_6=\varnothing$ and reach it progressively. The one visible
difference, that Costas arrays carry slightly fewer survivors at levels 2 and 3,
is the expected consequence of $S_6=\varnothing$ rather than an independent
signal.
\end{observation}

This is the stopping criterion the dyadic programme set for itself: if Costas
arrays and random permutations are indistinguishable under rank, consistency and
balanced-lift statistics, the branch closes. They are, and it does. \section{The corrected statistic does separate}
\label{p7-sec:muk}

Demanding $s=\mathbf0$ was the error; the exact transition is
$H_k b = s+\mathbf1+\kappa_k$ with a free survivor mask $s$, subject to
$\sum_{i:a_i=t}b_i = 2^{4-k}$. The right statistic is then not whether all
survivors can be killed at once, but how aggressively the actual lift reduces
them. Fixing $p \bmod 2^k$ and re-lifting with 250 random balanced bit vectors,
at order 28 over 20 Costas arrays and 20 random permutations:

\begin{center}
\begin{tabular}{@{}rlrrrrr@{}}
\toprule
$k$ & arm & $|S_k|$ & actual $|S_{k+1}|$ & sampled min & sampled mean & actual percentile \\
\midrule
1 & Costas & 1628.2 & 759.8 & 766.9 & 817.8 & \textbf{4.0} \\
1 & random & 1653.2 & 825.1 & 780.5 & 830.9 & 45.9 \\
2 & Costas & 764.5 & 350.1 & 347.2 & 386.0 & \textbf{4.7} \\
2 & random & 837.4 & 426.1 & 381.1 & 421.8 & 56.4 \\
\bottomrule
\end{tabular}
\end{center}

\begin{observation}
\label{p7-obs:muk}
Costas arrays lift near-optimally. Their actual bit sits at roughly the fourth
percentile of balanced re-lifts of the same base, and at $k=1$ their mean
next-level count (759.8) is below the mean \emph{minimum} over 250 random
balanced re-lifts (766.9). Random permutations sit at the median, 45.9 and 56.4.
This is a genuine topology-sensitive separation of the kind the all-escape test
failed to produce.
\end{observation}

Two cautions before this is treated as leverage. First, it may be tautological:
a Costas array must reach $S_6=\varnothing$, so its filtration is obliged to
decay faster than a generic lift's, and the percentile may be measuring only that
obligation. The test that distinguishes the cases is whether \emph{near}-Costas
permutations, of defect 1 or 2, also lift at the low percentile; if they do, the
statistic is predictive and could drive a search heuristic, and if only exact
Costas arrays do, it is a restatement of the terminal condition. Second, the
signal is a percentile, not a bound, so it yields a heuristic rather than a sound
prune: an array may lift badly at one level and recover.

What survives unconditionally is Theorem~\ref{p7-thm:parityfloor} and
Proposition~\ref{p7-prop:m2}, both exact and both about the filtration rather than
about existence.

%=====================================================================

\section{Dyadic residue caps do not characterise the Costas property}
\label{p7-sec:caps}

The filtration above counts defect conditions modulo powers of two. A natural extension works with the
displacement vectors themselves: count the pairs of dots whose displacement falls in each residue class
modulo $2^r$, and use the fact that a Costas array has every integer displacement at most once. This
section shows that at order 32 no argument of that kind can exclude Costas arrays.

For a permutation $h$ of order 32 and pairs $i<j$, write $(d,e)=(j-i,\,h_j-h_i)$, so that $1\le d\le31$
and $1\le|e|\le31$. For $r\in\{1,\dots,5\}$, let $c_r(a,b)$ be the number of pairs whose displacement is
congruent to $(a,b)$ modulo $2^r$, and let $C_r(a,b)$ be the number of admissible integer vectors in
that class. A Costas array satisfies $c_r(a,b)\le C_r(a,b)$ for every $r$ and every class.

\begin{proposition}[Dyadic residue caps do not characterise Costas uniqueness]
\label{p7-prop:caps}
There is an order-32 permutation $h^\ast$ with $c_r(a,b)\le C_r(a,b)$ for every $r\in\{1,\dots,5\}$ and
every class $(a,b)$, whose defect is 55.
\end{proposition}
\begin{proof}
The witness is
\[
\begin{aligned}
h^\ast=(&5,19,25,28,29,8,0,26,14,3,31,1,23,30,2,20,\\
        &16,13,18,9,15,10,22,27,12,21,7,4,6,17,11,24).
\end{aligned}
\]
Checking it is a finite computation from the definitions: $h^\ast$ is a permutation, it satisfies every
cap at every level (at $r=5$ the fullest class is exactly at capacity), and its repeated displacements,
55 in excess, are spread over 20 of the 31 lags. The witness was found by local search, but the proof
does not depend on the search, only on the check, which the released script
\texttt{verify\_dyadic\_caps.py} performs.
\end{proof}

\begin{corollary}
\label{p7-cor:caps}
Let $E$ be any statistic determined by the occupancies $c_r$. No contradiction of the form
$L\le E(h)\le U$ with $L>U$ can be obtained when $L$ holds for every order-32 permutation and $U$ is
derived from the caps alone.
\end{corollary}
\begin{proof}
$h^\ast$ is a permutation satisfying every cap, so $L\le E(h^\ast)\le U$.
\end{proof}

The corollary covers the dyadic collision energies of the relative trajectories
$r_{ij}(t)=d+te$, the cyclic autocorrelation, the power spectrum $|\hat A|^2$, which is the Fourier
transform of the occupancies, and multiresolution summaries computed in perfect-shuffle or bit-reversed
order. It does not cover arguments that record which pair occupies which class. But at modulus 32 that
labelled data determines the permutation up to a cyclic shift of its values, so such an argument is the
original problem in other coordinates.

\paragraph{M10 and its radix-2 analogues.} M10 is the address permutation used to organise a radix-7
fast Fourier transform on the spiral honeycomb image algebra~\cite{p7-sheridan2005}: a radix-7 perfect
shuffle of hexagonal-lattice addresses, which gives geometric meaning to Cooley--Tukey regrouping but
supplies no invariant of the image. Its square-grid analogues, the rotation of an index's binary digits
and bit reversal, were applied to the eight dihedral images of the defect-2 near-solution $N_{32}$ in
three ways (Table~\ref{p7-tab:shuffle}). None produced a Costas array. Rotating the whole 10-bit pixel
address keeps one dot per row and column only when it reduces to the transpose, a symmetry of the
square. Acting on the columns alone retains the most structure, but never brings the defect below 11.
Bit reversal gives defects of 67--74, close to the best of random permutations, whose median defect is
90.

\begin{table}[H]
\centering
\small
\begin{tabular}{@{}lrrrr@{}}
\toprule
action & images & permutations & defect & Costas \\
\midrule
rows and columns, rotation & 32 & 32 & 30--44 & 0 \\
columns only, rotation & 32 & 32 & 11--32 & 0 \\
whole 10-bit address, rotation & 72 & 8 & 2 & 0 \\
rows and columns, bit reversal & 8 & 8 & 74 & 0 \\
columns only, bit reversal & 8 & 8 & 67--68 & 0 \\
whole 10-bit address, bit reversal & 8 & 8 & 74 & 0 \\
\bottomrule
\end{tabular}
\caption{Radix-2 shuffles applied to the eight dihedral images of $N_{32}$. The eight permutations from
rotating the whole address are $N_{32}$'s own dihedral images.}
\label{p7-tab:shuffle}
\end{table}

\noindent Together with Proposition~\ref{p7-prop:caps}, this closes M10, and unlabelled dyadic Fourier
summaries generally, as candidate order-32 obstructions.

\section{Status of the claims}\label{p7-sec:status}
The parity floor is proved, and was checked identically against 400 random permutations, with
the floor attained. The mod-4 result rests on an exhaustive computation over signed cliques, which is
part of the project's self-test suite. Neither result bears on the existence of an order-32 Costas
array.

\paragraph{Relation to prior work.} An elementary theorem and a computation, new as far as we know; their novelty has not been established by a systematic literature search.

\begin{center}
\begin{tabular}{@{}ll@{}}
\toprule
result & standing \\
\midrule
at least 112 even diagonal second differences & proved \\
$m_2\ge1$ (no zero-survivor assignment) & computational, exhaustive \\
$m_2\le16$ & a witness \\
the exact value of $m_2$ & open \\
dyadic caps do not characterise Costas & certificate, checked from the definitions \\
\bottomrule
\end{tabular}
\end{center} The residue-cap witness was found by local search; the proposition rests only on the released check of it.
\begin{paperrefs}
\bibitem{p7-golombtaylor} S.~W.~Golomb and H.~Taylor, \emph{Constructions and
properties of Costas arrays}, Proc.\ IEEE \textbf{72} (1984).
\bibitem{p7-sheridan2005} P.~Sheridan, D.~M.~Alexander and K.~S.~Nunn-Clark, \emph{Fast Fourier transform in the spiral honeycomb image algebra}, Proceedings of the Third International Conference on Information Technology and Applications (ICITA'05), IEEE, 2005, \href{https://doi.org/10.1109/ICITA.2005.143}{\texttt{doi:10.1109/ICITA.2005.143}}.
\end{paperrefs}

\clearpage
\chapter{An Independent SAT Verification of the Known Non-existence of Order-32 Transpose-Symmetric Costas Arrays}\label{p8-paper}
\paperauthor
\paperabstract{A Costas array is a permutation matrix in which no two pairs of dots are separated by the
same vector. No Costas array of order 32 is known, and order 32 is the smallest order for
which existence is open. This paper studies the arrays that are symmetric about their main
diagonal, which correspond to permutations that are their own inverse. Their non-existence at
order 32 is already known: exhaustive enumerations of symmetric Costas arrays, published in
1993 through order 32 and extended in 2009 to order 36, find none of order 32. We verify it independently, by an entirely different method: a complete
search built on a new decomposition.
The search rests on a decomposition. The fixed points of such an array form a Golomb ruler,
so there are 0, 2, 4 or 6 of them, and up to reflection the problem splits into
44{,}947 disjoint regions, one per fixed-point set. Each region is encoded exactly
as a SAT instance. Easy regions were decided by the solver Kissat on one core each; hard
regions by the distributed solver Dagster on a 104-core node, using a native split
whose split coordinates are chosen so that the fixed points refute the subproblems that
would otherwise dominate the run. The single region with no fixed points was decided by a
further case split on one transposition, with reflection symmetry breaking. All 44{,}947
regions received UNSAT verdicts and none produced a validated solution; the listed production
runs used 16{,}887 service units in our HPC environment. The verification is computational, under explicit trust assumptions: the negative verdicts
are solver results, not yet supported by checked proof certificates. Its main interest is
methodological: it shows how the structure of a combinatorial problem can drive the
decomposition of its hardest parts.}


\section{Introduction}
This section sets out the question this paper addresses, why it is of interest, and how the paper is organised.

%======================================================================

Costas arrays are a small, concrete object with a stubborn open problem. They are easy to
define and to check, they exist at every order up to 31, and yet nobody knows whether one
exists at order 32. A natural sub-question asks whether one exists that is symmetric about its
main diagonal. The literature answers it: exhaustive enumerations published in 1993 and 2009 find
none of order 32 (Section~\ref{p8-sec:censuses}). This paper answers it again,
independently, by a complete computational search organised in an entirely different way.

The main interest is less the answer than how it was reached. Rather than searching the symmetric
arrays directly, the problem is decomposed, using structure specific to symmetric arrays, into 44{,}947
independent pieces, each an exact Boolean satisfiability (SAT) problem. A distributed
solver splits the hardest pieces further, in a way designed so that the fixed points of the
array do most of the work. For this particular problem the decomposition is not cheaper than a
direct search, which, run to completion, needed about five hundred times less computation (Section~\ref{p8-sec:direct}).
Its value lies in what it makes exact and auditable, and in its potential for problems that direct
search cannot reach. The report is self-contained: it assumes no prior knowledge of Costas arrays.

\paragraph{How to read this paper.}
\begin{itemize}
\item Part I introduces Costas arrays (Section~\ref{p8-sec:primer}) and their
symmetries (Section~\ref{p8-sec:symmetry}), and recalls the earlier structural results of this
project that motivate the approach (Section~\ref{p8-sec:earlier}).
\item Part II is the method: the partition into regions, the SAT encoding, the
distributed decomposition, and the \emph{case split} for the region without fixed points.
\item Part III gives the results, an analysis of the decomposition's performance,
and the costs.
\end{itemize}

%======================================================================
\paragraph{Dependencies.} This paper uses only the shared notation at the front of the volume, and
compares its results with the published enumerations of 1993 and 2009. It is independent of
Papers~1--7.

\section{Costas arrays from first principles}\label{p8-sec:primer}
This section introduces Costas arrays for readers new to them. It fixes the notation used throughout, works through the defining condition on examples, and summarises what the literature establishes about their existence.

%======================================================================

\subsection{Definition}

Take an $n\times n$ grid and place $n$ dots, one in each column and one in each row. Such a
placement is a \emph{permutation}: number columns and rows $0,\dots,n-1$, and let $\pi(i)$ be
the row of the dot in column $i$. Between any two dots, in columns $i<j$, draw the arrow from
the left one to the right one; its \emph{displacement} is the vector
\[
\bigl(j-i,\ \pi(j)-\pi(i)\bigr),
\]
horizontal distance first. There are $\binom n2$ such arrows.

\begin{definition}
A permutation $\pi$ of $\{0,\dots,n-1\}$ is a Costas array of order $n$ if its
$\binom n2$ displacement vectors are pairwise distinct.
\end{definition}

In words: no two arrows between dots are equal as vectors, so no pattern of two dots ever
repeats anywhere else in the grid, even shifted.

\begin{figure}[H]
\centering
\begin{minipage}{0.47\linewidth}\centering
\begin{tikzpicture}[x=0.55cm,y=0.55cm]
\draw[gray!45,very thin] (0,0) grid (6,6);
\fill[black] (0.5,2.5) circle (0.28);
\fill[black] (1.5,1.5) circle (0.28);
\fill[black] (2.5,5.5) circle (0.28);
\fill[black] (3.5,3.5) circle (0.28);
\fill[black] (4.5,4.5) circle (0.28);
\fill[black] (5.5,0.5) circle (0.28);
\draw[->,thick,blue] (0.5,2.5) -- (1.5,1.5);
\draw[->,thick,teal] (1.5,1.5) -- (4.5,4.5);
\node[font=\tiny,below] at (0.5,0) {0};
\node[font=\tiny,left] at (0,0.5) {0};
\node[font=\tiny,below] at (1.5,0) {1};
\node[font=\tiny,left] at (0,1.5) {1};
\node[font=\tiny,below] at (2.5,0) {2};
\node[font=\tiny,left] at (0,2.5) {2};
\node[font=\tiny,below] at (3.5,0) {3};
\node[font=\tiny,left] at (0,3.5) {3};
\node[font=\tiny,below] at (4.5,0) {4};
\node[font=\tiny,left] at (0,4.5) {4};
\node[font=\tiny,below] at (5.5,0) {5};
\node[font=\tiny,left] at (0,5.5) {5};
\end{tikzpicture}
\par\small (a) $\pi=(2,1,5,3,4,0)$, a Costas array
\end{minipage}\hfill
\begin{minipage}{0.47\linewidth}\centering
\begin{tikzpicture}[x=0.7cm,y=0.7cm]
\draw[gray!45,very thin] (0,0) grid (4,4);
\fill[black] (0.5,0.5) circle (0.28);
\fill[black] (1.5,2.5) circle (0.28);
\fill[black] (2.5,1.5) circle (0.28);
\fill[black] (3.5,3.5) circle (0.28);
\draw[->,thick,red] (0.5,0.5) -- (1.5,2.5);
\draw[->,thick,red] (2.5,1.5) -- (3.5,3.5);
\node[font=\tiny,below] at (0.5,0) {0};
\node[font=\tiny,left] at (0,0.5) {0};
\node[font=\tiny,below] at (1.5,0) {1};
\node[font=\tiny,left] at (0,1.5) {1};
\node[font=\tiny,below] at (2.5,0) {2};
\node[font=\tiny,left] at (0,2.5) {2};
\node[font=\tiny,below] at (3.5,0) {3};
\node[font=\tiny,left] at (0,3.5) {3};
\end{tikzpicture}
\par\small (b) $\pi=(0,2,1,3)$, not Costas
\end{minipage}
\caption{(a) An order-6 Costas array. The two arrows shown, $(1,-1)$ and $(3,3)$, are two of
its 15 displacements, and all 15 differ. (b) A permutation that fails: the arrows from column
0 to column 1 and from column 2 to column 3 are both $(1,2)$. Dots are at $(i,\pi(i))$ with
row 0 at the bottom.}
\label{p8-fig:primer}
\end{figure}

\subsection{Checking the condition lag by lag}

It is convenient to group the displacements by their horizontal component, the \emph{lag}
$s=j-i$. For each lag $s$ the condition says that the vertical differences
$\pi(x+s)-\pi(x)$, for $x=0,\dots,n-1-s$, are all different. Writing these rows out gives
the \emph{difference triangle} of $\pi$. For Figure~\ref{p8-fig:primer}(a):
\begin{center}
\begin{tabular}{@{}r l@{}}
lag 1: & $-1,\ \ 4,\ -2,\ \ 1,\ -4$\\
lag 2: & $\ \ 3,\ \ 2,\ -1,\ -3$\\
lag 3: & $\ \ 1,\ \ 3,\ -5$\\
lag 4: & $\ \ 2,\ -1$\\
lag 5: & $-2$
\end{tabular}
\end{center}
Every row has distinct entries, so the array is Costas. (Entries may repeat \emph{between}
rows, since those arrows have different horizontal components.)

\subsection{What is known}

Costas arrays were introduced in the 1960s for sonar and radar: a frequency-hopping signal
whose hop pattern is a Costas array has an ideal ``thumbtack'' ambiguity function, since no
shift in time and frequency makes more than one hop coincide~\cite{p8-costas1984}. Two algebraic
constructions produce infinite families~\cite{p8-golombtaylor1984}:
\begin{itemize}
\item \emph{Welch}: for a prime $p$ and a primitive root $g$ modulo $p$,
$\pi(i)=g^{\,i+1}\bmod p\;-\;1$ for $i=0,\dots,p-2$ is a Costas array of order $p-1$. With
$p=7$ and $g=3$ this gives exactly Figure~\ref{p8-fig:primer}(a).
\item \emph{Lempel--Golomb}: from primitive elements of the finite field of $q$ elements, an
array of order $q-2$; variants give $q-3$ and a few others.
\end{itemize}
Every order up to 31 is reached by a construction or by search, and all Costas arrays have
been enumerated exhaustively up to order 29~\cite{p8-drakakis2011}. Orders 32 and 33 are the
smallest with no known example~\cite{p8-drakakis2006}. Neither construction reaches them, and
exhaustive search is far beyond reach: there are $32!\approx2.6\times10^{35}$ permutations,
and the cost of exhaustive enumeration grows rapidly with the order.

%======================================================================
\section{Symmetry, and the transpose-symmetric subclass}\label{p8-sec:symmetry}
This section explains why transpose symmetry is the natural subclass to study, being the only symmetry of the square that a Costas array can have, and what its members look like. It separates what is known from the literature from what was computed in this project.

%======================================================================

\subsection{The symmetries of the square}

Rotating the grid by a multiple of $90^\circ$, or reflecting it in a horizontal, vertical or
diagonal line, maps permutation matrices to permutation matrices. These eight motions form
the dihedral group $D_4$. Each of them maps Costas arrays to Costas arrays, because it maps
the set of displacement vectors one-to-one onto another set of vectors. So Costas arrays come
in families of up to eight. An array is \emph{geometrically symmetric} if some non-identity motion maps it
to itself.

The following classification is elementary background, and we make no claim to its novelty.
Taylor, Rickard and Drakakis enumerate diagonal symmetry alongside symmetry types that are not
motions of the square~\cite{p8-taylor2009}; the classification below explains why, among the motions,
diagonal symmetry is the only one that needs enumerating. We include a short proof because it fixes the scope of the present search, and to keep
the paper self-contained.

\begin{proposition}[Only transpose symmetry is possible]\label{p8-prop:onlytranspose}
For $n\ge3$, the only possible non-identity elements of the stabiliser of a Costas array in $D_4$ are the two diagonal reflections.
\end{proposition}
\begin{proof}
A reflection in the vertical axis fixes $\pi$ only if $\pi(i)=\pi(n-1-i)$ for every $i$. For any
$i\ne n-1-i$ that puts two dots in one row, which is impossible for a permutation. A reflection in
the horizontal axis requires $\pi(i)=n-1-\pi(i)$ for every $i$, putting every dot in the same
central row, which is also impossible.

For the $180^\circ$ rotation, invariance gives $\pi(n-1-i)=n-1-\pi(i)$ for every $i$. Take
columns $i<j$, and their rotated pair $n-1-j<n-1-i$. The two pairs have the same displacement:
\[
\begin{aligned}
(n-1-i)-(n-1-j) &= j-i,\\
\pi(n-1-i)-\pi(n-1-j) &= \bigl(n-1-\pi(i)\bigr)-\bigl(n-1-\pi(j)\bigr)=\pi(j)-\pi(i).
\end{aligned}
\]
The two pairs of columns coincide only when $i+j=n-1$. For $n\ge3$, the pair $i=0$, $j=1$ has
$i+j=1\ne n-1$, so two distinct pairs of dots have the same displacement, contradicting the Costas
condition. Rotations by $90^\circ$ and $270^\circ$ square to the $180^\circ$ rotation, so they are
excluded too. That leaves the two diagonal reflections, which are conjugate under the group.
\end{proof}

Reflection in the main diagonal swaps the roles of rows and columns: it maps $\pi$ to
$\pi^{-1}$. So an array is transpose-symmetric exactly when $\pi=\pi^{-1}$, i.e.\ when
$\pi$ is an involution. By the proposition, studying involutions covers every geometrically symmetric
Costas array, since the other diagonal is conjugate to it.

\subsection{What a transpose-symmetric array looks like}

An involution consists of fixed points, coordinates $c$ with $\pi(c)=c$ whose dots lie
on the main diagonal, and transpositions $(a\,b)$ with $\pi(a)=b$, $\pi(b)=a$, whose
two dots are mirror images across the diagonal. Figure~\ref{p8-fig:sym} shows two real examples.

\begin{figure}[H]
\centering
\begin{minipage}{0.5\linewidth}\centering
\begin{tikzpicture}[x=0.32cm,y=0.32cm]
\draw[gray!45,very thin] (0,0) grid (16,16);
\draw[gray,dashed] (0,0) -- (16,16);
\fill[red!80!black] (0.5,0.5) circle (0.336);
\fill[black] (1.5,5.5) circle (0.28);
\fill[black] (2.5,15.5) circle (0.28);
\fill[black] (3.5,4.5) circle (0.28);
\fill[black] (4.5,3.5) circle (0.28);
\fill[black] (5.5,1.5) circle (0.28);
\fill[black] (6.5,9.5) circle (0.28);
\fill[black] (7.5,13.5) circle (0.28);
\fill[black] (8.5,10.5) circle (0.28);
\fill[black] (9.5,6.5) circle (0.28);
\fill[black] (10.5,8.5) circle (0.28);
\fill[red!80!black] (11.5,11.5) circle (0.336);
\fill[red!80!black] (12.5,12.5) circle (0.336);
\fill[black] (13.5,7.5) circle (0.28);
\fill[red!80!black] (14.5,14.5) circle (0.336);
\fill[black] (15.5,2.5) circle (0.28);
\node[font=\tiny,below] at (0.5,0) {0};
\node[font=\tiny,left] at (0,0.5) {0};
\node[font=\tiny,below] at (1.5,0) {1};
\node[font=\tiny,left] at (0,1.5) {1};
\node[font=\tiny,below] at (2.5,0) {2};
\node[font=\tiny,left] at (0,2.5) {2};
\node[font=\tiny,below] at (3.5,0) {3};
\node[font=\tiny,left] at (0,3.5) {3};
\node[font=\tiny,below] at (4.5,0) {4};
\node[font=\tiny,left] at (0,4.5) {4};
\node[font=\tiny,below] at (5.5,0) {5};
\node[font=\tiny,left] at (0,5.5) {5};
\node[font=\tiny,below] at (6.5,0) {6};
\node[font=\tiny,left] at (0,6.5) {6};
\node[font=\tiny,below] at (7.5,0) {7};
\node[font=\tiny,left] at (0,7.5) {7};
\node[font=\tiny,below] at (8.5,0) {8};
\node[font=\tiny,left] at (0,8.5) {8};
\node[font=\tiny,below] at (9.5,0) {9};
\node[font=\tiny,left] at (0,9.5) {9};
\node[font=\tiny,below] at (10.5,0) {10};
\node[font=\tiny,left] at (0,10.5) {10};
\node[font=\tiny,below] at (11.5,0) {11};
\node[font=\tiny,left] at (0,11.5) {11};
\node[font=\tiny,below] at (12.5,0) {12};
\node[font=\tiny,left] at (0,12.5) {12};
\node[font=\tiny,below] at (13.5,0) {13};
\node[font=\tiny,left] at (0,13.5) {13};
\node[font=\tiny,below] at (14.5,0) {14};
\node[font=\tiny,left] at (0,14.5) {14};
\node[font=\tiny,below] at (15.5,0) {15};
\node[font=\tiny,left] at (0,15.5) {15};
\end{tikzpicture}
\par\small (a) order 16, fixed points $\{0,11,12,14\}$
\end{minipage}\hfill
\begin{minipage}{0.46\linewidth}\centering
\begin{tikzpicture}[x=0.35cm,y=0.35cm]
\draw[gray!45,very thin] (0,0) grid (14,14);
\draw[gray,dashed] (0,0) -- (14,14);
\fill[black] (0.5,3.5) circle (0.28);
\fill[black] (1.5,5.5) circle (0.28);
\fill[black] (2.5,9.5) circle (0.28);
\fill[black] (3.5,0.5) circle (0.28);
\fill[black] (4.5,12.5) circle (0.28);
\fill[black] (5.5,1.5) circle (0.28);
\fill[black] (6.5,8.5) circle (0.28);
\fill[black] (7.5,13.5) circle (0.28);
\fill[black] (8.5,6.5) circle (0.28);
\fill[black] (9.5,2.5) circle (0.28);
\fill[black] (10.5,11.5) circle (0.28);
\fill[black] (11.5,10.5) circle (0.28);
\fill[black] (12.5,4.5) circle (0.28);
\fill[black] (13.5,7.5) circle (0.28);
\node[font=\tiny,below] at (0.5,0) {0};
\node[font=\tiny,left] at (0,0.5) {0};
\node[font=\tiny,below] at (1.5,0) {1};
\node[font=\tiny,left] at (0,1.5) {1};
\node[font=\tiny,below] at (2.5,0) {2};
\node[font=\tiny,left] at (0,2.5) {2};
\node[font=\tiny,below] at (3.5,0) {3};
\node[font=\tiny,left] at (0,3.5) {3};
\node[font=\tiny,below] at (4.5,0) {4};
\node[font=\tiny,left] at (0,4.5) {4};
\node[font=\tiny,below] at (5.5,0) {5};
\node[font=\tiny,left] at (0,5.5) {5};
\node[font=\tiny,below] at (6.5,0) {6};
\node[font=\tiny,left] at (0,6.5) {6};
\node[font=\tiny,below] at (7.5,0) {7};
\node[font=\tiny,left] at (0,7.5) {7};
\node[font=\tiny,below] at (8.5,0) {8};
\node[font=\tiny,left] at (0,8.5) {8};
\node[font=\tiny,below] at (9.5,0) {9};
\node[font=\tiny,left] at (0,9.5) {9};
\node[font=\tiny,below] at (10.5,0) {10};
\node[font=\tiny,left] at (0,10.5) {10};
\node[font=\tiny,below] at (11.5,0) {11};
\node[font=\tiny,left] at (0,11.5) {11};
\node[font=\tiny,below] at (12.5,0) {12};
\node[font=\tiny,left] at (0,12.5) {12};
\node[font=\tiny,below] at (13.5,0) {13};
\node[font=\tiny,left] at (0,13.5) {13};
\end{tikzpicture}
\par\small (b) order 14, no fixed points
\end{minipage}
\caption{Two transpose-symmetric Costas arrays. Fixed points are red, on the dashed diagonal;
every other dot has its mirror image across the diagonal. In (a) the fixed points have
pairwise distances 1, 2, 3, 11, 12 and 14, all different, which Lemma~\ref{p8-lem:golombprimer}
says must happen.}
\label{p8-fig:sym}
\end{figure}

\begin{lemma}[Fixed points form a Golomb ruler]\label{p8-lem:golombprimer}
In a transpose-symmetric Costas array, the pairwise differences between fixed points are all
distinct.
\end{lemma}
\begin{proof}
Fixed points $a<b$ have dots $(a,a)$ and $(b,b)$, with displacement $(b-a,b-a)$. Two pairs of
fixed points at the same distance would give the same displacement.
\end{proof}

A set of integers with all pairwise differences distinct is a Golomb ruler. Short
Golomb rulers with many marks are impossible: the shortest ruler with 8 marks has length 34, a known value in the literature on Golomb
rulers that Proposition~\ref{p8-prop:f} confirms independently by exhaustive search,
more than the 31 available on a $32\times32$ board. Since the number $f$ of fixed points has
the same parity as $n$, at order 32
\[
f\in\{0,2,4,6\}.
\]

\subsection{Constructions in the literature}
Some algebraic constructions in the literature produce symmetric arrays directly. The Lempel
construction, the special case of the Lempel--Golomb construction in which the two primitive
elements coincide, sets $\pi(i)=j$ exactly when $\alpha^i+\alpha^j=1$ in the field of $q$
elements. That relation is symmetric in $i$ and $j$, so every Lempel array is
transpose-symmetric. It has order $q-2$ for each prime power $q$~\cite{p8-golombtaylor1984}.
Drakakis, Gow and O'Carroll identify a second symmetric subfamily of the Lempel--Golomb
construction: when $q=r^2$ is a square, the arrays with $\beta=\alpha^r$~\cite{p8-dgo2006}. Both
families have order $q-2$. Order 32 would require $q=34$, which is not a prime power, so neither
family, and no other known construction, gives a symmetric array of order 32.

\subsection{Enumerations of symmetric arrays in the literature}\label{p8-sec:censuses}
Transpose-symmetric arrays are called \emph{diagonal} in the literature. Their numbers have been
determined by exhaustive search twice.
\begin{itemize}
\item Brown, Cenki, Games, Rushanan, Moreno and Pei tabulated the numbers of diagonal,
anti-reflective and consecutive Costas arrays of each order up to 32~\cite{p8-brown1993}.
\item Taylor, Rickard and Drakakis recomputed these counts~\cite{p8-taylor2009}. Their exhaustive
search is a backtracking search over permutations, using a midpoint test for the Costas condition
(the Ladendorf method) and symmetry reduction. For orders up to 27 they also tested every array in
the database of all known Costas arrays: the paper's text says the counts for those orders were
obtained that way, and its table caption says they were cross-checked that way. Above order 27,
including order 32, the counts come from the exhaustive search alone. They corrected a few entries
of the 1993 table, extended it to order 36 and to further symmetry types, and identified which
arrays are sporadic.
They found no sporadic symmetric array at any order from 24 to 36.
\item In a 2026 preprint, Dumitru and Budala report an exhaustive census of transpose-symmetric
(main-diagonal symmetric) arrays of orders 37 to 42, and no sporadic symmetric array at any order from
24 to 42~\cite{p8-dumitru2026symmetriccensus}.
\end{itemize}
Both tables record zero diagonal Costas arrays of order 32. In the 2009 table that entry is
marked as previously known, so it is a 1993 result confirmed in 2009. The non-existence of a
transpose-symmetric Costas array of order 32 has therefore been in the literature since 1993, and
it is independently confirmed by the search reported here, which shares neither method nor
software with either enumeration.

The 2009 table counts \emph{unique} arrays, one per equivalence family under the motions of the
square. Let $A$ be transpose-symmetric, that is, fixed by the main-diagonal reflection $D$. By
Proposition~\ref{p8-prop:onlytranspose} its stabiliser in $D_4$ is exactly $\{\mathrm{id},D\}$: it
cannot also be fixed by the other diagonal reflection, whose product with $D$ is the $180^\circ$
rotation. So its orbit has $8/2=4$ members. The $180^\circ$ rotation commutes with $D$, so it maps
$A$ to a second, different transpose-symmetric array; the reflections in the horizontal and vertical
axes conjugate $D$ to the anti-diagonal reflection, so they give the two anti-diagonally symmetric
members. Each class of diagonal arrays therefore contains exactly two transpose-symmetric arrays,
and the 2009 counts, one per class, should be exactly half of counts of individual arrays, which is what Table~\ref{p8-tab:symcounts}
compares.

\subsection{How rare they are}

In this project, transpose-symmetric Costas arrays were enumerated exhaustively up to order 24
by three independent implementations of an involution search, which agree on the size of the
search tree. The counts in Table~\ref{p8-tab:symcounts} come from those searches, not from the
literature. At every order in the table they are exactly twice the 2009 class counts, as the
counting conventions require, and both give zero at order 24. The 158 arrays of orders 12--22
are used throughout this project as test data.

\begin{table}[H]
\centering
\begin{tabular}{@{}rrrrrr@{}}
\toprule
order & $f=0$ & $f=2$ & $f=4$ & total & classes in~\cite{p8-taylor2009} \\
\midrule
12 & 2 & 28 & 4 & 34 & 17 \\
14 & 4 & 30 & 12 & 46 & 23 \\
16 & 4 & 28 & 8 & 40 & 20 \\
18 & 0 & 14 & 6 & 20 & 10 \\
20 & 0 & 6 & 2 & 8 & 4 \\
22 & 0 & 8 & 2 & 10 & 5 \\
24 & & & & 0 & 0 \\
\bottomrule
\end{tabular}
\caption{Transpose-symmetric Costas arrays by order and number of fixed points, from this
project's exhaustive searches, with the number of equivalence classes of diagonal arrays reported
in~\cite{p8-taylor2009}. Every count is exactly twice the class count. None has
six fixed points, and there are none at all at order 24. A direct search at order 32, run to completion,
visits $1.355\times10^{11}$ nodes (Section~\ref{p8-sec:direct}).}
\label{p8-tab:symcounts}
\end{table}

%======================================================================
\section{Earlier structural results in this project}\label{p8-sec:earlier}
This section summarises structural results obtained earlier in this project, before the campaign reported here. They motivate the change of decomposition adopted in Part~II, and they provide independent checks on its outcome.

%======================================================================

Before this campaign, the project attacked the order-32 symmetric problem through the
geometry of long displacements. This section records those results, which the present
search extends and with which it is consistent.

\paragraph{Large-large pairs.} Call a pair of coordinates $\{x,y\}$ large-large if
both $|x-y|\ge16$ and $|\pi(x)-\pi(y)|\ge16$: the two dots are far apart both horizontally and
vertically. Let $Z(\pi)$ be the number of large-large pairs. For an involution $\pi$ the map
$\tau\{x,y\}=\{\pi(x),\pi(y)\}$ permutes large-large pairs, and since $\pi^2=\mathrm{id}$ its
orbits have size 1 or 2. A pair is fixed by $\tau$ exactly when both points are fixed points,
or when $\{x,y\}$ is itself a transposition. This gives the \emph{orbit decomposition}
\[
Z = I_F + q + 2C,
\]
where $I_F$ counts large-large pairs of fixed points, $q$ the transpositions of length at
least 16, and $C$ the orbits of size two. It was verified on 12{,}000 random involutions and
on all 158 arrays of Table~\ref{p8-tab:symcounts}.

\paragraph{$Z=0$ is impossible.} Sort the displacements into four ``quadrants'' by whether
their horizontal and vertical components are short ($<16$) or long. A counting identity shows
that the short-short quadrant cannot be empty at order 32, and an automorphism argument shows
that the two mixed quadrants cannot both be empty. So the large-large quadrant is the only one
that could be empty, and that is the case $Z=0$. It forces all transposition lengths and all
fixed-point separations below 16, and hence $f\in\{2,4\}$. Both cases were closed by
exhaustive structural search: for $f=2$ by the chain $32{,}769{,}242\to9{,}180\to4\to0$
candidates under successively stronger sieves, and for $f=4$ by
$2{,}735\to20{,}360\to20{,}360\to0$.

\paragraph{$Z=1$ is impossible.} With $Z=1$ there is exactly one large-large pair: either a
single long transposition (branch T) or a single long pair of fixed points (branch F). That
yields seven labels by branch and $f$, all closed.

\begin{table}[H]
\centering
\begin{tabular}{@{}llll@{}}
\toprule
branch & $f$ & certificate & method \\
\midrule
T & 0 & $36\to1\to0$ & one residue, eliminated by the Costas condition \\
F & 2 & $36\to0$ & structural \\
T & 2 & $72\to0$ & full Costas search \\
T & 6 & --- & deductive: no 6-mark ruler fits a window of 15 \\
F & 6 & $728\to0$ & structural and full Costas \\
F & 4 & $19{,}405\to0$ & structural and full Costas \\
T & 4 & $2{,}735\to0$ & full Costas search \\
\bottomrule
\end{tabular}
\caption{The seven $Z=1$ labels and their closures. Each count is the number of root
configurations searched, every one ending in a fail-closed terminal status.}
\label{p8-tab:z1}
\end{table}

\paragraph{Why a different decomposition.} After these results, every remaining candidate
had $Z\ge2$. The next layer, $Z=2$, splits by the orbit decomposition into two structural
types and ten feasible cells, each needing its own bespoke search, and the layers grow
quickly from there. The present campaign takes a different route. It partitions by the
\emph{fixed-point set} rather than by $Z$, and hands each region to general SAT
technology, with a decomposition designed around the fixed points. Two consistency checks link
the two programmes:
\begin{itemize}
\item the $Z=2$ analysis derived, independently, a catalogue of reflection-canonical
fixed-point sets with totals 1, 256, 13{,}587 and 31{,}103 for $f=0,2,4,6$, exactly the region
counts of Section~\ref{p8-sec:framework};
\item every candidate excluded by the $Z=0$ and $Z=1$ theorems lies in some region, and all
regions are now closed, so the two programmes agree.
\end{itemize}
The earlier results are finite computations with fail-closed status, not formally certified,
like the present ones.

\section{The fixed-point partition}\label{p8-sec:framework}
This section turns the Golomb-ruler property into an exhaustive partition of the order-32 problem. It gives the formal definitions used in the rest of the paper and derives the 44{,}947 regions.

%======================================================================

Throughout, $n=32$ and coordinates are $[n]=\{0,\dots,31\}$, starting at zero.

\paragraph{Geometric convention.} The dot of column $i$ is at $(i,\pi(i))$: the first
coordinate is the column, the second the row. For two columns $i<j$ the displacement is
$(j-i,\ \pi(j)-\pi(i))$, whose first component, the lag $s=j-i$, is positive.
Distinctness of all displacements is therefore the same as requiring, for each lag
$s=1,\dots,31$ separately, that the vertical differences $\pi(x+s)-\pi(x)$ over
$x=0,\dots,31-s$ be pairwise distinct. The SAT encoding uses exactly this lag-by-lag form.

\begin{definition}
A permutation $\pi$ of $[n]$ is \emph{Costas} if the vectors $(j-i,\pi(j)-\pi(i))$,
$0\le i<j\le 31$, are pairwise distinct; it is a \emph{Costas involution} if also
$\pi^2=\mathrm{id}$. Its fixed-point set is $\Fix(\pi)=\{i:\pi(i)=i\}$, with $f=|\Fix(\pi)|$.
\end{definition}

\begin{example}[A collision]
If $\pi(i+s)-\pi(i)=\delta=\pi(j+s)-\pi(j)$ for some $i\ne j$, the column pairs $(i,i+s)$
and $(j,j+s)$ both have displacement $(s,\delta)$, and $\pi$ is not Costas. In the encoding
of Section~\ref{p8-sec:encoding}, the indicators $p_{s,i,\delta}$ and $p_{s,j,\delta}$ are then
both forced true, violating the at-most-one constraint for $(s,\delta)$.
\end{example}

\begin{lemma}\label{p8-lem:sym}
The permutation matrix of $\pi$ is transpose-symmetric if and only if $\pi$ is an involution.
\end{lemma}
\begin{proof}
The matrix has dots $(i,\pi(i))$; its transpose has dots $(\pi(i),i)$, which form the matrix
of $\pi^{-1}$. They coincide iff $\pi=\pi^{-1}$.
\end{proof}

An involution splits $[n]$ into fixed points and disjoint transpositions $(a\,b)$ with
$\pi(a)=b$, $\pi(b)=a$. Since $n$ is even, $f$ is even.

\begin{lemma}[Fixed points form a Golomb ruler]\label{p8-lem:golomb}
If $\pi$ is a Costas involution, the pairwise differences of $\Fix(\pi)$ are distinct.
\end{lemma}
\begin{proof}
For fixed points $a<b$ the dots $(a,a)$ and $(b,b)$ have displacement $(b-a,b-a)$. Two
different pairs with the same difference would repeat a displacement.
\end{proof}

\begin{proposition}\label{p8-prop:f}
At order 32, $f\in\{0,2,4,6\}$.
\end{proposition}
\begin{proof}
$f$ is even. The shortest Golomb ruler with eight marks has length 34~\cite{p8-oeisgolomb}, while
marks in $[32]$ span at most 31, so no eight of them have distinct pairwise differences, and any
ruler with more marks contains an 8-mark one. An exhaustive backtracking search over 8-element
subsets of $[32]$, run to completion, confirms this independently. Hence $f\le 6$.
\end{proof}

\begin{lemma}[Reflection]\label{p8-lem:refl}
Let $\rho(i)=31-i$. If $\pi$ is a Costas involution, so is $\rho\pi\rho$, and
$\Fix(\rho\pi\rho)=\rho(\Fix\pi)$.
\end{lemma}
\begin{proof}
$(\rho\pi\rho)^2=\rho\pi^2\rho=\mathrm{id}$. The array of $\rho\pi\rho$ is that of $\pi$
rotated by $180^\circ$, which negates every displacement and so preserves distinctness.
Finally $\rho\pi\rho(i)=i\iff\pi(\rho(i))=\rho(i)$.
\end{proof}

\begin{definition}[Region]
For a Golomb ruler $F\subseteq[n]$ with $|F|\in\{0,2,4,6\}$, the region $R_F$ is the set of
Costas involutions $\pi$ with $\Fix(\pi)=F$ exactly. $F$ is reflection-canonical if,
as a sorted tuple, $F\le\rho(F)$ lexicographically. Regions are named \code{R00000},
\code{R00001}, \dots{} in the order of that enumeration; for a named region,
$F(\code{R}\ldots)$ denotes its fixed-point set.
\end{definition}

\begin{theorem}[Partition]\label{p8-thm:partition}
Every Costas involution of order 32 lies in exactly one region and, after applying $\rho$ if
necessary, in a reflection-canonical one. Hence a transpose-symmetric Costas array of order
32 exists iff some reflection-canonical region is non-empty.
\end{theorem}
\begin{proof}
$\Fix(\pi)$ is unique, a Golomb ruler (Lemma~\ref{p8-lem:golomb}), and of admissible size
(Proposition~\ref{p8-prop:f}); Lemma~\ref{p8-lem:refl} maps a non-canonical case to the canonical
region of $\rho(F)$.
\end{proof}

\begin{figure}[H]
\centering
\begin{minipage}{0.9\linewidth}
\small
Algorithm 1 (region enumeration, \code{regions.py})\label{p8-alg:enum}
\begin{enumerate}[nosep, label=\arabic*.]
\item emit the empty set as \code{R00000};
\item for $f=2,4,6$, and for each $f$-subset $F$ of $[32]$ in lexicographic order:
\item \quad reject $F$ unless its $\binom{f}{2}$ positive differences are distinct;
\item \quad reject $F$ unless $F\le\rho(F)$ as sorted tuples;
\item \quad emit $F$ with the next region number, its size, diameter $\max F-\min F$, and
its number of even elements.
\end{enumerate}
For the 8-mark check of Proposition~\ref{p8-prop:f} the same distinct-difference test is applied
by backtracking from mark 0 with span at most 31; it returns no ruler.
\end{minipage}
\end{figure}

\begin{table}[H]
\centering
\begin{tabular}{@{}rrrr@{}}
\toprule
$f$ & Golomb rulers in $[32]$ & fixed by $\rho$ & canonical regions \\
\midrule
0 & 1 & 1 & 1 \\
2 & 496 & 16 & 256 \\
4 & 27{,}174 & 0 & 13{,}587 \\
6 & 62{,}206 & 0 & 31{,}103 \\
\midrule
total & 89{,}877 & 17 & 44{,}947 \\
\bottomrule
\end{tabular}
\caption{Region counts. Canonical $=(\text{raw}+\text{fixed})/2$. A $\rho$-fixed set with
four or more marks contains $a<b$ and $31-b<31-a$, which repeat the difference $b-a$, so
none is a Golomb ruler. Script and output digests are in the released software archive.}
\label{p8-tab:regions}
\end{table}

\begin{definition}[Seed refutation]\label{p8-def:seed}
For free $x<y$, form the multiset of displacement vectors between every pair of dots in
$\{(c,c):c\in F\}\cup\{(x,y),(y,x)\}$, each pair taken in column order. The pair $\{x,y\}$ is
\emph{seed-refuted} for $F$ if this multiset contains a repeated vector.
\end{definition}

\begin{lemma}\label{p8-lem:seed}
If $\{x,y\}$ is seed-refuted, no member of $R_F$ contains the transposition $(x\,y)$, and the
encoding of Section~\ref{p8-sec:encoding} refutes $x_{xy}$ by unit propagation alone.
\end{lemma}
\begin{proof}
Any such $\pi$ contains all of those dots. In the encoding, $x_{xy}$ forces $h_{xy},h_{yx}$
(clause family 1), which together with the constant fixed-point values force two
displacement indicators with the same $(s,\delta)$ true (family 5), and their binary
at-most-one clause (family 6) is then falsified.
\end{proof}

%======================================================================
\section{SAT encoding of a region}\label{p8-sec:encoding}
This section gives the propositional encoding of a single region, proves it exact, and records how it was validated. Every solver verdict in the paper rests on this encoding.

%======================================================================

\paragraph{Notation.} For literals $q_1,\dots,q_m$:
$\ALO(q_1,\dots,q_m)=(q_1\vee\dots\vee q_m)$,
$\AMO(q_1,\dots,q_m)=\bigwedge_{1\le i<j\le m}(\neg q_i\vee\neg q_j)$ (pairwise), and
$\EO=\ALO\wedge\AMO$.

Fix a region $F$ and let $V=[n]\setminus F$.

\paragraph{Variables.}
\begin{itemize}
\item $x_{ab}$ for $a,b\in V$, $a<b$: transposition $(a\,b)$ present.
\item $h_{ij}$ for $i,j\in V$, $i\ne j$: $\pi(i)=j$.
\item $p_{s,x,\delta}$ for $1\le s\le 31$, $0\le x\le 31-s$, $-31\le\delta\le 31$, created only
when \emph{realisable}: there are permitted values $u$ for column $x$ and $w$ for column $x+s$
with $w-u=\delta$. A fixed column $c\in F$ permits only $u=c$; a free column permits every
$u\in V$ other than itself.
\end{itemize}
The literal $[\pi(i)=j]$ means $h_{ij}$ for $i,j\in V$; it is the constant \emph{true} if
$i=j\in F$, and \emph{false} if exactly one of $i,j$ is in $F$.

\paragraph{Clauses.}
\begin{enumerate}
\item \emph{Edge--map link:} $(\neg x_{ab}\vee h_{ab})$, $(\neg x_{ab}\vee h_{ba})$,
$(\neg h_{ab}\vee x_{ab})$ for $a<b$.
\item \emph{One image per free coordinate:} $\EO(\{h_{ij}\}_{j\in V\setminus\{i\}})$ for
each $i\in V$.
\item \emph{One transposition per free coordinate:} $\EO(\{x_e: c\in e\})$ for each $c\in V$.
\item \emph{Distinct transposition lengths:} $\AMO(\{x_{a,a+d}\})$ for each length $d$.
\item \emph{Displacement indicators:} for each column pair $(x,x+s)$ and permitted values
$(u,w)$, the clause $\neg[\pi(x)=u]\vee\neg[\pi(x+s)=w]\vee p_{s,x,w-u}$, with constant-true
literals removed. There are three cases: both columns free (a three-literal clause); one
column fixed (a two-literal clause); both fixed (the unit clause $p_{s,x,w-u}$, so
fixed-to-fixed displacements are asserted outright and can collide with displacements of
transpositions).
\item \emph{Costas:} $\AMO(\{p_{s,x,\delta}\}_x)$ for each $(s,\delta)$.
\end{enumerate}

\begin{remark}
Given family (1), families (2) and (3) imply each other: an image $h_{ij}$ forces the edge
$x_{ij}$, and an edge forces both images. Both are kept because they propagate differently,
through the map and through the matching. Family (4) is implied by the Costas property,
since a transposition of length $d$ has internal displacement $(d,-d)$, and is kept for
propagation. The proof of Theorem~\ref{p8-thm:exact} uses (1), (3), (5) and (6).
\end{remark}

\begin{theorem}[Exactness]\label{p8-thm:exact}
The set of edge-variable projections of the models of the encoding is in bijection with the
set of Costas involutions $\pi$ with $\Fix(\pi)=F$.
\end{theorem}
\begin{proof}
Let $\alpha$ be a model. By (3) its true edges form a perfect matching of $V$; with $F$ fixed,
this defines an involution $\pi$ with $\Fix(\pi)=F$ exactly, because no variable expresses
$\pi(i)=i$ for $i\in V$ and (3) places every $i\in V$ in a transposition. By (1) and (2), the
$h$-variables agree with $\pi$. By (5), every displacement that $\pi$ realises has its
indicator true, so (6) forbids two column pairs at one lag with the same vertical difference:
$\pi$ is Costas. Conversely, a Costas involution with fixed set $F$, with $p$ set to true
exactly on realised displacements, satisfies every clause. The map from projections to
involutions is injective, since the projection is the matching itself. Indicators for
displacements that are not realised are not forced false, so several models can share one
projection; the bijection is therefore stated for projections.
\end{proof}

\begin{remark}[Enumeration blocks edge projections]\label{p8-rem:projblock}
The theorem is a bijection between involutions and edge-variable projections, not between
involutions and complete assignments. Any validation by model enumeration must therefore block the
edge projection, not the complete Boolean model: at order 10, a single edge projection has at least
64 complete models that differ only in unrealised indicators.
\end{remark}

\begin{table}[H]
\centering
\begin{tabular}{@{}rrrrrrr@{}}
\toprule
$f$ & $|V|$ & edge vars & map vars & indicator vars & total vars & clauses \\
\midrule
0 & 32 & 496 & 992 & 31{,}185 & 32{,}673 & 824{,}506 \\
2 & 30 & 435 & 870 & 29{,}030--29{,}087 & 30{,}335--30{,}392 & 666{,}858--669{,}798 \\
4 & 28 & 378 & 756 & 25{,}277--26{,}789 & 26{,}411--27{,}923 & 527{,}716--533{,}367 \\
6 & 26 & 325 & 650 & 22{,}943--24{,}290 & 23{,}918--25{,}265 & 412{,}138--418{,}785 \\
\bottomrule
\end{tabular}
\caption{Formula sizes. Edge and map counts are $\binom{|V|}{2}$ and $|V|(|V|-1)$ for every
region; indicator and clause counts depend on $F$ and are given as ranges over five sampled
regions per class (the $f=0$ class has one region).}
\label{p8-tab:sizes}
\end{table}

\paragraph{Validation.}
\begin{itemize}
\item \emph{Exhaustive equivalence at orders 8 and 10.} For each of the 39 and 106 fixed-point
sets, the comparison is made at the level of edge projections, not complete assignments
(Remark~\ref{p8-rem:projblock}). After each model $\alpha$ the enumerator adds the
projection-blocking clause
\[
\bigvee_{\alpha(x_{ab})=1}\neg x_{ab}\;\vee\bigvee_{\alpha(x_{ab})=0}x_{ab},
\]
over all edge variables of the region, so each matching is returned once, however its indicators
are completed. Each decoded involution is checked independently for being a permutation and an
involution, for its exact fixed-point set and for distinct displacements, and any duplicate is
rejected outside the solver. Conversely, every Costas involution found by brute-force enumeration of
the region's perfect matchings is checked to satisfy every clause under its canonical assignment, in
which each indicator is true exactly when its displacement is realised. The first direction places
the projected set inside the direct set, and the second places the direct set inside the projected
set; comparing the sorted sets confirms that they are equal in every region
(Table~\ref{p8-tab:projval}).
\item \emph{Regression.} All 158 known transpose-symmetric Costas arrays of orders 12--22
satisfy their region's CNF under direct clause evaluation.
\item \emph{Determinism.} All 225 Stage-1 CNFs (the first stage of the campaign) were regenerated
on the cluster and matched their recorded SHA-256 digests byte for byte.
\end{itemize}

\begin{table}[H]
\centering
\small
\begin{tabular}{@{}rrrrrcl@{}}
\toprule
order & regions & solver calls & projected & direct & failures & set digest \\
\midrule
8 & 39 & 57 & 18 & 18 & 0 / 0 / 0 / 0 & \texttt{2247552e2506108c} \\
10 & 106 & 134 & 28 & 28 & 0 / 0 / 0 / 0 & \texttt{c5067410de45442a} \\
\bottomrule
\end{tabular}
\caption{Projection-level validation with Kissat 4.0.4. The failures column counts invalid,
duplicate, missing and extra involutions, in that order. Digests are SHA-256 prefixes of the sorted
sets of decoded involutions, and the projected and direct digests agree. A regression test also
confirms that one order-10 projection has at least 64 complete models differing only in unrealised
indicators, none of which survives the projection blocker.}
\label{p8-tab:projval}
\end{table}


%======================================================================
\section{Solvers, verdicts and units}\label{p8-sec:solvers}
This section fixes how solver outcomes become verdicts, so that every reported result has a precise operational meaning, and describes the computing environment and how resources are counted.

%======================================================================

\paragraph{Kissat} 4.0.4~\cite{p8-kissat} runs on one core with a time limit. Its outcome maps
to verdicts as follows, all fail-closed:
\begin{itemize}
\item \st{UNSAT-UNCERTIFIED}: exit 20 together with \code{s UNSATISFIABLE};
\item \st{SAT-VALIDATED}: exit 10, \code{s SATISFIABLE}, and a model that passes independent
validation;
\item \st{TIMEOUT}: exit 0 with \code{s UNKNOWN}, i.e.\ stopped at the time limit;
\item \st{ERROR}: anything else.
\end{itemize}

\paragraph{Dagster}~\cite{p8-dagster} (AFSAT-Integration branch) prints no status line.
\begin{itemize}
\item \st{UNSAT-UNCERTIFIED} requires exit 0 and an existing, empty solution file, and in
addition the message-accounting audit of Section~\ref{p8-sec:gates}.
\item A non-empty solution file is decoded, and every model is validated.
\item Any other outcome is \st{TIMEOUT} or \st{ERROR}.
\end{itemize}
The ledgers record internal codes, mapped as: \code{UNSAT} (Kissat) and \code{UNSAT\_UNPROVEN}
(Dagster) to \st{UNSAT-UNCERTIFIED}; \code{SAT} with \code{model\_validated} to
\st{SAT-VALIDATED}; and \code{TIMEOUT} to \st{TIMEOUT}.

\paragraph{Model validation.} A reported model is decoded into $\pi$ and checked directly:
that it is a permutation and an involution, that $\Fix(\pi)=F$, and that its displacements
are pairwise distinct. A failure stops the run as a defect.

\paragraph{The HPC environment.} All runs used a PBS-scheduled cluster of dual-socket Intel Xeon
(Sapphire Rapids) nodes, each with 104 cores and 512\,GB of memory, with node-local scratch
storage and a shared parallel file system. Its accounting unit, the \emph{service unit} (SU), was
charged at about 2 SU per core-hour on all requested cores for the wall time used. All core-time
comparisons below count every allocated core, including Dagster's master.

%======================================================================
\section{The Dagster native split decomposition}\label{p8-sec:decomp}
This section specifies the decomposition used for every region too hard for a single core. It describes Dagster's execution model, defines the native split and proves it exhaustive, explains the coordinate-selection rule on which its performance depends, and defines the audit that accepts its results.

%======================================================================

\subsection{Dagster's execution model}

Dagster takes one CNF and a DAG. Each vertex names a subset of the CNF's clauses, by clause
index, and each arc is labelled with \emph{interface variables}. A master process schedules
work; each worker runs a CDCL solver. With enumeration mode \code{-e 0}:
\begin{itemize}
\item the source vertex enumerates every satisfying assignment of its clauses, projected onto
the interface, and each distinct projection is sent as a \emph{message};
\item the target vertex is solved once per message, with the message's literals imposed;
\item any satisfying assignment of the target ends the whole run (a solution exists);
\item otherwise the run ends normally only after the source is exhausted and every message
sent has been completed.
\end{itemize}

\subsection{The native split DAG}

Choose $k$ distinct free \emph{split coordinates} $c_1,\dots,c_k$, and let
$E_i=\{x_e: c_i\in e\}$ be the \emph{edge variables} incident to $c_i$.

\begin{figure}[H]
\centering
\begin{tikzpicture}[every node/.style={font=\small}, arr/.style={-{Stealth[length=2.4mm]}, thick},
  box/.style={draw, rounded corners, align=left, text width=11.5cm, inner sep=6pt}]
\node[box] (v0) {Vertex 0 --- split generator\\
clauses: $\EO(E_1),\dots,\EO(E_k)$ (copies of family (3) for $c_1,\dots,c_k$)\\
output: every assignment to $E_1\cup\dots\cup E_k$ satisfying them --- $M(D,k)$ messages};
\node[box, below=14mm of v0] (v1) {Vertex 1 --- exact regional solver\\
clauses: the complete regional CNF (families 1--6)\\
per message: solved with the message's literals imposed $\Rightarrow$ a solution
(SAT, ends the run) or UNSAT for that message};
\draw[arr] (v0) -- node[right, align=left]{\scriptsize one message per partner assignment\\[-1pt]
\scriptsize interface: $E_1\cup\dots\cup E_k$} (v1);
\end{tikzpicture}
\caption{The native split DAG. The region is UNSAT-UNCERTIFIED exactly when vertex 0 is
exhausted and every message is UNSAT at vertex 1 with no solution found.}
\label{p8-fig:dag}
\end{figure}

\paragraph{Why the clauses are copied into vertex 0.} Dagster stores one global CNF, and
vertices refer to clauses by index. Vertex 1 must contain the whole regional formula,
including the $\EO$ constraints for $c_1,\dots,c_k$. Appending a second copy of those $\EO$
clauses and giving vertex 0 only the copies makes the source a small, local enumerator,
while the terminal formula stays complete. Duplicate clauses do not change the model set.

\begin{theorem}[Exhaustive disjoint cover]\label{p8-thm:split}
The messages are pairwise distinct assignments to the interface. The sets of regional models
compatible with each message are pairwise disjoint and together cover every model of the
region; some of these sets are empty. Hence the region is UNSAT iff vertex 1 is UNSAT under
every message.
\end{theorem}
\begin{proof}
Let $\pi$ be a model. Each $c_i$ is in exactly one transposition, so $\pi$'s restriction to
$E_1\cup\dots\cup E_k$ satisfies every $\EO(E_i)$ and is enumerated as some message $m$;
$\pi$ is compatible with $m$ and with no other message, since distinct messages differ on an
interface variable. Conversely, a model of vertex 1 under $m$ is a model of the region,
because vertex 1 holds every clause. A message whose literals contradict the matching, for
example two split coordinates choosing the same outside partner, has an empty compatible set
and is refuted at vertex 1 almost at once.
\end{proof}

The theorem holds for any choice of split coordinates and any $k$; the choice affects only
the performance.

\subsection{Number of messages}

Each split coordinate has $D=|V|-1$ potential partners ($D=27$ for $f=4$, $D=29$ for $f=2$).
The constraints $\EO(E_i)$ and $\EO(E_j)$ share exactly one variable, $x_{c_ic_j}$: if $c_i$
chooses $c_j$, both constraints are satisfied by that one variable, so $c_j$ must choose
$c_i$. Vertex 0 contains only these local constraints, not the full matching, so two
unpaired split coordinates may choose the same outside partner. Such messages are generated
deliberately and refuted immediately downstream.

\begin{proposition}\label{p8-prop:M}
The number of messages is
\[
M(D,k)=\sum_{m=0}^{\lfloor k/2\rfloor} N_m(k)\,\bigl(D-(k-1)\bigr)^{k-2m},\qquad
N_m(k)=\frac{k!}{(k-2m)!\,2^m\,m!},
\]
where $N_m(k)$ counts the ways to choose $m$ disjoint pairs among the $k$ split coordinates.
\end{proposition}
\begin{proof}
A message pairs some $m$ disjoint pairs of split coordinates with each other; each of the
remaining $k-2m$ chooses independently among its $D-(k-1)$ partners outside the split set.
\end{proof}

Thus $M(D,2)=(D-1)^2+1$ and $M(D,3)=(D-2)^3+3(D-2)$, giving $M(27,2)=677$,
$M(29,2)=785$ and $M(29,3)=19{,}764$. These are exactly the counts observed in every Dagster
run, and they are re-verified by the audit of Section~\ref{p8-sec:gates}.

\subsection{Unsplit branches and coordinate selection}\label{p8-sec:rule}

A message in which $c_i$ and $c_j$ are paired with each other fixes one transposition among
them instead of two, so it is a subproblem one level shallower than the rest, typically
about an order of magnitude harder. Because a message is never divided across workers, such
an unsplit branch becomes the critical path of the run. If $\{c_i,c_j\}$ is
seed-refuted, the branch is refuted at once (Lemma~\ref{p8-lem:seed}). The selection rule
exploits this.

\begin{figure}[H]
\centering
\begin{minipage}{0.92\linewidth}\small
Algorithm 2 (coordinate selection, \code{dag\_native\_split.py})\label{p8-alg:sel}
\begin{enumerate}[nosep, label=\arabic*.]
\item order $\leftarrow$ \emph{free coordinates} sorted by (partner-domain size, index); all domains
are equal in this encoding, so this is ascending index;
\item $c_1\leftarrow$ first in order; record rule \code{smallest\_domain};
\item for $t=2,\dots,k$: $c_t\leftarrow$ the first coordinate in order, not yet chosen, that is
seed-refuted (Definition~\ref{p8-def:seed}) against every chosen coordinate; record
\code{seed\_refutes\_pairing};
\item if none exists: $c_t\leftarrow\arg\max$ of the number of chosen coordinates against which
it is seed-refuted, ties broken by the earliest position in order; record
\code{fallback\_refutes\_}$r$\code{\_of\_}$t{-}1$, where $r$ is that number;
\item store $(c_1,\dots,c_k)$ and the rules in \code{split\_meta.json}, and copy them into the
result ledger.
\end{enumerate}
\end{minipage}
\end{figure}

\paragraph{Effect on the sample.}
\begin{itemize}
\item $k=2$: a fully seed-refuted $c_2$ exists for all 224 sampled regions with
$f\in\{2,4,6\}$; only the $f=0$ region, which has no fixed points, falls back.
\item $k=3$, $f=2$: two fixed points seed-refute only about 20 of the 435 pairs. No sampled
$f=2$ region has a $c_3$ refuted against both $c_1$ and $c_2$, and only 6 of the 24 contain
three mutually seed-refuted coordinates at all. The fallback picks a $c_3$ refuted against
one of them, leaving $D-2=27$ unsplit branches of depth-2 size among 19{,}764 messages.
\end{itemize}

\subsection{Choice of depth}

The campaign uses depth $k=2$ for $f=4$ (677 messages, about 6.6 per worker) and $k=3$ for
$f=2$ (19{,}764 messages, about 190 per worker). A deeper split makes messages smaller and
shortens the critical path. It can also lower the total work, plausibly because Dagster's TiniSAT workers trail Kissat by
little on easy instances and by much on hard ones; this is a hypothesis, supported by a single
comparison. On $F(\code{R00043})=\{1,13\}$,
going from depth 2 to depth 3 reduced both the wall time (3{,}086\,s to 872\,s) and the
total message work (about 75 to 18 core-hours).

\subsection{Configuration}

\begin{table}[H]
\centering
\begin{tabular}{@{}lll@{}}
\toprule
setting & value & meaning \\
\midrule
\code{-m} & 0 & each worker runs TiniSAT (CDCL) \\
\code{-e} & 0 & enumerate all source solutions; stop at the first full solution \\
\code{-b} & 0 & depth-first allocation of messages \\
\code{-g} & 0 & table-based master \\
MPI ranks & 104 & 1 master + 103 workers, on one 104-core node \\
cap & 3{,}600\,s & per region, campaign default \\
\bottomrule
\end{tabular}
\caption{Dagster configuration. The worker count is $\min(M(D,k),103)$. Mode \code{-e 3}
stops each worker after the first solution per message, which is unsound for exhausting a
multi-vertex DAG; the driver rejects it.}
\label{p8-tab:config}
\end{table}

\subsection{Scheduler behaviour that shapes performance}

These properties, established from Dagster's source, explain the metrics below.
\begin{itemize}
\item \emph{Per-message cost.} A TiniSAT worker builds a fresh copy of vertex 1's clauses and a
new solver for each message, so nothing learned carries over. The campaign is therefore closer to cube-and-conquer than to one
globally learning solver: a nogood useful in many messages is rediscovered in each. This costs about 0.1\,s per
message at order 32.
\item \emph{No intra-message parallelism.} One worker solves a message. When fewer messages
remain than idle workers, idle workers receive \emph{copies} of unfinished messages, a
portfolio differing only in solver phase.
\item \emph{Pre-emption.} Workers poll the master every 20{,}000 decisions. If the numbers of
workers on two messages differ by two or more, the master reassigns a busy worker, and its
search on the old message is discarded.
\item \emph{Source duplication.} At the start only the source message exists, so every worker
receives it; with depth-first allocation they move to vertex-1 messages as these appear.
\end{itemize}

\subsection{Instrumentation and metrics}

The master logs, with timestamps, each message it sends to a worker and each report that a
worker has finished. Pairing each send with that worker's next report gives each message's
duration and identifies pre-empted work. The raw master log of every run is kept, compressed.

\begin{table}[H]
\centering
\small
\begin{tabular}{@{}l p{12cm}@{}}
\toprule
\code{msgs} & distinct vertex-1 messages sent, against the upper bound $D^k$ \\
\code{dup} & vertex-1 sends beyond the distinct messages: copies to idle workers, and re-sends \\
\code{start} & time to the first vertex-1 send; below 0.5\,s in every reported run, shown as 0\,s \\
\code{disp} & from then until every vertex-1 message has been sent at least once \\
\code{tail} & from then until the last \emph{completion report}. Duplicate copies still in
flight when the run ends are not waited for \\
\code{longest}, \code{p90} & longest and 90th-percentile completed message duration \\
\code{busy} & total completed message time / (103 workers $\times$ wall time) \\
\code{pre-empt} & mid-message reassignments, and the search time they discarded \\
overhead & wall $-$ (start $+$ disp $+$ tail): MPI launch, 104 processes loading the CNF,
and shutdown \\
\bottomrule
\end{tabular}
\caption{Dagster run metrics.}
\label{p8-tab:metrics}
\end{table}

\subsection{UNSAT acceptance gates}\label{p8-sec:gates}

With no proof and no status line, a Dagster \st{UNSAT-UNCERTIFIED} verdict is accepted only if:
\begin{enumerate}
\item the process exits 0 with an existing, empty solution file, without timeout or a killed rank;
\item the number of distinct vertex-1 messages equals $M(D,k)$ for the recorded split;
\item every distinct message has at least one completion report;
\item every message assigns exactly the interface variables of the recorded split, each once;
\item any worker still holding a message at shutdown holds a copy of a completed one;
\item the CNF, DAG and generator digests match the region and the release.
\end{enumerate}
Gates 2--5 are implemented by \code{hpc/audit\_dagster.py} on the raw master log. Two versions
matter.
\begin{itemize}
\item v1, the version run over the campaign logs on the cluster, checks gate 2, gate 3,
gate 5, and that all messages of a run carry the same \emph{number} of interface literals. Every
run that closed a region passes:
\begin{itemize}
\item the 242 depth-3 $f=2$ runs, $(D,k)=(29,3)$;
\item the depth-2 runs closing R00043 and R00046, $(29,2)$;
\item every $f=4$ depth-2 run, $(27,2)$;
\item all 31 cases of $R_\emptyset$, $(31,3)$.
\end{itemize}
\item v2, released with this paper, implements gate 4 in full. It rebuilds the
interface the run actually used, every edge variable incident to the split coordinates recorded
in its ledger, and requires each message to assign exactly those variables, once each. It passes
on the two archived logs examined here (677 of 677 messages match the 53-variable interface) and
fails when given the wrong coordinates.
\end{itemize}
Campaign audit with v2. Run over the archived results, v2 found 317 Dagster runs that
reported UNSAT.
\begin{itemize}
\item 291 have a master log, and every one passes: 244 $f=2$ runs (242 at depth 3, plus
R00043 and R00046 at depth 2), 16 $f=4$ runs, and all 31 cases of $R_\emptyset$.
\item The other 26 have no log. All come from the four earliest jobs, the $f=4$ scaling tests on
R00473, R00621, R01430, R02114, R08990 and R09219, which ran before raw master logs were kept;
the first job to keep them, 181304590, followed directly. All six regions were also closed by
Kissat, and the final inventory credits them to Kissat, so these runs carry no weight.
\end{itemize}
Every region credited to Dagster, and every case, therefore rests on a passing v2 audit
(see the released manifests).
Logs of runs that ended at their cap fail, as they should. The one exception is R00030's first
run, which passes: its search had in fact finished, and the process was killed during shutdown
(Section~\ref{p8-sec:f2done}).



%======================================================================
\paragraph{Knowingly inconsistent messages.} The source does not encode the matching constraints
between the outside partners of split coordinates, which keeps it simple and its message count in
closed form. Some messages therefore give two split coordinates the same outside partner, and the
terminal solver rejects them at once. At depth 3 with $D=29$, $27^3-27\cdot26\cdot25=2{,}133$ of the
$19{,}764$ messages, about $11\%$, are of this kind.

\section{Case splitting by a fixed transposition, and symmetry breaking}\label{p8-sec:cases}
The native split relies on fixed points, and one region has none. This section introduces the two devices that close it, a case split on the partner of one coordinate and reflection symmetry breaking, proves both sound, and reports their validation.

%======================================================================

The native split relies on fixed points to refute unsplit branches. The region with no fixed
points, $R_\emptyset$ (\code{R00000}), offers none: no pairing is seed-refuted, the coordinate
rule falls back at every level, and a depth-3 split of the whole region leaves 87 unsplit
branches. In a one-hour run on one node, 84 of the 103 workers spent the entire hour on those
branches and none finished. Two devices remove the problem.

\subsection{Cases on one coordinate's partner}

\begin{theorem}[Case split]\label{p8-thm:case}
For a coordinate $c\in V$, let $R_F^{(c,x)}$ be the set of members of $R_F$ containing the
transposition $(c\,x)$, for $x\in V\setminus\{c\}$. These sets are pairwise disjoint and their
union is $R_F$. Hence $R_F$ is empty iff every $R_F^{(c,x)}$ is empty.
\end{theorem}
\begin{proof}
In any $\pi\in R_F$ the free coordinate $c$ lies in exactly one transposition, $(c\,\pi(c))$.
\end{proof}

A case is encoded by adding the unit clause $x_{cx}$ to the regional CNF, and is written
\code{R00000@0-}$x$ for $c=0$. For $R_\emptyset$ this gives 31 cases, $x=1,\dots,31$.

\paragraph{The fixed transposition is a seed.} Within a case, the two dots $(c,x)$ and $(x,c)$
are known, just as fixed points are known in other regions. Definition~\ref{p8-def:seed} extends
directly: a pairing $\{y,z\}$ is seed-refuted in the case if the dots of $F$, of $(c\,x)$ and of
$(y\,z)$ together repeat a displacement, and Lemma~\ref{p8-lem:seed} applies unchanged.
Algorithm~2 runs with the extended seed and with $c,x$ removed from the candidates. In all 31
cases it finds two coordinates whose pairing is seed-refuted, so a depth-2 split inside any
case has no unsplit branches. No case admits three mutually seed-refuted coordinates, so at
depth 3 about 29 unsplit branches remain per case, now only the size of a depth-2 piece.

\subsection{Reflection symmetry breaking}

\begin{lemma}\label{p8-lem:sb}
Suppose $\rho(F)=F$, where $\rho(i)=n-1-i$. Then $R_F$ is empty iff
$R_F\cap\{\pi:\pi(0)+\pi(n-1)\le n-1\}$ is empty.
\end{lemma}
\begin{proof}
If $\pi\in R_F$, then $\sigma=\rho\pi\rho\in R_F$ by Lemma~\ref{p8-lem:refl}, and
$\sigma(0)+\sigma(n-1)=2(n-1)-(\pi(0)+\pi(n-1))$. So one of $\pi,\sigma$ satisfies the
inequality.
\end{proof}

The constraint is encoded as the binary clauses $\neg h_{0,u}\vee\neg h_{n-1,w}$ for every
pair of permitted images with $u+w>n-1$; at order 32 there are 465. It applies to
$R_\emptyset$, and is combined with the case split. Soundness of the combination follows from
the two results. If $R_\emptyset$ had a member, Lemma~\ref{p8-lem:sb} would give one, $\sigma$,
satisfying the constraint, and $\sigma$ lies in case $x=\sigma(0)$, which would therefore be
satisfiable. Inside case $x$ the constraint is simply $\pi(31)\le31-x$: a direct cut of
coordinate 31's domain, stronger for larger $x$. In case $x=31$, $\pi(31)=0$ is forced and the
constraint removes nothing.

\begin{figure}[H]
\centering
\begin{minipage}{0.47\linewidth}\centering
\begin{tikzpicture}[x=0.205cm,y=0.205cm]
\fill[gray!22] (0,1) rectangle (1,2);
\fill[gray!22] (1,0) rectangle (2,1);
\fill[gray!22] (0,2) rectangle (1,3);
\fill[gray!22] (2,0) rectangle (3,1);
\fill[gray!22] (0,3) rectangle (1,4);
\fill[gray!22] (3,0) rectangle (4,1);
\fill[gray!22] (0,4) rectangle (1,5);
\fill[gray!22] (4,0) rectangle (5,1);
\fill[gray!22] (0,5) rectangle (1,6);
\fill[gray!22] (5,0) rectangle (6,1);
\fill[gray!22] (0,6) rectangle (1,7);
\fill[gray!22] (6,0) rectangle (7,1);
\fill[gray!22] (0,7) rectangle (1,8);
\fill[gray!22] (7,0) rectangle (8,1);
\fill[gray!22] (0,8) rectangle (1,9);
\fill[gray!22] (8,0) rectangle (9,1);
\fill[gray!22] (0,9) rectangle (1,10);
\fill[gray!22] (9,0) rectangle (10,1);
\fill[gray!22] (0,10) rectangle (1,11);
\fill[gray!22] (10,0) rectangle (11,1);
\fill[gray!22] (0,11) rectangle (1,12);
\fill[gray!22] (11,0) rectangle (12,1);
\fill[gray!22] (0,12) rectangle (1,13);
\fill[gray!22] (12,0) rectangle (13,1);
\fill[gray!22] (0,13) rectangle (1,14);
\fill[gray!22] (13,0) rectangle (14,1);
\fill[gray!22] (0,14) rectangle (1,15);
\fill[gray!22] (14,0) rectangle (15,1);
\fill[gray!22] (0,15) rectangle (1,16);
\fill[gray!22] (15,0) rectangle (16,1);
\fill[gray!22] (0,16) rectangle (1,17);
\fill[gray!22] (16,0) rectangle (17,1);
\fill[gray!22] (0,17) rectangle (1,18);
\fill[gray!22] (17,0) rectangle (18,1);
\fill[gray!22] (0,18) rectangle (1,19);
\fill[gray!22] (18,0) rectangle (19,1);
\fill[gray!22] (0,19) rectangle (1,20);
\fill[gray!22] (19,0) rectangle (20,1);
\fill[gray!22] (0,20) rectangle (1,21);
\fill[gray!22] (20,0) rectangle (21,1);
\fill[gray!22] (0,21) rectangle (1,22);
\fill[gray!22] (21,0) rectangle (22,1);
\fill[gray!22] (0,22) rectangle (1,23);
\fill[gray!22] (22,0) rectangle (23,1);
\fill[gray!22] (0,23) rectangle (1,24);
\fill[gray!22] (23,0) rectangle (24,1);
\fill[gray!22] (0,24) rectangle (1,25);
\fill[gray!22] (24,0) rectangle (25,1);
\fill[gray!22] (0,25) rectangle (1,26);
\fill[gray!22] (25,0) rectangle (26,1);
\fill[gray!22] (0,26) rectangle (1,27);
\fill[gray!22] (26,0) rectangle (27,1);
\fill[gray!22] (0,27) rectangle (1,28);
\fill[gray!22] (27,0) rectangle (28,1);
\fill[gray!22] (0,28) rectangle (1,29);
\fill[gray!22] (28,0) rectangle (29,1);
\fill[gray!22] (0,29) rectangle (1,30);
\fill[gray!22] (29,0) rectangle (30,1);
\fill[gray!22] (0,30) rectangle (1,31);
\fill[gray!22] (30,0) rectangle (31,1);
\fill[gray!22] (0,31) rectangle (1,32);
\fill[gray!22] (31,0) rectangle (32,1);
\fill[gray!22] (1,0) rectangle (2,1);
\fill[gray!22] (0,1) rectangle (1,2);
\fill[gray!22] (1,2) rectangle (2,3);
\fill[gray!22] (2,1) rectangle (3,2);
\fill[gray!22] (1,3) rectangle (2,4);
\fill[gray!22] (3,1) rectangle (4,2);
\fill[gray!22] (1,4) rectangle (2,5);
\fill[gray!22] (4,1) rectangle (5,2);
\fill[gray!22] (1,5) rectangle (2,6);
\fill[gray!22] (5,1) rectangle (6,2);
\fill[gray!22] (1,6) rectangle (2,7);
\fill[gray!22] (6,1) rectangle (7,2);
\fill[gray!22] (1,7) rectangle (2,8);
\fill[gray!22] (7,1) rectangle (8,2);
\fill[gray!22] (1,8) rectangle (2,9);
\fill[gray!22] (8,1) rectangle (9,2);
\fill[gray!22] (1,9) rectangle (2,10);
\fill[gray!22] (9,1) rectangle (10,2);
\fill[gray!22] (1,10) rectangle (2,11);
\fill[gray!22] (10,1) rectangle (11,2);
\fill[gray!22] (1,11) rectangle (2,12);
\fill[gray!22] (11,1) rectangle (12,2);
\fill[gray!22] (1,12) rectangle (2,13);
\fill[gray!22] (12,1) rectangle (13,2);
\fill[gray!22] (1,13) rectangle (2,14);
\fill[gray!22] (13,1) rectangle (14,2);
\fill[gray!22] (1,14) rectangle (2,15);
\fill[gray!22] (14,1) rectangle (15,2);
\fill[gray!22] (1,15) rectangle (2,16);
\fill[gray!22] (15,1) rectangle (16,2);
\fill[gray!22] (1,16) rectangle (2,17);
\fill[gray!22] (16,1) rectangle (17,2);
\fill[gray!22] (1,17) rectangle (2,18);
\fill[gray!22] (17,1) rectangle (18,2);
\fill[gray!22] (1,18) rectangle (2,19);
\fill[gray!22] (18,1) rectangle (19,2);
\fill[gray!22] (1,19) rectangle (2,20);
\fill[gray!22] (19,1) rectangle (20,2);
\fill[gray!22] (1,20) rectangle (2,21);
\fill[gray!22] (20,1) rectangle (21,2);
\fill[gray!22] (1,21) rectangle (2,22);
\fill[gray!22] (21,1) rectangle (22,2);
\fill[gray!22] (1,22) rectangle (2,23);
\fill[gray!22] (22,1) rectangle (23,2);
\fill[gray!22] (1,23) rectangle (2,24);
\fill[gray!22] (23,1) rectangle (24,2);
\fill[gray!22] (1,24) rectangle (2,25);
\fill[gray!22] (24,1) rectangle (25,2);
\fill[gray!22] (1,25) rectangle (2,26);
\fill[gray!22] (25,1) rectangle (26,2);
\fill[gray!22] (1,26) rectangle (2,27);
\fill[gray!22] (26,1) rectangle (27,2);
\fill[gray!22] (1,27) rectangle (2,28);
\fill[gray!22] (27,1) rectangle (28,2);
\fill[gray!22] (1,28) rectangle (2,29);
\fill[gray!22] (28,1) rectangle (29,2);
\fill[gray!22] (1,29) rectangle (2,30);
\fill[gray!22] (29,1) rectangle (30,2);
\fill[gray!22] (1,30) rectangle (2,31);
\fill[gray!22] (30,1) rectangle (31,2);
\fill[gray!22] (1,31) rectangle (2,32);
\fill[gray!22] (31,1) rectangle (32,2);
\fill[gray!22] (3,0) rectangle (4,1);
\fill[gray!22] (0,3) rectangle (1,4);
\fill[gray!22] (3,1) rectangle (4,2);
\fill[gray!22] (1,3) rectangle (2,4);
\fill[gray!22] (3,2) rectangle (4,3);
\fill[gray!22] (2,3) rectangle (3,4);
\fill[gray!22] (3,4) rectangle (4,5);
\fill[gray!22] (4,3) rectangle (5,4);
\fill[gray!22] (3,5) rectangle (4,6);
\fill[gray!22] (5,3) rectangle (6,4);
\fill[gray!22] (3,6) rectangle (4,7);
\fill[gray!22] (6,3) rectangle (7,4);
\fill[gray!22] (3,7) rectangle (4,8);
\fill[gray!22] (7,3) rectangle (8,4);
\fill[gray!22] (3,8) rectangle (4,9);
\fill[gray!22] (8,3) rectangle (9,4);
\fill[gray!22] (3,9) rectangle (4,10);
\fill[gray!22] (9,3) rectangle (10,4);
\fill[gray!22] (3,10) rectangle (4,11);
\fill[gray!22] (10,3) rectangle (11,4);
\fill[gray!22] (3,11) rectangle (4,12);
\fill[gray!22] (11,3) rectangle (12,4);
\fill[gray!22] (3,12) rectangle (4,13);
\fill[gray!22] (12,3) rectangle (13,4);
\fill[gray!22] (3,13) rectangle (4,14);
\fill[gray!22] (13,3) rectangle (14,4);
\fill[gray!22] (3,14) rectangle (4,15);
\fill[gray!22] (14,3) rectangle (15,4);
\fill[gray!22] (3,15) rectangle (4,16);
\fill[gray!22] (15,3) rectangle (16,4);
\fill[gray!22] (3,16) rectangle (4,17);
\fill[gray!22] (16,3) rectangle (17,4);
\fill[gray!22] (3,17) rectangle (4,18);
\fill[gray!22] (17,3) rectangle (18,4);
\fill[gray!22] (3,18) rectangle (4,19);
\fill[gray!22] (18,3) rectangle (19,4);
\fill[gray!22] (3,19) rectangle (4,20);
\fill[gray!22] (19,3) rectangle (20,4);
\fill[gray!22] (3,20) rectangle (4,21);
\fill[gray!22] (20,3) rectangle (21,4);
\fill[gray!22] (3,21) rectangle (4,22);
\fill[gray!22] (21,3) rectangle (22,4);
\fill[gray!22] (3,22) rectangle (4,23);
\fill[gray!22] (22,3) rectangle (23,4);
\fill[gray!22] (3,23) rectangle (4,24);
\fill[gray!22] (23,3) rectangle (24,4);
\fill[gray!22] (3,24) rectangle (4,25);
\fill[gray!22] (24,3) rectangle (25,4);
\fill[gray!22] (3,25) rectangle (4,26);
\fill[gray!22] (25,3) rectangle (26,4);
\fill[gray!22] (3,26) rectangle (4,27);
\fill[gray!22] (26,3) rectangle (27,4);
\fill[gray!22] (3,27) rectangle (4,28);
\fill[gray!22] (27,3) rectangle (28,4);
\fill[gray!22] (3,28) rectangle (4,29);
\fill[gray!22] (28,3) rectangle (29,4);
\fill[gray!22] (3,29) rectangle (4,30);
\fill[gray!22] (29,3) rectangle (30,4);
\fill[gray!22] (3,30) rectangle (4,31);
\fill[gray!22] (30,3) rectangle (31,4);
\fill[gray!22] (3,31) rectangle (4,32);
\fill[gray!22] (31,3) rectangle (32,4);
\fill[gray!22] (7,0) rectangle (8,1);
\fill[gray!22] (0,7) rectangle (1,8);
\fill[gray!22] (7,1) rectangle (8,2);
\fill[gray!22] (1,7) rectangle (2,8);
\fill[gray!22] (7,2) rectangle (8,3);
\fill[gray!22] (2,7) rectangle (3,8);
\fill[gray!22] (7,3) rectangle (8,4);
\fill[gray!22] (3,7) rectangle (4,8);
\fill[gray!22] (7,4) rectangle (8,5);
\fill[gray!22] (4,7) rectangle (5,8);
\fill[gray!22] (7,5) rectangle (8,6);
\fill[gray!22] (5,7) rectangle (6,8);
\fill[gray!22] (7,6) rectangle (8,7);
\fill[gray!22] (6,7) rectangle (7,8);
\fill[gray!22] (7,8) rectangle (8,9);
\fill[gray!22] (8,7) rectangle (9,8);
\fill[gray!22] (7,9) rectangle (8,10);
\fill[gray!22] (9,7) rectangle (10,8);
\fill[gray!22] (7,10) rectangle (8,11);
\fill[gray!22] (10,7) rectangle (11,8);
\fill[gray!22] (7,11) rectangle (8,12);
\fill[gray!22] (11,7) rectangle (12,8);
\fill[gray!22] (7,12) rectangle (8,13);
\fill[gray!22] (12,7) rectangle (13,8);
\fill[gray!22] (7,13) rectangle (8,14);
\fill[gray!22] (13,7) rectangle (14,8);
\fill[gray!22] (7,14) rectangle (8,15);
\fill[gray!22] (14,7) rectangle (15,8);
\fill[gray!22] (7,15) rectangle (8,16);
\fill[gray!22] (15,7) rectangle (16,8);
\fill[gray!22] (7,16) rectangle (8,17);
\fill[gray!22] (16,7) rectangle (17,8);
\fill[gray!22] (7,17) rectangle (8,18);
\fill[gray!22] (17,7) rectangle (18,8);
\fill[gray!22] (7,18) rectangle (8,19);
\fill[gray!22] (18,7) rectangle (19,8);
\fill[gray!22] (7,19) rectangle (8,20);
\fill[gray!22] (19,7) rectangle (20,8);
\fill[gray!22] (7,20) rectangle (8,21);
\fill[gray!22] (20,7) rectangle (21,8);
\fill[gray!22] (7,21) rectangle (8,22);
\fill[gray!22] (21,7) rectangle (22,8);
\fill[gray!22] (7,22) rectangle (8,23);
\fill[gray!22] (22,7) rectangle (23,8);
\fill[gray!22] (7,23) rectangle (8,24);
\fill[gray!22] (23,7) rectangle (24,8);
\fill[gray!22] (7,24) rectangle (8,25);
\fill[gray!22] (24,7) rectangle (25,8);
\fill[gray!22] (7,25) rectangle (8,26);
\fill[gray!22] (25,7) rectangle (26,8);
\fill[gray!22] (7,26) rectangle (8,27);
\fill[gray!22] (26,7) rectangle (27,8);
\fill[gray!22] (7,27) rectangle (8,28);
\fill[gray!22] (27,7) rectangle (28,8);
\fill[gray!22] (7,28) rectangle (8,29);
\fill[gray!22] (28,7) rectangle (29,8);
\fill[gray!22] (7,29) rectangle (8,30);
\fill[gray!22] (29,7) rectangle (30,8);
\fill[gray!22] (7,30) rectangle (8,31);
\fill[gray!22] (30,7) rectangle (31,8);
\fill[gray!22] (7,31) rectangle (8,32);
\fill[gray!22] (31,7) rectangle (32,8);
\draw[gray!45,very thin] (0,0) grid (32,32);
\draw[gray,dashed] (0,0) -- (32,32);
\fill[red!80!black] (0.5,0.5) circle (0.66);
\fill[red!80!black] (1.5,1.5) circle (0.66);
\fill[red!80!black] (3.5,3.5) circle (0.66);
\fill[red!80!black] (7.5,7.5) circle (0.66);
\node[font=\tiny,below] at (0.5,0) {0};
\node[font=\tiny,left] at (0,0.5) {0};
\node[font=\tiny,below] at (4.5,0) {4};
\node[font=\tiny,left] at (0,4.5) {4};
\node[font=\tiny,below] at (8.5,0) {8};
\node[font=\tiny,left] at (0,8.5) {8};
\node[font=\tiny,below] at (12.5,0) {12};
\node[font=\tiny,left] at (0,12.5) {12};
\node[font=\tiny,below] at (16.5,0) {16};
\node[font=\tiny,left] at (0,16.5) {16};
\node[font=\tiny,below] at (20.5,0) {20};
\node[font=\tiny,left] at (0,20.5) {20};
\node[font=\tiny,below] at (24.5,0) {24};
\node[font=\tiny,left] at (0,24.5) {24};
\node[font=\tiny,below] at (28.5,0) {28};
\node[font=\tiny,left] at (0,28.5) {28};
\end{tikzpicture}
\par\small (a) region $F=\{0,1,3,7\}$
\end{minipage}\hfill
\begin{minipage}{0.47\linewidth}\centering
\begin{tikzpicture}[x=0.205cm,y=0.205cm]
\fill[red!25] (0,0) rectangle (1,1);
\fill[red!25] (1,1) rectangle (2,2);
\fill[red!25] (2,2) rectangle (3,3);
\fill[red!25] (3,3) rectangle (4,4);
\fill[red!25] (4,4) rectangle (5,5);
\fill[red!25] (5,5) rectangle (6,6);
\fill[red!25] (6,6) rectangle (7,7);
\fill[red!25] (7,7) rectangle (8,8);
\fill[red!25] (8,8) rectangle (9,9);
\fill[red!25] (9,9) rectangle (10,10);
\fill[red!25] (10,10) rectangle (11,11);
\fill[red!25] (11,11) rectangle (12,12);
\fill[red!25] (12,12) rectangle (13,13);
\fill[red!25] (13,13) rectangle (14,14);
\fill[red!25] (14,14) rectangle (15,15);
\fill[red!25] (15,15) rectangle (16,16);
\fill[red!25] (16,16) rectangle (17,17);
\fill[red!25] (17,17) rectangle (18,18);
\fill[red!25] (18,18) rectangle (19,19);
\fill[red!25] (19,19) rectangle (20,20);
\fill[red!25] (20,20) rectangle (21,21);
\fill[red!25] (21,21) rectangle (22,22);
\fill[red!25] (22,22) rectangle (23,23);
\fill[red!25] (23,23) rectangle (24,24);
\fill[red!25] (24,24) rectangle (25,25);
\fill[red!25] (25,25) rectangle (26,26);
\fill[red!25] (26,26) rectangle (27,27);
\fill[red!25] (27,27) rectangle (28,28);
\fill[red!25] (28,28) rectangle (29,29);
\fill[red!25] (29,29) rectangle (30,30);
\fill[red!25] (30,30) rectangle (31,31);
\fill[red!25] (31,31) rectangle (32,32);
\fill[gray!22] (0,0) rectangle (1,1);
\fill[gray!22] (0,0) rectangle (1,1);
\fill[gray!22] (0,1) rectangle (1,2);
\fill[gray!22] (1,0) rectangle (2,1);
\fill[gray!22] (0,2) rectangle (1,3);
\fill[gray!22] (2,0) rectangle (3,1);
\fill[gray!22] (0,3) rectangle (1,4);
\fill[gray!22] (3,0) rectangle (4,1);
\fill[gray!22] (0,4) rectangle (1,5);
\fill[gray!22] (4,0) rectangle (5,1);
\fill[gray!22] (0,5) rectangle (1,6);
\fill[gray!22] (5,0) rectangle (6,1);
\fill[gray!22] (0,6) rectangle (1,7);
\fill[gray!22] (6,0) rectangle (7,1);
\fill[gray!22] (0,7) rectangle (1,8);
\fill[gray!22] (7,0) rectangle (8,1);
\fill[gray!22] (0,8) rectangle (1,9);
\fill[gray!22] (8,0) rectangle (9,1);
\fill[gray!22] (0,9) rectangle (1,10);
\fill[gray!22] (9,0) rectangle (10,1);
\fill[gray!22] (0,10) rectangle (1,11);
\fill[gray!22] (10,0) rectangle (11,1);
\fill[gray!22] (0,11) rectangle (1,12);
\fill[gray!22] (11,0) rectangle (12,1);
\fill[gray!22] (0,12) rectangle (1,13);
\fill[gray!22] (12,0) rectangle (13,1);
\fill[gray!22] (0,13) rectangle (1,14);
\fill[gray!22] (13,0) rectangle (14,1);
\fill[gray!22] (0,14) rectangle (1,15);
\fill[gray!22] (14,0) rectangle (15,1);
\fill[gray!22] (0,15) rectangle (1,16);
\fill[gray!22] (15,0) rectangle (16,1);
\fill[gray!22] (0,16) rectangle (1,17);
\fill[gray!22] (16,0) rectangle (17,1);
\fill[gray!22] (0,17) rectangle (1,18);
\fill[gray!22] (17,0) rectangle (18,1);
\fill[gray!22] (0,18) rectangle (1,19);
\fill[gray!22] (18,0) rectangle (19,1);
\fill[gray!22] (0,19) rectangle (1,20);
\fill[gray!22] (19,0) rectangle (20,1);
\fill[gray!22] (0,20) rectangle (1,21);
\fill[gray!22] (20,0) rectangle (21,1);
\fill[gray!22] (0,21) rectangle (1,22);
\fill[gray!22] (21,0) rectangle (22,1);
\fill[gray!22] (0,22) rectangle (1,23);
\fill[gray!22] (22,0) rectangle (23,1);
\fill[gray!22] (0,23) rectangle (1,24);
\fill[gray!22] (23,0) rectangle (24,1);
\fill[gray!22] (0,24) rectangle (1,25);
\fill[gray!22] (24,0) rectangle (25,1);
\fill[gray!22] (0,25) rectangle (1,26);
\fill[gray!22] (25,0) rectangle (26,1);
\fill[gray!22] (0,26) rectangle (1,27);
\fill[gray!22] (26,0) rectangle (27,1);
\fill[gray!22] (0,27) rectangle (1,28);
\fill[gray!22] (27,0) rectangle (28,1);
\fill[gray!22] (0,28) rectangle (1,29);
\fill[gray!22] (28,0) rectangle (29,1);
\fill[gray!22] (0,29) rectangle (1,30);
\fill[gray!22] (29,0) rectangle (30,1);
\fill[gray!22] (0,30) rectangle (1,31);
\fill[gray!22] (30,0) rectangle (31,1);
\fill[gray!22] (0,31) rectangle (1,32);
\fill[gray!22] (31,0) rectangle (32,1);
\fill[gray!22] (5,0) rectangle (6,1);
\fill[gray!22] (0,5) rectangle (1,6);
\fill[gray!22] (5,1) rectangle (6,2);
\fill[gray!22] (1,5) rectangle (2,6);
\fill[gray!22] (5,2) rectangle (6,3);
\fill[gray!22] (2,5) rectangle (3,6);
\fill[gray!22] (5,3) rectangle (6,4);
\fill[gray!22] (3,5) rectangle (4,6);
\fill[gray!22] (5,4) rectangle (6,5);
\fill[gray!22] (4,5) rectangle (5,6);
\fill[gray!22] (5,5) rectangle (6,6);
\fill[gray!22] (5,5) rectangle (6,6);
\fill[gray!22] (5,6) rectangle (6,7);
\fill[gray!22] (6,5) rectangle (7,6);
\fill[gray!22] (5,7) rectangle (6,8);
\fill[gray!22] (7,5) rectangle (8,6);
\fill[gray!22] (5,8) rectangle (6,9);
\fill[gray!22] (8,5) rectangle (9,6);
\fill[gray!22] (5,9) rectangle (6,10);
\fill[gray!22] (9,5) rectangle (10,6);
\fill[gray!22] (5,10) rectangle (6,11);
\fill[gray!22] (10,5) rectangle (11,6);
\fill[gray!22] (5,11) rectangle (6,12);
\fill[gray!22] (11,5) rectangle (12,6);
\fill[gray!22] (5,12) rectangle (6,13);
\fill[gray!22] (12,5) rectangle (13,6);
\fill[gray!22] (5,13) rectangle (6,14);
\fill[gray!22] (13,5) rectangle (14,6);
\fill[gray!22] (5,14) rectangle (6,15);
\fill[gray!22] (14,5) rectangle (15,6);
\fill[gray!22] (5,15) rectangle (6,16);
\fill[gray!22] (15,5) rectangle (16,6);
\fill[gray!22] (5,16) rectangle (6,17);
\fill[gray!22] (16,5) rectangle (17,6);
\fill[gray!22] (5,17) rectangle (6,18);
\fill[gray!22] (17,5) rectangle (18,6);
\fill[gray!22] (5,18) rectangle (6,19);
\fill[gray!22] (18,5) rectangle (19,6);
\fill[gray!22] (5,19) rectangle (6,20);
\fill[gray!22] (19,5) rectangle (20,6);
\fill[gray!22] (5,20) rectangle (6,21);
\fill[gray!22] (20,5) rectangle (21,6);
\fill[gray!22] (5,21) rectangle (6,22);
\fill[gray!22] (21,5) rectangle (22,6);
\fill[gray!22] (5,22) rectangle (6,23);
\fill[gray!22] (22,5) rectangle (23,6);
\fill[gray!22] (5,23) rectangle (6,24);
\fill[gray!22] (23,5) rectangle (24,6);
\fill[gray!22] (5,24) rectangle (6,25);
\fill[gray!22] (24,5) rectangle (25,6);
\fill[gray!22] (5,25) rectangle (6,26);
\fill[gray!22] (25,5) rectangle (26,6);
\fill[gray!22] (5,26) rectangle (6,27);
\fill[gray!22] (26,5) rectangle (27,6);
\fill[gray!22] (5,27) rectangle (6,28);
\fill[gray!22] (27,5) rectangle (28,6);
\fill[gray!22] (5,28) rectangle (6,29);
\fill[gray!22] (28,5) rectangle (29,6);
\fill[gray!22] (5,29) rectangle (6,30);
\fill[gray!22] (29,5) rectangle (30,6);
\fill[gray!22] (5,30) rectangle (6,31);
\fill[gray!22] (30,5) rectangle (31,6);
\fill[gray!22] (5,31) rectangle (6,32);
\fill[gray!22] (31,5) rectangle (32,6);
\fill[orange!60] (31,27) rectangle (32,28);
\fill[orange!60] (31,28) rectangle (32,29);
\fill[orange!60] (31,29) rectangle (32,30);
\fill[orange!60] (31,30) rectangle (32,31);
\draw[gray!45,very thin] (0,0) grid (32,32);
\draw[gray,dashed] (0,0) -- (32,32);
\fill[blue!75!black] (0.5,5.5) circle (0.6);
\fill[blue!75!black] (5.5,0.5) circle (0.6);
\node[font=\tiny,below] at (0.5,0) {0};
\node[font=\tiny,left] at (0,0.5) {0};
\node[font=\tiny,below] at (4.5,0) {4};
\node[font=\tiny,left] at (0,4.5) {4};
\node[font=\tiny,below] at (8.5,0) {8};
\node[font=\tiny,left] at (0,8.5) {8};
\node[font=\tiny,below] at (12.5,0) {12};
\node[font=\tiny,left] at (0,12.5) {12};
\node[font=\tiny,below] at (16.5,0) {16};
\node[font=\tiny,left] at (0,16.5) {16};
\node[font=\tiny,below] at (20.5,0) {20};
\node[font=\tiny,left] at (0,20.5) {20};
\node[font=\tiny,below] at (24.5,0) {24};
\node[font=\tiny,left] at (0,24.5) {24};
\node[font=\tiny,below] at (28.5,0) {28};
\node[font=\tiny,left] at (0,28.5) {28};
\end{tikzpicture}
\par\small (b) region $F=\emptyset$, case $(0\,5)$
\end{minipage}
\caption{What a region looks like at order 32. (a) The four fixed points sit on the
diagonal; their rows and columns (grey) cannot hold any other dot, and the other 28
coordinates must be paired into 14 mirror-image transpositions with all displacements
distinct. (b) With no fixed points the diagonal is excluded (red). In case $(0\,5)$ the dots
$(0,5)$ and $(5,0)$ are fixed (blue), and symmetry breaking removes the orange cells:
$\pi(31)\le 26$.}
\label{p8-fig:regions32}
\end{figure}

\subsection{Validation}
\begingroup\sloppy
\begin{itemize}
\item \emph{Symmetry breaking, against brute force.} At orders 8, 10, 12 and 14 every
fixed-point-free Costas involution was enumerated directly: 4, 2, 2 and 4 of them, in mirror
pairs. In every case the CNF has exactly as many models, the CNF with the new clauses has
exactly the number satisfying the constraint (half), and every mirror pair keeps a member.
\item \emph{Case split, at order 10.} The fixed-point-free region has two solutions,
$\pi=(5,9,4,6,2,0,3,8,7,1)$ and its mirror, with $\pi(0)=5$ and $8$. Without symmetry
breaking exactly cases $(0\,5)$ and $(0\,8)$ are SAT; with it only $(0\,5)$ is, through both
Kissat and the native split. All nine case CNFs regenerate to their recorded digests.
\item \emph{Bookkeeping.} A region counts as closed by cases only when, for one coordinate,
every possible partner has an UNSAT case; one missing case leaves it open, and one SAT case
makes it SAT. These three behaviours were tested.
\end{itemize}
\endgroup

%======================================================================
\section{Results}\label{p8-sec:results}
This section reports the outcome class by class, and ends with the overall result and its accounting by method.


\subsection{Kissat on the sample}

Table~\ref{p8-tab:kissat} gives Stage~1 and the tail. In this sample the fixed-point count is the
strongest observed predictor of runtime. All $f=6$ regions closed within seconds, $f=4$
regions were intermediate, and $f=2$ regions were much harder. Variation within each class
remains material and has not been modelled. Candidate predictors include diameter, the number
of seed-refuted pairs, and the number of unavoidable unsplit branches.

\begin{table}[H]
\centering
\begin{tabular}{@{}rrrrrrr@{}}
\toprule
$f$ & sampled & closed in 600\,s & median & quartiles & max & closed in 4\,h tail \\
\midrule
6 & 100 & 100 & 7.7\,s & 6\,s / 11\,s & 33.9\,s & --- \\
4 & 100 & 78 & 380\,s & 317\,s / 440\,s & 592\,s & 22 of 22 \\
2 & 24 & 0 & --- & --- & --- & 13 of 24 \\
0 & 1 & 0 & --- & --- & --- & 0 of 1 \\
\bottomrule
\end{tabular}
\caption{Kissat on one core. Medians, quartiles and maxima are over regions closed within
600\,s. Every completed run was \st{UNSAT-UNCERTIFIED}, and there were no errors. In the tail,
the $f=4$ regions closed with median 1{,}066\,s and the 13 closed $f=2$ regions with median
11{,}135\,s.}
\label{p8-tab:kissat}
\end{table}

\subsection{Dagster, $f=4$, depth 2}

Two of the hardest $f=4$ regions Kissat closed in Stage~1 were rerun with Dagster, so each
Dagster time has a Kissat baseline.

\begin{table}[H]
\centering
\small
\begin{tabular}{@{}lrrrrrrrrrr@{}}
\toprule
region & Kissat & Dagster & msgs & disp & tail & longest & p90 & busy & wall gain & core cost \\
\midrule
R02114 & 584\,s & 39\,s & 677/729 & 10\,s & 7\,s & 9\,s & 4\,s & 31\% & 15.1$\times$ & 6.9$\times$ \\
R09219 & 592\,s & 37\,s & 677/729 & 9\,s & 5\,s & 9\,s & 4\,s & 30\% & 16.0$\times$ & 6.5$\times$ \\
\bottomrule
\end{tabular}
\caption{$F(\code{R02114})=\{0,7,15,28\}$, $F(\code{R09219})=\{3,7,10,26\}$. Core cost is
$104\times$ the Dagster wall time divided by Kissat's single-core time. Overhead was about 22\,s
in each run.}
\label{p8-tab:f4}
\end{table}

\subsection{Dagster, $f=2$}

All 11 sampled $f=2$ regions that Kissat left open after 4\,h were rerun with Dagster.
Every one was search-closed.

\begin{table}[H]
\centering
\small
\begin{tabular}{@{}lrrrrrrrr@{}}
\toprule
region & depth & wall & disp & tail & longest & p90 & busy & pre-empted \\
\midrule
R00046 & 2 & 2{,}849\,s & 2{,}118\,s & 694\,s & 1{,}071\,s & 665\,s & 85\% & 34{,}378\,s \\
R00043 & 2 & 3{,}086\,s & 2{,}283\,s & 763\,s & 1{,}142\,s & 755\,s & 85\% & 36{,}684\,s \\
\midrule
R00149 & 3 & 277\,s & 250\,s & 5\,s & 99\,s & 3\,s & 91\% & 120\,s \\
R00121 & 3 & 397\,s & 364\,s & 10\,s & 226\,s & 4\,s & 93\% & 420\,s \\
R00097 & 3 & 469\,s & 437\,s & 9\,s & 312\,s & 5\,s & 94\% & 311\,s \\
R00069 & 3 & 611\,s & 576\,s & 11\,s & 493\,s & 6\,s & 95\% & 453\,s \\
R00068 & 3 & 681\,s & 644\,s & 12\,s & 587\,s & 7\,s & 95\% & 500\,s \\
R00254 & 3 & 726\,s & 532\,s & 160\,s & 693\,s & 5\,s & 74\% & 1{,}444\,s \\
R00043 & 3 & 872\,s & 625\,s & 211\,s & 836\,s & 7\,s & 73\% & 10{,}814\,s \\
R00114 & 3 & 882\,s & 846\,s & 14\,s & 684\,s & 10\,s & 96\% & 650\,s \\
R00063 & 3 & 1{,}496\,s & 1{,}243\,s & 217\,s & 1{,}460\,s & 15\,s & 84\% & 9{,}706\,s \\
R00087 & 3 & 2{,}306\,s & 1{,}271\,s & 986\,s & 2{,}257\,s & 12\,s & 57\% & 36{,}084\,s \\
\bottomrule
\end{tabular}
\caption{103 workers. Depth 2 sends 785 messages and depth 3 sends 19{,}764.
$F(\code{R00043})=\{1,13\}$, $F(\code{R00046})=\{1,16\}$. Kissat did not close any of these
regions within 14{,}400\,s.}
\label{p8-tab:f2}
\end{table}
\subsection{The $f=2$ class: complete}\label{p8-sec:f2done}

Besides the 24 sampled regions, the other 232 ran at depth 3 in 29 batches of eight, in region order.
The first five batches, \emph{wave~1}, held R00001--R00040. One region, R00030, reached the 1-hour
cap and was re-run with a 3-hour cap.

\begin{itemize}
\item All 256 $f=2$ regions are \st{search-closed}, and none is satisfiable: 13 by
Kissat within 4\,h, and 243 by Dagster (242 at depth 3, R00046 at depth 2). The cluster's
campaign status reports 256 closed, 0 open, 0 failed and 0 never run.
\item All 192 regions in batches 005--028 closed within the 1-hour cap, the slowest
being $\{1,29\}$ in 2{,}871\,s. R00030 closed on re-run in 3{,}554\,s, making $\{0,30\}$ the
hardest $f=2$ region.
\item Cost: batches 005--028 used 6{,}410 SU, about 33 SU per region, half the wave-1
rate. With wave~1 and R00030's re-run, the campaign used 9{,}212 SU for 232 regions.
\item Audit: every depth-3 master log on the cluster passes gates 2--5.
\end{itemize}

Subject to the claim boundary, this closes one fixed-point class within the present search:
\begin{quote}
\emph{No transpose-symmetric Costas array of order 32 has exactly two fixed points}
(computational, \st{search-closed}; not yet certified).
\end{quote}

\begin{figure}[H]
\centering
\begin{tikzpicture}
\begin{axis}[width=13cm, height=6.2cm, xlabel={separation $b-a$ of the fixed points $\{a,b\}$},
  ylabel={Dagster wall time (s)}, xmin=0, xmax=32, ymin=0, ymax=3800, grid=major,
  colormap/viridis, colorbar, colorbar style={title={$a$}, width=3mm},
  point meta min=0, point meta max=15, tick label style={font=\small}]
\addplot[scatter, only marks, mark=*, mark size=1.3pt, point meta=explicit] coordinates {(1,2225) [0] (2,2084) [0] (3,1883) [0] (4,1493) [0] (5,1278) [0] (6,1171) [0] (7,751) [0] (8,627) [0] (9,539) [0] (10,467) [0] (11,409) [0] (12,375) [0] (13,347) [0] (14,342) [0] (15,328) [0] (16,321) [0] (17,316) [0] (18,309) [0] (19,289) [0] (20,295) [0] (21,311) [0] (22,356) [0] (23,481) [0] (24,519) [0] (25,703) [0] (26,822) [0] (27,1071) [0] (28,1723) [0] (29,2194) [0] (30,3554) [0] (31,3315) [0] (1,2043) [1] (2,1858) [1] (3,1838) [1] (4,1586) [1] (5,1719) [1] (6,1629) [1] (7,1356) [1] (8,1028) [1] (9,798) [1] (10,734) [1] (11,674) [1] (12,872) [1] (13,716) [1] (14,765) [1] (16,599) [1] (17,561) [1] (18,554) [1] (19,520) [1] (20,747) [1] (22,961) [1] (23,1089) [1] (24,1092) [1] (25,1401) [1] (26,1772) [1] (27,2095) [1] (28,2871) [1] (29,2563) [1] (1,1945) [2] (2,1436) [2] (3,1496) [2] (4,855) [2] (5,1111) [2] (6,915) [2] (7,727) [2] (8,681) [2] (9,611) [2] (10,583) [2] (11,537) [2] (12,526) [2] (13,495) [2] (14,502) [2] (15,473) [2] (16,468) [2] (18,458) [2] (21,769) [2] (22,1016) [2] (23,1232) [2] (24,1727) [2] (25,1896) [2] (26,2421) [2] (27,2306) [2] (1,1417) [3] (2,1192) [3] (3,630) [3] (4,904) [3] (5,757) [3] (6,512) [3] (7,570) [3] (8,522) [3] (9,491) [3] (10,469) [3] (11,431) [3] (12,422) [3] (13,416) [3] (14,402) [3] (15,380) [3] (16,395) [3] (17,388) [3] (19,447) [3] (20,698) [3] (21,858) [3] (22,1378) [3] (23,1591) [3] (24,2199) [3] (25,1984) [3] (1,1287) [4] (2,882) [4] (3,750) [4] (4,428) [4] (5,540) [4] (6,509) [4] (7,452) [4] (8,365) [4] (9,397) [4] (10,377) [4] (11,351) [4] (12,364) [4] (13,335) [4] (14,320) [4] (15,312) [4] (16,326) [4] (17,333) [4] (18,365) [4] (19,411) [4] (20,596) [4] (21,859) [4] (22,1131) [4] (23,1474) [4] (1,844) [5] (2,711) [5] (3,600) [5] (4,523) [5] (5,330) [5] (6,413) [5] (7,375) [5] (8,356) [5] (9,327) [5] (10,329) [5] (12,294) [5] (13,278) [5] (14,277) [5] (15,278) [5] (16,288) [5] (17,305) [5] (18,347) [5] (19,388) [5] (20,482) [5] (21,654) [5] (1,693) [6] (2,578) [6] (4,435) [6] (5,389) [6] (6,341) [6] (7,319) [6] (8,308) [6] (9,282) [6] (10,283) [6] (11,259) [6] (12,333) [6] (13,238) [6] (14,236) [6] (15,240) [6] (16,260) [6] (17,288) [6] (18,332) [6] (19,374) [6] (1,557) [7] (3,428) [7] (5,349) [7] (6,332) [7] (7,458) [7] (8,275) [7] (9,261) [7] (10,250) [7] (11,229) [7] (12,224) [7] (13,214) [7] (14,409) [7] (15,231) [7] (16,260) [7] (17,283) [7] (1,461) [8] (2,402) [8] (3,372) [8] (4,339) [8] (5,320) [8] (6,304) [8] (7,281) [8] (8,572) [8] (9,247) [8] (11,206) [8] (12,203) [8] (13,198) [8] (15,229) [8] (1,392) [9] (2,360) [9] (3,338) [9] (4,326) [9] (5,312) [9] (6,298) [9] (7,285) [9] (8,265) [9] (10,213) [9] (11,197) [9] (12,192) [9] (13,195) [9] (1,351) [10] (2,334) [10] (3,333) [10] (4,317) [10] (5,308) [10] (6,308) [10] (7,286) [10] (8,259) [10] (9,238) [10] (10,714) [10] (11,198) [10] (1,352) [11] (2,354) [11] (3,359) [11] (4,352) [11] (6,340) [11] (7,305) [11] (8,285) [11] (9,261) [11] (1,388) [12] (2,392) [12] (3,396) [12] (4,403) [12] (5,389) [12] (6,350) [12] (7,337) [12] (1,457) [13] (2,467) [13] (3,471) [13] (4,454) [13] (5,425) [13] (1,561) [14] (2,726) [14] (3,557) [14] (1,715) [15]};
\end{axis}
\end{tikzpicture}
\caption{All 242 depth-3 Dagster runs on $f=2$ regions: wall time against the separation of
the two fixed points, coloured by the smaller fixed point $a$. Within each row the shape is
a U shape. Rows with larger $a$ are cheaper and span fewer separations,
since canonical sets satisfy $a+b\le 31$.}
\label{p8-fig:f2all}
\end{figure}

\paragraph{The seven fully split regions.} From the fixed points alone, Algorithm~2 finds three
split coordinates with all pairings seed-refuted in exactly seven $f=2$ regions: R00033,
R00064, R00093, R00120, R00145, R00168 and R00189. These are exactly the seven regions whose
longest message was conspicuously short, 10--52\,s, where comparable regions ran hundreds to
thousands of seconds: none is missing and none is extra. The regions without unsplit branches
are exactly those without long messages. This confirms the account of Section~\ref{p8-sec:rule}:
in this decomposition, unsplit branches set the critical path.

\paragraph{R00030's first run.} That run was recorded \st{TIMEOUT}, yet its master log passes
gates 2--5. Its dispatch (2{,}090\,s) and tail (1{,}488\,s) place the last completed message at
about 3{,}578\,s, so the search had finished before the 3{,}600\,s cap, and the process was
killed during Dagster's shutdown, which took 20--60\,s in every run. Gate~1, normal exit, is what
correctly kept it from being accepted. The re-run then completed normally.

\subsection{The $f=6$ and $f=4$ classes}\label{p8-sec:bulk}

\begin{table}[H]
\centering
\begin{tabular}{@{}lrrrrr@{}}
\toprule
batch & regions & search-closed & timeouts & median & SU \\
\midrule
$f=6$, 000 & 10{,}400 & 10{,}400 & 0 & 13.3\,s & 92.91 \\
$f=6$, 001 & 10{,}400 & 10{,}400 & 0 & 9.7\,s & 66.91 \\
$f=6$, 002 & 10{,}203 & 10{,}203 & 0 & 7.1\,s & 49.05 \\
\midrule
$f=4$, 000 & 2{,}496 & 2{,}486 & 10 & 822\,s & 1{,}372.22 \\
$f=4$, 001 & 2{,}496 & 2{,}494 & 2 & 674\,s & 1{,}098.30 \\
$f=4$, 002 & 2{,}496 & 2{,}496 & 0 & 552\,s & 906.82 \\
$f=4$, 003 & 2{,}496 & 2{,}496 & 0 & 451\,s & 731.76 \\
$f=4$, 004 & 2{,}496 & 2{,}496 & 0 & 359\,s & 545.65 \\
$f=4$, 005 & 1{,}007 & 1{,}007 & 0 & 217\,s & 129.54 \\
\bottomrule
\end{tabular}
\caption{Kissat class batches (one core per region). With the 100 sampled regions of each
class, all 31{,}103 $f=6$ regions are search-closed, and 13{,}575 of the 13{,}587 $f=4$
regions; the other 12 went to Dagster.}
\label{p8-tab:bulk}
\end{table}

Median times fall steadily from batch to batch. Regions are numbered in lexicographic order of
$F$, so the earliest batches hold the sets whose marks start near 0, which pack the fixed
points into one corner of the board. These are hardest, as they are for $f=2$
(Figure~\ref{p8-fig:f2all}).

\paragraph{The 12 $f=4$ stragglers.} All 12 have their four fixed points within a span of 6 to
9, among the most compact four-mark Golomb rulers. Two are the optimal ruler $\{0,1,4,6\}$
and its shift $\{1,2,5,7\}$. For every one of them, Algorithm~2 found a second coordinate whose
pairing with the first is seed-refuted, so the depth-2 split had no unsplit branch. Dagster
closed all 12 in 255--541\,s each (Table~\ref{p8-tab:f4strag}), against more than an hour for
Kissat on one core.

\begin{table}[H]
\centering
\small
\begin{tabular}{@{}llrrrr|llrrrr@{}}
\toprule
region & $F$ & wall & longest & p90 & busy & region & $F$ & wall & longest & p90 & busy \\
\midrule
R00257 & \{0,1,3,7\} & 418 & 247 & 99 & 57\% & R00307 & \{0,1,5,8\} & 266 & 132 & 58 & 56\% \\
R00258 & \{0,1,3,8\} & 334 & 177 & 73 & 56\% & R00608 & \{0,2,3,7\} & 466 & 275 & 108 & 56\% \\
R00259 & \{0,1,3,9\} & 255 & 136 & 60 & 61\% & R00609 & \{0,2,3,8\} & 362 & 181 & 82 & 56\% \\
R00282 & \{0,1,4,6\} & 541 & 299 & 124 & 54\% & R00633 & \{0,2,5,6\} & 494 & 291 & 116 & 54\% \\
R00283 & \{0,1,4,9\} & 259 & 118 & 60 & 56\% & R03729 & \{1,2,5,7\} & 412 & 224 & 96 & 56\% \\
R00306 & \{0,1,5,7\} & 343 & 199 & 82 & 57\% & R04006 & \{1,3,4,8\} & 373 & 207 & 92 & 54\% \\
\bottomrule
\end{tabular}
\caption{The 12 $f=4$ regions Kissat left open, closed by Dagster at depth 2 (times in
seconds; 677 messages each). They used 263 SU in all.}
\label{p8-tab:f4strag}
\end{table}

\subsection{The region with no fixed points}\label{p8-sec:f0}

$R_\emptyset$ was decided through its 31 cases (Section~\ref{p8-sec:cases}), with symmetry
breaking, by Dagster at depth 3: 24{,}476 messages per case, $=M(31,3)$. A depth-2 trial on
cases $(0\,1)$--$(0\,3)$ had first failed to finish in an hour. Every one of its 901 messages
was dispatched, but the 90th-percentile message took about 800\,s. At depth 3, case $(0\,1)$
took 933\,s and about a quarter of the total work, with the 90th percentile at 10\,s, the same
effect of smaller pieces seen on $f=2$.

\begin{figure}[H]
\centering
\begin{tikzpicture}
\begin{axis}[width=13cm, height=5.4cm, xlabel={case: partner $x=\pi(0)$},
  ylabel={wall time (s)}, xmin=0, xmax=32, ymin=0, ymax=2000, grid=major, ybar, bar width=4.5pt,
  tick label style={font=\small}]
\addplot[fill=blue!55, draw=blue!70!black] coordinates {(1,933) (2,922) (3,1003) (4,846) (5,742) (6,665) (7,569) (8,491) (9,427) (10,369) (11,331) (12,303) (13,288) (14,281) (15,268) (16,264) (17,249) (18,247) (19,247) (20,278) (21,267) (22,281) (23,321) (24,448) (25,627) (26,882) (27,1232) (28,1271) (29,1254) (30,49) (31,1801)};
\end{axis}
\end{tikzpicture}
\caption{The 31 cases of $R_\emptyset$ (depth 3, 103 workers). Cases 1--23 fall as symmetry
breaking ($\pi(31)\le31-x$) cuts deeper. Cases 24--29 are dominated by one long message
while almost every other piece is trivial: their 90th-percentile message took 0--2\,s.
Case 30 takes 49\,s. In case 31, $\pi(31)=0$ is forced, so symmetry breaking removes
nothing; the transposition $(0\,31)$ is itself mirror-symmetric, so that case still has the full
reflection symmetry.}
\label{p8-fig:f0}
\end{figure}

All 31 cases are \st{UNSAT-UNCERTIFIED}, so $R_\emptyset$ is \st{search-closed}
(Theorem~\ref{p8-thm:case}). The status tool reports \emph{all 31 cases on coordinate 0
$\to$ closed}, and every case log passes the message audit with $(D,k)=(31,3)$. The 31
cases took about 5.1 node-hours and 1{,}057 SU.

\subsection{The result}\label{p8-sec:result}

\begin{table}[H]
\centering
\begin{tabular}{@{}rrrrr@{}}
\toprule
$f$ & regions & closed by Kissat & closed by Dagster & satisfiable \\
\midrule
6 & 31{,}103 & 31{,}103 & --- & 0 \\
4 & 13{,}587 & 13{,}575 & 12 & 0 \\
2 & 256 & 13 & 243 & 0 \\
0 & 1 & --- & 1 (31 cases) & 0 \\
\midrule
all & 44{,}947 & 44{,}691 & 256 & 0 \\
\bottomrule
\end{tabular}
\caption{Every region is \st{search-closed}, confirmed for each class by the status tool on
the cluster. Regions closed by both methods are counted once, under the first method listed.}
\label{p8-tab:final}
\end{table}

\noindent In terms of methods,
\[
\underbrace{44{,}691}_{\text{Kissat}}
+\underbrace{243}_{\text{Dagster, }f=2}
+\underbrace{12}_{\text{Dagster, }f=4}
+\underbrace{1}_{\text{Dagster, 31 cases, }f=0}
=44{,}947 .
\]
The 31 cases are not counted as regions.

\begin{quote}
Computational result. Under the released generator, the exact regional encoding, the
accepted solver rules and the audited execution protocol, all 44{,}947 canonical regions are
\st{search-closed} and no region contains a Costas array. By Theorem~\ref{p8-thm:partition},
\emph{there is no transpose-symmetric Costas array of order 32}. This agrees with the published
enumerations of Section~\ref{p8-sec:censuses}, which the search confirms independently.

Certification status. The negative verdicts are not supported by independently
checked UNSAT proofs. The result is a complete computational exclusion under the trust
assumptions of Section~\ref{p8-sec:claim}, not a formally certified theorem.
\end{quote}

Equivalently: no Costas array of order 32 has any non-trivial symmetry
(Proposition~\ref{p8-prop:onlytranspose}).

\section{Performance analysis}\label{p8-sec:perf}
This section examines where time went in the distributed runs and what the campaign revealed about the decomposition, keeping latency and aggregate cost separate.

%======================================================================

\subsection{Where the wall time goes}
\begin{itemize}
\item \emph{$f=4$, depth 2.} Dispatch takes about 10\,s and the tail 5--7\,s; the longest message
is 9\,s. The other $\approx$22\,s of a 37--39\,s run is overhead, more than half the wall time.
Busy is about 30\%: at this size a full node exceeds what one region needs.
\item \emph{$f=2$, depth 3.} Most regions are throughput-bound, with busy 91--96\% and a tail of
5--14\,s. In R00043, R00063 and R00087 the longest message is almost dispatch plus tail: one
message sets the wall time, most probably one of the 27 unsplit branches (the correlation strongly implicates them; mapping
each longest message back to its partner assignment would settle it) of depth-2 size
(Section~\ref{p8-sec:rule}). In R00087 that message ran 2{,}257\,s, busy fell to 57\%, and
36{,}084\,s of search was discarded by pre-emption while idle workers raced copies of it.
\end{itemize}

\subsection{Depth}
On R00043, depth 3 cut the wall time 3.5$\times$, the total message work about 4$\times$, and
the pre-empted time 3.4$\times$. The 90th-percentile message fell from 755\,s to 7\,s. The fall
in total work is consistent with TiniSAT keeping pace with Kissat on small instances and falling
behind on large ones; on a 20\,s child of an $f=4$ region the two took 19.8\,s and 19.9\,s.

\subsection{Against Kissat}
For a region taking $T_K$ seconds with Kissat on one core and $T_D$ with Dagster on a
104-core node, the wall-time gain is $T_K/T_D$ and the core-cost ratio is $104\,T_D/T_K$.
Dagster lowers latency on hard regions at an increased aggregate core cost.

\begin{table}[H]
\centering
\begin{tabular}{@{}lrrrr@{}}
\toprule
& Kissat, 1 core & Dagster, 104 cores & wall-time gain & cost \\
\midrule
$f=4$ (R02114, R09219) & 584--592\,s & 37--39\,s & 15--16$\times$ & 6.5--6.9$\times$ core-s \\
$f=2$ hard (10, depth 3) & $>$14{,}400\,s & 277--2{,}306\,s & $>$6--52$\times$ & $\approx$51 SU \\
\bottomrule
\end{tabular}
\caption{For the $f=2$ regions Kissat's true times are unknown beyond 4\,h, so gains are lower bounds.}
\label{p8-tab:vs}
\end{table}

The Dagster configuration measured here gives much lower wall time on the hard regions tested.
The comparison does not isolate the decomposition from the core count or from the
difference between solvers (TiniSAT under Dagster, Kissat alone). Isolating them would need
the same leaf solver run outside Dagster over the same messages. For throughput across many
easy regions, Kissat on one core per region remains cheaper per region.



\subsection{What the decomposition taught}\label{p8-sec:lessons}

\paragraph{Unsplit branches set the critical path.} A message in which two split coordinates
pair with each other is a whole level shallower than the rest. In the $f=2$ class, Algorithm~2
finds a fully seed-refuted triple in exactly seven regions (R00033, R00064, R00093, R00120,
R00145, R00168, R00189). These were exactly the seven regions whose longest message was short,
10--52\,s, while comparable regions ran hundreds to thousands of seconds; none was missing and
none extra. In $R_\emptyset$ at depth 3, without the case split, 84 of the 103 workers spent
the whole hour on unsplit branches. Choosing coordinates whose pairings the seed refutes is
the single most important design rule.

\paragraph{Smaller pieces, less work.} Dagster's TiniSAT workers keep pace with Kissat on
small subproblems but fall behind on large ones. Going one level deeper reduced the
\emph{total} work about fourfold twice: on the $f=2$ region R00043 (about 75 to 18
core-hours), and on case $(0\,1)$ of $R_\emptyset$.

\paragraph{Order matters, structure decides.} Within every class, the compact fixed-point
sets near one corner of the board are the hardest; separations around 20 in the $f=2$ class
are about seven times cheaper. Difficulty is governed by structure, not region size: the
formula sizes within a class vary by only a few per cent.

\subsection{Comparison with a direct search}\label{p8-sec:direct}

The comparisons above are against Kissat on single regions. A more basic question is how the whole
campaign compares with a direct search of the kind used historically, which needs no SAT solver at
all. This project has such a search: a backtracking enumeration of involutions that fixes each
coordinate's partner in turn and rejects repeated displacements incrementally with bitmasks. It
uses the same principle as the Ladendorf midpoint test of the 2009 enumeration, though it is
implemented differently and without symmetry reduction. We timed it on a single 2.8\,GHz Intel Xeon core of a virtual machine. Clock frequency and turbo
behaviour were controlled neither here nor on the cluster, so absolute times are indicative; the
conclusions rest on ratios far larger than such effects.

\begin{table}[H]
\centering
\begin{tabular}{@{}rrrrrr@{}}
\toprule
order & nodes & ratio & time & arrays & twice the 2009 classes \\
\midrule
16 & 160{,}797 & & 0.05\,s & 40 & 40 \\
18 & 864{,}774 & 5.38 & 0.26\,s & 20 & 20 \\
20 & 4{,}710{,}666 & 5.45 & 1.57\,s & 8 & 8 \\
22 & 25{,}853{,}570 & 5.49 & 9.56\,s & 10 & 10 \\
24 & 142{,}822{,}399 & 5.52 & 58.8\,s & 0 & 0 \\
26 & 791{,}747{,}546 & 5.54 & 382\,s & 4 & 4 \\
\bottomrule
\end{tabular}
\caption{The direct involution search on one core. Node counts agree with those recorded earlier in
the project by three independent implementations. Array counts are exactly twice the 2009 class
counts~\cite{p8-taylor2009}. Ratios are between successive even orders.}
\label{p8-tab:direct}
\end{table}

The node count grows by a factor of about 5.5 per step of two orders, rising slowly, and the time by
about 6.5. Order 32 is three such steps beyond the last measured order. The projection is
about $1.4\times10^{11}$ nodes (between $1.3$ and $1.7\times10^{11}$ for plausible ratios), and 30 to 45
core-hours on this core.

\paragraph{Measured.} The direct search was then run to completion on one 104-core node of our
environment. It was split by the partners of the first three coordinates into 21{,}727 independent tasks,
and preceded by a self-test in which the split reproduced the unsplit search's node and array counts
exactly at orders 16 to 24; at orders 16 to 22 its arrays are also identical, as sorted lists, to those of an
independent unsplit search. Every task completed, and the fail-closed merge found no transpose-symmetric
Costas array of order 32 among $135{,}486{,}910{,}196$ nodes, within $5\%$ of the projection. The search took
$62{,}164$ core-seconds, 17.3 core-hours, and 22 minutes of wall time: faster than projected. It is a
fourth determination of the result, after the enumerations of 1993 and 2009 and the SAT campaign of this
paper, and independent of the last in method and software.

Our campaign's listed production runs used 8{,}443 core-hours (Table~\ref{p8-tab:costs}). For this
problem, then, the direct search needed about 490 times less computation. On latency, it divides
trivially by the partners of the first few coordinates, and on the 104 cores of one node it finished in
22 minutes, against hours of wall time for the campaign. The large gains reported in Table~\ref{p8-tab:vs} are relative to single-core Kissat on the
hardest regions, not to a direct search.

The historical runtime cannot be compared numerically. The 2009 paper reports approximately one year
of processing for its whole extension, covering orders 28 to 36, several symmetry types and the
hardware of the time, without separating order 32. At about $6.5\times$ per two orders, order 36 alone
costs some forty times order 32, so a total of that magnitude is consistent with the measurements
here.

Two conclusions follow.
\begin{itemize}
\item For this problem the decomposition does not save computation. What it adds is an exact
encoding, an audited region-by-region record, and a route to proof certificates. It also offers a
method intended for problems where direct search does not scale, such as structured subclasses of
the general order-32 problem.
\item An earlier projection in this project, $6\times10^{14}$ nodes at order 32 (``a few
core-years''), was wrong. It applied the growth factor once per order rather than once per two orders,
overstating the cost by a factor of about four thousand.
\end{itemize}

\paragraph{Note: latency in search steps, under a common decomposition.}
Suppose the direct search were split into exactly the same pieces as a Dagster run, every piece
had its own processor, and only search steps were counted, ignoring input and output, start-up
and the cost of each step. Both methods would then finish when their hardest piece finishes, so
the comparison comes down to how many steps each needs on a single piece.

\begin{proposition}\label{p8-prop:simulation}
Fix a piece, and consider a \emph{multiway propagation search} on its regional formula. At each node
it selects the coordinate the direct search would select, creates one child for each candidate
partner, asserts that partner's edge variable, and applies unit propagation, never branching on
auxiliary variables. Every node it visits is also visited by the direct search.
\end{proposition}
\begin{proof}[Proof sketch]
Every rejection by the direct search is a repeated displacement among placed dots. In the
formula, deciding $x_{cr}$ forces $h_{cr}$ and $h_{rc}$ (clause family~1). With the image literals
of the existing dots true, the indicator clauses (family~5) force the displacement indicator of
each new pair of dots, and a repeated displacement falsifies an at-most-one clause (family~6).
A repeat of a transposition's own displacement $(d,-d)$ is caught the same way, or by the
constraint on transposition lengths (family~4). So each rejection is a propagation conflict at the same node, and by induction on depth the
propagation search visits only nodes that the direct search visits.
\end{proof}

\noindent The proposition concerns this multiway procedure. An ordinary DPLL or CDCL solver
branches on individual edge literals, and on others, so its tree is not directly comparable; what
carries over to it is the fact used in the proof, that every immediate displacement rejection of the
direct search is a unit-propagation conflict.

Counted in steps, then, the propagation search never needs more nodes than the direct search, and it needs fewer
whenever propagation detects a failure earlier. The main such case is a coordinate left with no
possible partner. A partner whose dots would repeat a displacement has its edge variable propagated
false, and once all are false the coordinate's at-least-one clause fails, before the search ever
branches on that coordinate. Adding exactly this check to the direct search measures how much it
is worth.

\begin{table}[H]
\centering
\begin{tabular}{@{}rrrr@{}}
\toprule
order & nodes, plain & nodes, with the check & reduction \\
\midrule
12 & 5{,}990 & 2{,}312 & 2.6$\times$ \\
14 & 30{,}510 & 10{,}944 & 2.8$\times$ \\
16 & 160{,}797 & 54{,}760 & 2.9$\times$ \\
18 & 864{,}774 & 282{,}811 & 3.1$\times$ \\
20 & 4{,}710{,}666 & 1{,}484{,}965 & 3.2$\times$ \\
\bottomrule
\end{tabular}
\caption{The direct search with and without rejecting any state in which an unmatched coordinate
has no compatible partner. Array counts are identical in both versions.}
\label{p8-tab:wipeout}
\end{table}

The reduction grows ever more slowly, by 0.20, 0.15, 0.12 and then 0.11 per step, which suggests
about $4\times$ at order 32. This measures one inference, not everything that propagation and clause learning can do, so it is
not an upper bound on SAT's advantage, and it does not establish a ratio at order 32. On its own,
this inference suggests a factor of about four. An exponential separation of the kind Rintanen
established between SAT-based planning and A$^*$-style state-space search~\cite{p8-rintanen2011}
would need the encoding to detect future impossibilities long before the direct search reaches
them. The measurements show no sign of that at these orders, although a worst-case family with
such behaviour could coexist with modest averages.

In practice the steps do not cost the same. The direct search processes about 2.4 million nodes
per second with bit operations, while each SAT step propagates through a formula of half a million
clauses or more. That difference far exceeds a factor of four, which is why the direct search was
measured as about five hundred times cheaper overall. Even within the direct search, this check
costs more time than it saves at these orders: the run at order 20 takes 2.6 times longer while
visiting a third as many nodes.

Two qualifications apply. The check measures one inference, empty partner domains, which is the
principal one available here; clause learning and the CDCL solver's own branching order lie outside
the proposition and could in principle prune further, although learnt clauses are not shared
between the pieces of a split. And the measurement stops at order 20.


%======================================================================
\section{Cost}\label{p8-sec:cost}
This section accounts for the computational resources the campaign used.

%======================================================================

\begin{table}[H]
\centering
\begin{tabular}{@{}lrl@{}}
\toprule
work & SU & notes \\
\midrule
$f=6$, Kissat & 209 & 3 batches \\
$f=4$, Kissat & 4{,}784 & 6 batches \\
$f=4$, Dagster stragglers & 263 & 12 regions \\
$f=2$, Dagster & 10{,}062 & 243 regions: 11 in the sample, 232 in the campaign \\
$f=0$, Dagster cases & 1{,}057 & 31 cases \\
Kissat tail runs (sample) & 512 & 47 regions at 4\,h \\
\midrule
production total & 16{,}887 & \\
\midrule
exploratory runs & $\approx$1{,}530 & scaling tests, depth trials, the $R_\emptyset$ pilots \\
\bottomrule
\end{tabular}
\caption{Service units in the HPC environment. The listed production runs consumed
16{,}887 SU; the whole research campaign used more, and no complete total is given because
some early and diagnostic runs lack accounting. Charges are about 2 SU per core-hour, on all requested cores for the wall time used. The Stage-1 screen of the sample, the restart sweep and a few early
diagnostic runs are not included; their resource summaries were not recorded.}
\label{p8-tab:costs}
\end{table}

%======================================================================
\section{Claim boundary and trust assumptions}\label{p8-sec:claim}
This section states precisely what the result establishes and what it does not, and lists the assumptions on which it rests.


\paragraph{Relation to the literature.} The non-existence verified here follows from the
exhaustive enumerations of 1993 and 2009 (Section~\ref{p8-sec:censuses}), so
the present work is an independent confirmation, not a first proof. Its value lies in the
independence of the confirmation, since the two methods share no software and no approach, and
in the decomposition that made it possible.

\paragraph{Scope.} The result concerns transpose-symmetric arrays and therefore, by
Proposition~\ref{p8-prop:onlytranspose}, every array invariant under a non-identity motion of the
square. Whether a Costas array of order 32
exists at all remains open.

\paragraph{Two layers.} The \emph{computational result} is that every canonical region is
\st{search-closed} and none contains a validated array. The \emph{certification status} is
that no UNSAT proof has been independently checked. The two are kept separate throughout.

\paragraph{Trust assumptions.} The exclusion assumes that the following are correct:
\begin{itemize}
\item region enumeration and reflection canonicalisation (Algorithm~1, Table~\ref{p8-tab:regions});
\item regional CNF generation (Theorem~\ref{p8-thm:exact});
\item symmetry breaking for $R_\emptyset$ (Lemma~\ref{p8-lem:sb});
\item native-split construction and the case split (Theorems~\ref{p8-thm:split} and~\ref{p8-thm:case});
\item the solvers, Kissat and Dagster's TiniSAT workers;
\item Dagster's message transport and completion accounting;
\item result merging, bookkeeping and the final inventory;
\item file integrity and provenance;
\item the independent model validator (relevant only to SAT results).
\end{itemize}
A SAT witness would be independently checkable. The UNSAT verdicts retain these assumptions
until proofs are produced.

\paragraph{What supports the verdicts.}
\begin{itemize}
\item The encoding is validated exhaustively at orders 8 and 10, and on 158 known arrays.
\item Every decomposition is proved complete.
\item Every Dagster closure passes the message audit (Section~\ref{p8-sec:gates}).
\item Every method reproduced known answers at small orders, including satisfiable regions.
\item Methods agree where both were run.
\item The result is consistent with the earlier structural theorems (Section~\ref{p8-sec:earlier}).
\end{itemize}

%======================================================================
\section{Conclusion}

This paper revisited the question whether a Costas array of order 32 can be symmetric about its
main diagonal. The question is natural because transpose symmetry is the only symmetry of the
square that a Costas array can have. The literature answers it: exhaustive enumerations published
in 1993 and 2009 find none of order 32. We set out to confirm that answer independently, and to do so with
a method that would carry over to problems that no enumeration reaches.

The approach rests on one structural observation: the fixed points of a transpose-symmetric
Costas array form a Golomb ruler. At order 32 this limits them to 0, 2, 4 or 6, and partitions the
problem into 44{,}947 reflection-canonical regions, each of which can be posed exactly as a
satisfiability problem. Most regions are easy, and a sequential solver closed them. The hard ones
needed a distributed decomposition, and there the same structure proved decisive a second time.
Splitting a region by the partners of a few coordinates works well only if no piece is left in
which two of those coordinates pair with each other, and the fixed points supply exactly the
information needed to refute such pieces in advance. Where there are no fixed points, fixing one
transposition restores the same leverage.

With these devices every region was closed, confirming independently that no transpose-symmetric,
and hence no geometrically symmetric, Costas array of order 32 exists. The confirmation is computational. It is
reproducible from the released software and audited for completeness, but it rests on the
correctness of the solvers and tooling rather than on checked proofs. It is not the cheapest route
to this particular answer: a direct search would have needed far less computation. It agrees with the published enumerations,
and with this project's earlier structural theorems, and it leaves the general order-32 question
open. Its broader lesson is methodological: in a large combinatorial search, the structure that
defines a problem's natural partition can also be made to drive the decomposition of its hardest
parts.



\begin{paperrefs}
\bibitem{p8-costas1984}
J.~P. Costas, A study of a class of detection waveforms having nearly ideal range--Doppler
ambiguity properties, \emph{Proc.\ IEEE} 72(8):996--1009, 1984.
\bibitem{p8-golombtaylor1984}
S.~W. Golomb and H.~Taylor, Constructions and properties of Costas arrays, \emph{Proc.\ IEEE}
72(9):1143--1163, 1984.
\bibitem{p8-drakakis2011}
K.~Drakakis, F.~Iorio, S.~Rickard and J.~Walsh, Results of the enumeration of Costas arrays of
order 29, \emph{Advances in Mathematics of Communications} 5(3):547--553, 2011.
\bibitem{p8-drakakis2006}
K.~Drakakis, A review of Costas arrays, \emph{Journal of Applied Mathematics}, 2006.
\bibitem{p8-taylor2009}
K.~Taylor, S.~Rickard and K.~Drakakis, On Costas arrays with various types of symmetry,
\emph{16th International Conference on Digital Signal Processing}, 2009,
doi:10.1109/ICDSP.2009.5201135.
\bibitem{p8-dgo2006}
K.~Drakakis, R.~Gow and L.~O'Carroll, On some properties of Costas arrays generated via finite
fields, \emph{40th Annual Conference on Information Sciences and Systems}, 801--805, 2006.
\bibitem{p8-brown1993}
C.~Brown, M.~Cenki, R.~Games, J.~Rushanan, O.~Moreno and P.~Pei, New enumeration results for
Costas arrays, \emph{IEEE International Symposium on Information Theory}, 405, 1993.
\bibitem{p8-dumitru2026symmetriccensus} B.~Dumitru and C.~Budala, \emph{An exhaustive census of main-diagonal symmetric Costas arrays of orders 37--42}, preprint arXiv:2608.28690, 2026. Not peer-reviewed.
\bibitem{p8-kissat}
A.~Biere, K.~Fazekas, M.~Fleury and M.~Heisinger, CaDiCaL, Kissat, Paracooba, Plingeling and
Treengeling entering the SAT Competition 2020, \emph{Proc.\ SAT Competition 2020}, 2020.
\bibitem{p8-dagster}
Dagster: distributed SAT solving over DAG decompositions, \url{https://github.com/ANU-HPC/dagster}.
\bibitem{p8-rintanen2011} J.~Rintanen, \emph{Planning with SAT, admissible heuristics and A*}, Proceedings of the 22nd International Joint Conference on Artificial Intelligence (IJCAI), 2011.
\bibitem{p8-oeisgolomb} OEIS A003022, length of the shortest Golomb ruler with $n$ marks.
\end{paperrefs}

\clearpage
\chapter{The Cost of an Unrestricted Order-32 Search: A Validated Estimate}\label{p9-paper}
\paperauthor
\paperabstract{How much computation would an exhaustive search of all order-32 permutations take? We measure
it for a straightforward column-by-column search, written afresh, which reproduces the project's exact
tree sizes at orders 11 to 14. Knuth's probe estimator and a stratified variant are validated against
exact counts at orders 12 to 16, and, on a cluster node, against 192 order-32 subtrees of up to about
$1.5\times10^9$ nodes counted exactly, where their aggregate error is $0.35\%$ or less. The order-32 tree
has $1.538\times10^{20}$ nodes, with a relative standard error under $0.1\%$ and no heavy tail. At the rate
measured on a fully loaded 104-core node, $3.14\times10^6$ nodes per core per second, this search would
take about $1.5\times10^6$ core-years, or 27 million KSU, and about 150 years on $10^4$ cores. The 2011
order-29 enumeration was at least 36 times more efficient, so Paper~1's projection of $10^3$--$10^4$
core-years is a statement about that method; scaled by the measured growth, it suggests
9{,}000--18{,}000 modern core-years under an assumed per-core speedup. Engineering the same search, with bit-parallel candidate rows ($5.0$ times faster) and pruning by all
eight symmetries of the square (a tree $2.13$ times smaller, validated exactly through order 13), lowers the
cost about eightfold, to about $1.9\times10^5$ core-years or 3.3 million KSU. The volume's own measurements
imply that a distributed SAT campaign would cost far more.}

\section{Introduction}\label{p9-sec:paperintro}
The conclusions of this volume name one route with a known path to a definite answer at order 32:
exhaustive search of all permutations. Whether that route is realistic depends on its cost, which until
now rested on a projection scaled from a published enumeration at order 29. This paper measures the cost
directly for one well-defined search, using estimators first validated where the answer is known, and
compares the result with that projection.

\paragraph{Dependencies.} This paper uses the shared notation, Paper~1's exact tree sizes as a cross-check,
and Paper~8's measured cost of SAT relative to direct search. No other paper depends on it.

\section{Preliminaries}\label{p9-sec:notation}
A \emph{Costas array} of order $n$ is an $n\times n$ permutation array in which the displacement
vectors between pairs of dots are all distinct. Writing the array as a permutation $h$ of
$\{0,\dots,n-1\}$, with the dot of column $i$ in row $h_i$, the condition is that for each
horizontal separation $m$ the vertical differences $h_{i+m}-h_i$ are distinct. Costas arrays are
known at every order up to 31; orders 32 and 33 are the smallest for which existence is open. The
dihedral group $D_4$ of the square maps Costas arrays to Costas arrays, and arrays are counted
either individually or by $D_4$ orbit (equivalence class). Further notation and terms are collected
in the notation table and glossary at the front of this volume.

\section{The search and its tree}
\label{p9-sec:tree}
The search fills columns from left to right. A node is a partial assignment $p_0,\dots,p_{c-1}$ that is
Costas-valid so far: its rows are distinct, and at every lag the vertical differences are distinct. For
each new column the search tries every unused row, checking the new dot's displacements against a
bitmask of the differences already used at each lag. The first column's row is restricted to
$p_0\le\lfloor(n-1)/2\rfloor$, which uses the vertical-complement symmetry $p\mapsto n-1-p$ to halve the
tree, as in the measurements of Paper~1. Nothing else is pruned.

The implementation, \texttt{estimate/costas\_tree.c}, was written afresh for this study and shares no code
with the project's earlier solvers. It reproduces their exact tree sizes at every order where both have
been run, differing by exactly one because Paper~1 also counts the empty root:

\begin{table}[H]
\centering
\small
\begin{tabular}{@{}rrrr@{}}
\toprule
order & nodes, this search & nodes, Paper~1 & solutions \\
\midrule
11 & 789{,}511 & 789{,}512 & 2{,}509 \\
12 & 3{,}488{,}680 & 3{,}488{,}681 & 3{,}926 \\
13 & 18{,}290{,}043 & 18{,}290{,}044 & 7{,}077 \\
14 & 83{,}176{,}084 & 83{,}176{,}085 & 8{,}626 \\
\bottomrule
\end{tabular}
\caption{Exact tree sizes. At even orders the solution counts are exactly half the published totals
(7{,}852 and 17{,}252), as the symmetry restriction requires; at odd orders arrays whose first dot is in
the middle row are kept whole, so the counts exceed half.}
\label{p9-tab:exact}
\end{table}

\section{Estimators and their validation}
\label{p9-sec:validation}
Knuth's estimator~\cite{p9-knuth1975} follows a random path from the root, choosing uniformly among the
valid children at each level. If the path meets $d_1,d_2,\dots$ children at successive levels, the
quantity $d_1+d_1d_2+d_1d_2d_3+\cdots$ is an unbiased estimate of the number of nodes, and averaging
independent probes reduces its variance. The \emph{stratified} variant enumerates every node at a fixed
depth $D$ exactly, counting the nodes above it exactly, and then probes beneath each frontier node,
which removes the variance contributed by the top of the tree. In both cases the standard error is
computed from the spread of the probes.

The estimators were validated against the exact counts at orders 12 to 16, three runs per order.
Every estimate is within $0.32\%$ of the exact count, and in all 15 runs the exact count lies within
two standard errors of the estimate.

\begin{table}[H]
\centering
\small
\begin{tabular}{@{}rrrr@{}}
\toprule
order & Knuth, $10^6$ probes & stratified, $D=3$ & stratified, $D=4$ \\
\midrule
12 & $-0.02\%$ ($\pm0.16\%$) & $+0.27\%$ ($\pm0.44\%$) & $-0.04\%$ ($\pm0.29\%$) \\
13 & $+0.03\%$ ($\pm0.17\%$) & $-0.13\%$ ($\pm0.40\%$) & $+0.04\%$ ($\pm0.26\%$) \\
14 & $-0.12\%$ ($\pm0.19\%$) & $+0.32\%$ ($\pm0.40\%$) & $+0.13\%$ ($\pm0.24\%$) \\
15 & $-0.15\%$ ($\pm0.20\%$) & $-0.09\%$ ($\pm0.37\%$) & $-0.20\%$ ($\pm0.22\%$) \\
16 & $-0.08\%$ ($\pm0.21\%$) & $-0.10\%$ ($\pm0.37\%$) & $+0.10\%$ ($\pm0.21\%$) \\
\bottomrule
\end{tabular}
\caption{Error of each estimate against the exact count, with two standard errors in parentheses. The
exact counts at orders 15 and 16 are those of Paper~1.}
\label{p9-tab:validation}
\end{table}

Kilby, Slaney, Thi\'ebaux and Walsh describe weighted-backtrack and recursive estimators that improve
on Knuth's for uneven trees~\cite{p9-kilby2006}. They were not implemented here. The Costas tree turns out
to be regular enough that plain probing is accurate and well calibrated, both at these orders and,
as Section~\ref{p9-sec:cluster} shows, in genuine order-32 subtrees.

\section{The order-32 tree}
\label{p9-sec:order32}
Five independent estimates at order 32 agree closely:

\begin{table}[H]
\centering
\small
\begin{tabular}{@{}lrr@{}}
\toprule
method & estimate & standard error \\
\midrule
Knuth, $2\times10^6$ probes, seed 21 & $1.5365\times10^{20}$ & $0.17\%$ \\
Knuth, $2\times10^6$ probes, seed 22 & $1.5409\times10^{20}$ & $0.17\%$ \\
Knuth, $2\times10^6$ probes, seed 23 & $1.5438\times10^{20}$ & $0.17\%$ \\
stratified, $D=3$, 14{,}640 frontier nodes & $1.5374\times10^{20}$ & $0.19\%$ \\
stratified, $D=4$, 408{,}470 frontier nodes & $1.5354\times10^{20}$ & $0.12\%$ \\
\bottomrule
\end{tabular}
\caption{Estimates of the order-32 tree. Their inverse-variance mean is $1.538\times10^{20}$ nodes, with a
relative standard error of $0.07\%$.}
\label{p9-tab:order32}
\end{table}

\noindent There is no sign of a heavy tail: no single probe contributes more than $0.02\%$ of any total.
At orders 29, 30 and 31 the estimates are $1.566\times10^{18}$, $7.097\times10^{18}$ and
$3.406\times10^{19}$, so the tree grows by factors of $4.53$, $4.80$ and $4.52$ per order from 29 to 32,
and by $98.2$ in all. The alternation comes from the symmetry restriction, which keeps the middle row
whole at odd orders. The growth agrees with Paper~1's measured ratio of $4.80$ per order at orders 11 to
16, which extrapolated from order 16 gives $1.6\times10^{20}$.

\section{Measurements on a cluster node}
\label{p9-sec:cluster}
The same program was run on one node of the cluster used for the campaign of Paper~8: two Intel Xeon
Sapphire Rapids processors, 104 cores in all, with the \texttt{powersave} frequency governor and turbo
enabled. The job cost 78 service units.

\begin{table}[H]
\centering
\small
\begin{tabular}{@{}lr@{}}
\toprule
measurement at order 32 & value \\
\midrule
node rate, one core alone & $3.72\times10^6$ nodes/s \\
node rate per core, all 104 cores busy & $3.14\times10^6$ nodes/s (range $3.13$--$3.16$) \\
node rate per node & $3.27\times10^8$ nodes/s \\
slowdown of a fully loaded node & $1.18$ \\
tree size, 8 independent Knuth runs & $1.5381\times10^{20}$ (relative s.e.\ $0.06\%$) \\
\bottomrule
\end{tabular}
\caption{Cluster measurements. The single-core rate equals the virtual machine's, and the tree estimate
agrees with Table~\ref{p9-tab:order32} to within $0.01\%$.}
\label{p9-tab:cluster}
\end{table}

The validation of Section~\ref{p9-sec:validation} used orders 12 to 16. The cluster run also tested the
estimator on genuine order-32 subtrees. Each of 192 cores took a random path to depth 8 or 9 of the
order-32 tree, estimated the subtree below with 5{,}000 probes, and then counted it exactly.

\begin{table}[H]
\centering
\small
\begin{tabular}{@{}lrr@{}}
\toprule
& depth 8 & depth 9 \\
\midrule
subtrees counted exactly & 96 of 96 & 96 of 96 \\
median subtree size & $1.51\times10^9$ nodes & $7.69\times10^7$ nodes \\
exact count within two standard errors & $93.8\%$ & $94.8\%$ \\
total of estimates / total of exact counts & $0.9965$ & $1.0002$ \\
\bottomrule
\end{tabular}
\caption{The estimator on order-32 subtrees. No subtree reached the 30-minute time cap, so none was
excluded.}
\label{p9-tab:deepval}
\end{table}

\noindent The coverage is close to the nominal $95\%$, and across the depth-8 subtrees, about
$1.5\times10^{11}$ nodes in all, the estimates' total is within $0.35\%$ of the exact one. The estimator is
therefore validated in the tree it is used on, not only at small orders.

\section{From nodes to cost}
\label{p9-sec:cost}
A campaign keeps every core busy, so the relevant rate is the fully loaded one, $3.14\times10^6$ nodes per
core per second. At that rate the order-32 tree takes
\[
\frac{1.538\times10^{20}}{3.14\times10^6}\ \text{s}\approx1.36\times10^{10}\ \text{core-hours}
\approx1.5\times10^{6}\ \text{core-years}.
\]
On a queue charging 2 service units per core-hour, that is about $2.7\times10^{10}$ service units, or
27 million thousand-unit blocks (KSU). Exploiting the full symmetry group of the square, rather than the
single symmetry used here, would save at most a further factor of about four. Neither figure is within
reach of any allocation.

\section{Latency}
\label{p9-sec:latency}
How long a campaign takes depends on how evenly its work divides. The cluster run estimated the subtree
below every one of the 408{,}470 nodes at depth 4, each with 20 probes, and scheduled them as tasks,
largest first, onto the least-loaded worker. The work is fairly even: the largest task is $34.7$ times
the mean, $1.15\times10^6$ core-hours, and the largest $1\%$ of tasks hold $3.8\%$ of the work.

\begin{table}[H]
\centering
\small
\begin{tabular}{@{}rrrr@{}}
\toprule
workers & wall time & efficiency & charged KSU \\
\midrule
104 & about 14{,}900 years & $100.0\%$ & $2.71\times10^7$ \\
1{,}040 & about 1{,}490 years & $100.0\%$ & $2.71\times10^7$ \\
10{,}400 & about 149 years & $99.7\%$ & $2.72\times10^7$ \\
104{,}000 & about 131 years & $11.3\%$ & $2.39\times10^8$ \\
\bottomrule
\end{tabular}
\caption{Wall time and charged cost of the straightforward search, with depth-4 tasks.}
\label{p9-tab:latency}
\end{table}

\noindent Up to about $10^4$ workers, depth-4 tasks keep every core busy. Beyond that the largest task
exceeds the ideal wall time, and efficiency collapses unless the work is split more finely. The wall
times put the straightforward search beyond reach at any scale. A method costing about $10^4$ core-years,
as Section~\ref{p9-sec:legacy} suggests for the 2011 method, would need about a year on $10^4$ cores,
provided its work divides as evenly.

\section{The 2011 enumeration was far more efficient}
\label{p9-sec:legacy}
Paper~1 projected the order-32 search at $10^3$--$10^4$ modern core-years by scaling the reported cost of
the 2011 order-29 enumeration, 366.55 CPU-years~\cite{p9-drakakis29}. The present search would need about
13{,}400 core-years at order 29, so the 2011 method was at least 36 times more efficient, and more once
the slower hardware of 2010 is allowed for. Paper~1's projection is therefore a statement about that
method, not about a straightforward column-by-column search. Section~\ref{p9-sec:engineering} shows that
most of the gap is recoverable by engineering.

Scaling the 2011 cost by the measured growth from order 29 to 32, a factor of 98 rather than the $5^3=125$
of a five-fold-per-order rule, gives about 36{,}000 CPU-years on 2010 hardware. If per-core speed has
improved two- to four-fold since then, an assumption that has not been measured, that is
9{,}000--18{,}000 modern core-years, or about 160{,}000--320{,}000 KSU on a 2-unit queue, and up to about 370{,}000 allowing for the
slowdown of a fully loaded node measured in Section~\ref{p9-sec:cluster}. This assumes that
the 2011 method's tree grows like the present one; its algorithm has not been reproduced here.

\section{Engineering the search}
\label{p9-sec:engineering}
Most of the 2011 method's advantage turns out to be recoverable by engineering the same search. Three
changes were tested, each first validated against exact counts.

\paragraph{Bit-parallel candidate rows.} The search of Section~\ref{p9-sec:tree} tests each candidate row
against every earlier column, $O(n\cdot c)$ work to expand a node at column $c$. Row $r$ conflicts with
column $i$ exactly when the difference $r-p_i$ is already used at lag $c-i$, so the rows forbidden by
column $i$ are the used-difference mask of that lag, shifted into row coordinates. ORing these masks gives
every forbidden row at once, in $O(c)$ word operations per node. The tree is unchanged: node and
solution counts are identical at orders 10 to 14. At order 32 the rate rises from $3.33\times10^6$ to
$1.68\times10^7$ nodes per second on the same core, a factor of $5.0$.

\paragraph{Column order.} Four orders were compared, all with the same symmetry restriction, which applies
to whichever column is placed first.

\begin{table}[H]
\centering
\small
\begin{tabular}{@{}rrrrr@{}}
\toprule
order & left to right & outside in & middle out & fewest legal rows \\
\midrule
12 & $3.49\times10^6$ ($0.36$\,s) & $1.22\times10^7$ ($1.39$\,s) & $3.49\times10^6$ ($0.44$\,s) & $2.42\times10^6$ ($1.22$\,s) \\
13 & $1.83\times10^7$ ($2.19$\,s) & $7.81\times10^7$ ($8.86$\,s) & $1.79\times10^7$ ($1.95$\,s) & $1.22\times10^7$ ($7.58$\,s) \\
14 & $8.32\times10^7$ ($12.5$\,s) & $4.38\times10^8$ ($55.9$\,s) & $8.32\times10^7$ ($11.5$\,s) & $5.38\times10^7$ ($35.6$\,s) \\
\bottomrule
\end{tabular}
\caption{Nodes and time for four column orders. ``Fewest legal rows'' chooses at each node the unplaced
column with fewest legal rows and abandons a branch when any column has none.}
\label{p9-tab:orders}
\end{table}

\noindent Any order that keeps the placed columns contiguous gives the same tree, because whether a block
of dots is Costas does not depend on where the block sits; at even orders the counts for left to right and
middle out are identical. Placing dots far apart, outside in, wastes the short-lag constraints, which prune
hardest, and gives a tree five times larger by order 14. Choosing the column with fewest legal rows shrinks
the tree increasingly, to $65\%$ of the baseline at order 14, and its growth per order is about $4.65$
against $4.84$. But each node costs about three times as much, so it is slower at every order measured.

Orders that place a column between two earlier ones need one more check than left to right. A new dot
forms two pairs at the same lag $L$, one on each side, and these collide exactly when its row is the
midpoint of the other two, $2r=p_{c-L}+p_{c+L}$: the Ladendorf condition of the 2009 search. A first
version of these searches omitted the check and accepted non-Costas arrays; the validation against exact
array counts caught it before any result was used.

\paragraph{Full symmetry pruning.} Using all eight symmetries of the square, the search keeps exactly one
array per class: the lexicographically smallest of its eight images. A prefix is cut when some image is
guaranteed to be smaller for every completion. An image entry is either known, or confined to an interval:
an unplaced row's column lies among the unplaced columns. The vertical flip compares known values; the four
symmetries involving transposition can often be decided early from the partly known inverse; the horizontal
flip and the half-turn are decided only near the leaves. The pruning is exact: at every order from 4 to
13 the search outputs precisely the class representatives found by enumerating every Costas array and
grouping them by symmetry, including 990 classes among 7{,}852 arrays at order 12 and 1{,}616 among 12{,}828
at order 13. On the pruned tree the estimator is within $0.3\%$ of the exact counts at orders 13 to 15.

\begin{table}[H]
\centering
\small
\begin{tabular}{@{}rrrr@{}}
\toprule
order & one symmetry & all eight & fewer nodes \\
\midrule
12 & $3{,}488{,}680$ & $1{,}544{,}404$ & $2.26\times$ \\
13 & $18{,}290{,}043$ & $7{,}509{,}524$ & $2.44\times$ \\
14 & $83{,}176{,}084$ & $37{,}132{,}082$ & $2.24\times$ \\
15 & $435{,}090{,}280$ & $182{,}409{,}077$ & $2.39\times$ \\
\bottomrule
\end{tabular}
\caption{Nodes with one and with all eight symmetries, in the bit-parallel search.}
\label{p9-tab:symmetry}
\end{table}

\noindent At order 32 three independent estimates put the pruned tree at $7.22\times10^{19}$ nodes, $2.13$
times smaller, while the extra comparisons lower the rate to $1.42\times10^7$ nodes per second. The search is
about $1.8$ times faster overall: less than the theoretical four, because two of the symmetries can be
decided only near the leaves.

\paragraph{The revised cost.} With the fully loaded slowdown of Section~\ref{p9-sec:cluster}:

\begin{table}[H]
\centering
\small
\begin{tabular}{@{}lrrr@{}}
\toprule
search & core-hours & core-years & KSU \\
\midrule
as in Section~\ref{p9-sec:cost} & $1.36\times10^{10}$ & $1.55\times10^6$ & $2.7\times10^7$ \\
bit-parallel candidate rows & $3.0\times10^{9}$ & $3.4\times10^5$ & $6.0\times10^6$ \\
and full symmetry pruning & $1.7\times10^{9}$ & $1.9\times10^5$ & $3.3\times10^6$ \\
\bottomrule
\end{tabular}
\caption{The cost of the order-32 search, before and after engineering.}
\label{p9-tab:revised}
\end{table}

\noindent The engineered search is about eight times cheaper, but still about 18 years on $10^4$ cores. At
order 29 it would need about 1{,}650 core-years against the 2011 enumeration's 366.55 CPU-years, so a factor
of four to five remains, more if 2010 hardware was slower per core. Whether that factor lies in node rate or
in a smaller tree is not known. For comparison, a 2026 preprint estimates by sampling that an
unrestricted order-28 search with its software would take 43.7 CPU-years, and projects 7.97 GPU-years
at order 30~\cite{p9-dumitru2026localpatterns}. Reductions in the growth per order compound: the fewest-legal-rows order
lowers it from $4.84$ to about $4.65$, which over the orders up to 32 is worth about three times, if its
per-node cost could be brought down to the baseline's. Paper~10 measures a related attempt: binary
displacement consistency, which maintains the domains that order needs, reduced the forward-checking tree
about fourfold at order 14 but made complete searches slower, and was not retained.

\section{Implications for a distributed SAT campaign}
\label{p9-sec:sat}
A planning model for a distributed SAT campaign can be written
\[
K_{\mathrm{SAT}}=K_{\mathrm{ref}}\,\frac{\mu}{S},
\]
where $K_{\mathrm{ref}}$ is the cost of a reference enumeration, $S$ the improvement of a modern direct
method over it, and $\mu$ the overhead of SAT and its decomposition over that direct method. Both
factors can now be bounded from measurements in this volume.
\begin{itemize}
\item Relative to the 2011 method, the straightforward modern search is not faster but at least 36
times slower, so $S>1$ requires an algorithm better than the 2011 one, which no one in this
investigation has shown.
\item For the transpose-symmetric subclass, the SAT campaign of Paper~8 used about 490 times the
computation of the direct search later run to completion, so $\mu$ was of that order, not a small constant. Even $\mu=5$ applied to
the 2011 method's scaled cost gives about 0.8--1.6 million KSU.
\end{itemize}
On cost, therefore, a distributed SAT campaign is not a viable route to the unrestricted problem. The
only candidate is a direct enumeration at least as efficient as the 2011 method, distributed through a
fail-closed task queue, at a cost of order $10^5$ KSU on present assumptions. Before any allocation is
requested, that method should be reproduced and its rate and tree measured at orders where they can be
checked, such as 20 to 24; the estimator validated here makes that calibration cheap.

\section{Limitations}
\label{p9-sec:limits}
The main estimates describe one search, with columns filled from left to right and one symmetry used;
Section~\ref{p9-sec:engineering} measures three changes to it. Other
filling orders, such as from the middle outwards, give different trees, and the 2011 method's tree is
not measured here. The node rates were measured on one virtual machine and one cluster node. The weighted-backtrack and recursive estimators were not implemented. The improvement in
per-core speed since 2010 is an assumption.


\section{Status of the claims}\label{p9-sec:status}
The tree sizes at orders 11 to 14 are exact and agree with Paper~1. The order-32 tree size is an
estimate, validated in all 15 runs at orders 12 to 16 and consistent across five independent runs at
order 32. The node rates are measurements on one virtual machine and on one 104-core cluster node, whose results are recorded in \texttt{estimate/cluster\_results/}; the order-32 validation was made on the cluster node. The comparison with the 2011 enumeration uses its
reported cost, and the conversion to modern core-years assumes a per-core speedup that has not been
measured. The study is reproduced by \texttt{estimate/run\_estimates.sh}, and its figures are recorded in
\texttt{estimate/estimate\_results.json}.

\paragraph{Relation to prior work.} Knuth's estimator is classical~\cite{p9-knuth1975}, and better estimators
for uneven trees exist~\cite{p9-kilby2006}. The measurements and the comparison with the 2011 enumeration are
this project's. The engineering measurements are reproduced by the programs in \texttt{ordering/}; the symmetry pruning is validated exactly through order 13.
\begin{paperrefs}
\bibitem{p9-knuth1975} D.~E.~Knuth, \emph{Estimating the efficiency of backtrack programs}, Mathematics of Computation 29 (1975), 121--136.
\bibitem{p9-kilby2006} P.~Kilby, J.~Slaney, S.~Thi\'ebaux and T.~Walsh, \emph{Estimating search tree size}, Proceedings of the 21st National Conference on Artificial Intelligence (AAAI), 2006.
\bibitem{p9-drakakis29} K.~Drakakis, F.~Iorio, S.~Rickard and J.~Walsh,
\emph{Results of the enumeration of Costas arrays of order 29}, IEEE Trans.\
Inform.\ Theory \textbf{57} (2011).
\bibitem{p9-dumitru2026localpatterns} B.~Dumitru and C.~Budala, \emph{Local patterns and the cost of Costas array search}, preprint hal-05731061v1, 2026, \url{https://hal.science/hal-05731061}. Not peer-reviewed.
\end{paperrefs}

\clearpage
\chapter{Searching for Short Obstructions to Partial Costas Arrays}\label{p10-paper}
\paperauthor
\paperabstract{Almost every prefix that an exhaustive Costas search explores at mid-depth passes every local
test yet has no completion. We ask whether those deaths have reasons shorter than the search, testing each
candidate against gates fixed in advance and against negative controls. Minimal infeasible dot sets are
mostly explained by the sparsity of Costas arrays, as an equally sparse random set shows; their residue is
corner structure, which corner reduction predicts exactly across orders. Clause learning, Hall's condition
and cutting planes refute dead prefixes only with effort comparable to the native search. The
pair-marginal linear relaxation, the first level of the Sherali--Adams hierarchy with displacement
capacities, is different: it is sound, and beyond a sharp transition near $0.51$--$0.56$ of the order it
certifies every sampled dead prefix, at orders 12 to 16, as an independent implementation confirms. But at
order 32 the transition falls where the search tree is widest, capping the benefit of any rule acting there
at about nine-fold. The next level of the hierarchy is stronger but moves the transition by only one column. Binary displacement consistency and chain-filtered Hall analysis, which also work on exact occupancy, become strong in the same band, which stays near $0.55n$ through order 20, beyond the peak of the Costas counts, while aggregate statistics of the differences carry little beyond sparsity.}

\section{Introduction}\label{p10-sec:paperintro}
The previous papers close structured routes to an order-32 Costas array and measure the cost of
the search that remains. This paper asks whether that search could be shortened by finding \emph{why} its
dead prefixes die: whether there are obstructions shorter than the subtrees they would replace. It reports
a sequence of experiments, from local rules through clause learning and cutting planes to a hierarchy of
linear relaxations, together with the gates, controls and independent replication applied to each.

\paragraph{Dependencies.} This paper uses the catalogues and search of Paper~9, its estimator, and the corner
surgery of Paper~2 as a positive control.

\section{Preliminaries}\label{p10-sec:notation}
A \emph{Costas array} of order $n$ is an $n\times n$ permutation array in which the displacement
vectors between pairs of dots are all distinct. Writing the array as a permutation $h$ of
$\{0,\dots,n-1\}$, with the dot of column $i$ in row $h_i$, the condition is that for each
horizontal separation $m$ the vertical differences $h_{i+m}-h_i$ are distinct. Costas arrays are
known at every order up to 31; orders 32 and 33 are the smallest for which existence is open. The
dihedral group $D_4$ of the square maps Costas arrays to Costas arrays, and arrays are counted
either individually or by $D_4$ orbit (equivalence class). Further notation and terms are collected
in the notation table and glossary at the front of this volume.

\section{The question}
\label{p10-sec:question}
An exhaustive search at order 32 spends almost all its effort in prefixes that pass every local test yet have
no completion. At mid-depth this is nearly universal: of the prefixes that survive the native checks and
forward checking, 90--100\% are dead (Section~\ref{p10-sec:local}). The question of this paper is whether those
deaths have \emph{reasons shorter than the search that discovers them}: rules, certificates or relaxations
that refute a dead prefix with far less work than exhausting its subtree, and that recur across positions
and orders. Each candidate was tested against a gate fixed in advance, against negative controls, and,
for the strongest results, by an independent implementation.

\section{Method}
\label{p10-sec:method}
\paragraph{Ground truth.} At orders 12 to 15 every Costas array is listed by the project's exact search, so a
prefix is completable exactly when some array contains it. At larger orders completability is decided by
a SAT solver and, for subtree sizes, by the native search; no relaxation is ever used as ground truth.

\paragraph{Prefixes.} A prefix fills columns from the left. It is \emph{locally admissible} if its rows are
distinct, no displacement repeats, and every unplaced column still has a legal row (forward checking).

\paragraph{Sampling.} Prefixes are drawn by random walks through the locally admissible tree. A walk reaches
a prefix with probability $\prod_i 1/b_i$, where $b_i$ is the number of legal choices at step $i$, so weighting
by $\prod_i b_i$ gives estimates for the uniform depth layer; weighting dead prefixes by their native subtree
size gives estimates of the share of \emph{work} affected. Weighted estimates are reported with their
effective sample sizes $N_{\mathrm{eff}}=(\sum w)^2/\sum w^2$.

\section{Local obstructions}
\label{p10-sec:local}
\paragraph{Minimal obstructions.} For dead order-12 prefixes at depth 5, deleting dots one at a time while the
set stays infeasible gives minimal infeasible dot sets: sets of dots with no repeated displacement that lie
in no Costas array. They are tiny, almost always four dots, but scattered: 60 obstructions came in 56
shapes, and none of twelve shapes tested is infeasible at every position at order 12 or 13.

\paragraph{Negative control.} The same pipeline was run against 7{,}852 random permutations, the number of
order-12 Costas arrays, on the same prefixes:

\begin{table}[H]
\centering
\small
\begin{tabular}{@{}lrrrr@{}}
\toprule
depth 5 & dead & 3-dot obstructions & 4-dot & recurring shapes \\
\midrule
Costas arrays & $86\%$ & 19 & 319 & $29\%$ \\
random permutations & $91\%$ & 6 & 352 & $22\%$ \\
\bottomrule
\end{tabular}
\caption{Minimal obstructions against the Costas arrays and against an equally sparse random set, on the same
400 sampled prefixes.}
\label{p10-tab:control}
\end{table}

\noindent The bulk of the obstruction statistics is reproduced by sparsity alone. The Costas-specific residue
is small: more three-dot obstructions, and somewhat more recurrence.

\paragraph{The residue is positional.} Exhaustively at order 12, 3{,}508 of the 288{,}600 three-dot sets with no
repeated displacement lie in no Costas array: $1.2\%$, against about $0.27\%$ expected for a random set of the
same size. The excess concentrates at corners: $0.84\%$ of triples with no corner dot are infeasible, $4.9\%$
with one, and $62\%$ with two. No three-dot shape is infeasible at every position at both orders 12 and 13.

\paragraph{A positive control: corner reduction.} Deleting a corner dot translates the other dots, so order-$n$
arrays with a dot at $(0,0)$ correspond exactly to order-$(n-1)$ arrays that admit a corner addition. This
predicts the feasibility of every corner triple at order $n$ from the catalogue at order $n-1$:

\begin{table}[H]
\centering
\small
\begin{tabular}{@{}rrrr@{}}
\toprule
order & arrays with a dot at $(0,0)$ & admissible arrays, order $n-1$ & triples predicted correctly \\
\midrule
13 & 509 & 509 & 8{,}676 of 8{,}676 \\
14 & 734 & 734 & 12{,}132 of 12{,}132 \\
15 & 616 & 616 & 16{,}513 of 16{,}513 \\
\bottomrule
\end{tabular}
\caption{Corner reduction as a lifted obstruction: exact at three orders, with no refitting.}
\label{p10-tab:corner}
\end{table}

\noindent Plain infeasibility one order lower explains none of the 1{,}599 infeasible corner triples; the whole
effect is the small size of the corner-admissible subclass. The apparatus recovers real cross-order
structure, and the structure it recovers is the corner surgery of Paper~2.

\section{Clause learning, matching and cutting planes}
\label{p10-sec:proofs}
\paragraph{Incremental clause learning.} Solving 36 order-16 sibling cubes with one persistent CDCL solver
per family, instead of a fresh solver per cube, saved $8$--$18\%$ of conflicts; related siblings benefited
slightly more than unrelated cubes.

\paragraph{The native/SAT frontier.} For dead prefixes that survive all local checks, the native search's
subtree size was compared with the conflicts a fresh CDCL solver needs to refute the prefix. At order 14 the
median ratio was $1.4$--$1.6$ native nodes per conflict, at order 16 $1.16$, and no prefix showed even a
five-fold advantage: CDCL finds no shorter reasons than the search.

\paragraph{Hall's condition.} A matching test of unplaced columns against unused rows catches $0$--$2\%$ of dead
prefixes at mid-depth, rising to $62\%$ only at depth 9 of 12, where little work remains.

\paragraph{Cutting planes.} Dead prefixes below the LP transition of Section~\ref{p10-sec:lp} were refuted by Exact,
a cutting-planes pseudo-Boolean solver descended from RoundingSat~\cite{p10-elffers2018}, on the native
integer model of that section. Every prefix was refuted and every live control satisfied, but the effort
tracked the native subtree: medians of $1.3$, $1.2$ and $0.9$ native nodes per conflict at orders 12, 13 and
14, falling as subtrees grow. The build used had no LP integration, and its proofs were not logged.

\section{The pair-marginal relaxation}
\label{p10-sec:lp}
\paragraph{Formulation.} Placement variables $x_{i,r}\in[0,1]$ satisfy the row and column equalities. For each
pair of columns $i<j$, variables $y_{ijrs}\in[0,1]$ over distinct rows carry one unit of mass, with marginals
$\sum_s y_{ijrs}=x_{i,r}$ and $\sum_r y_{ijrs}=x_{j,s}$: the first level of the Sherali--Adams
hierarchy~\cite{p10-sherali1990}. Each displacement class $(d,e)$ has capacity $\sum y\le1$ over the pairs at
lag $d$ and rows differing by $e$. A prefix fixes its placements to 1. Every constraint holds for every
Costas array, so the relaxation is sound at every order. A weaker form replaces the marginals by the
McCormick inequalities $y\le x_{i,r}$, $y\le x_{j,s}$, $y\ge x_{i,r}+x_{j,s}-1$.

\paragraph{Soundness, checked.} No prefix of a real Costas array was cut (60 of 60, at depths 3 to 10 of order
12), nor any of 483 completable prefixes in the confirmation below.

\paragraph{A sharp transition.} Across orders 12 to 16, with 60 to 160 dead prefixes per cell:

\begin{table}[H]
\centering
\small
\begin{tabular}{@{}rrrrrr@{}}
\toprule
order & depth & $k/n$ & weak form & strong form & strong, uniform layer \\
\midrule
12 & 6 & $0.500$ & $5\%$ & $28\%$ & $25\%$ \\
12 & 7 & $0.583$ & $53\%$ & $100\%$ & $100\%$ \\
13 & 7 & $0.538$ & $14\%$ & $74\%$ & $72\%$ \\
13 & 8 & $0.615$ & $75\%$ & $100\%$ & $100\%$ \\
14 & 8 & $0.571$ & $41\%$ & $98\%$ & $99\%$ \\
15 & 8 & $0.533$ & $8\%$ & $63\%$ & $68\%$ \\
15 & 9 & $0.600$ & $57\%$ & $100\%$ & $100\%$ \\
16 & 9 & $0.562$ & $30\%$ & $93\%$ & $94\%$ \\
16 & 10 & $0.625$ & $83\%$ & $100\%$ & $100\%$ \\
\bottomrule
\end{tabular}
\caption{Share of dead prefixes the relaxation proves infeasible, near the transition. Work-weighted rates
agree beyond the transition and are slightly lower inside it.}
\label{p10-tab:transition}
\end{table}

\noindent Logistic fits put the 90\% point $\alpha_{90}$ between $0.512$ and $0.557$ of the order at all five
orders, and a fixed-fraction model fits far better than a fixed number of remaining columns (negative
log-likelihood 314 against 606). Beyond the transition no miss occurred among 1{,}140 sampled dead
prefixes, an approximate 95\% upper bound of $0.26\%$ on the miss rate, with the weighted estimates agreeing
at effective sample sizes of 32 to 146 per cell. An independent implementation, with its own enumerator,
generator and sampler, reproduced the transition and the weak--strong contrast at orders 12 to 14.

\paragraph{The certificates are global.} Shrinking the displacement capacities to a minimal set of lags that
still refutes each prefix leaves 3 to 5 lags, but the sets vary: 17 distinct among 20 prefixes at order 12.
A per-lag matching test, Hall's condition on column pairs and differences, explains only 5 of 25 and 1 of 16
of the certified prefixes. The refutations couple several lags through shared columns.

\section{Order-32 geometry}
\label{p10-sec:geometry}
The validated estimator of Paper~9, applied per depth to the order-32 tree with all eight symmetries,
locates the transition at depths 17--18 and the tree's widest level at depth 20. Calibrated against exact
counts at orders 12 to 14, the per-depth estimates were within $0.3\%$.

\begin{table}[H]
\centering
\small
\begin{tabular}{@{}rrrr@{}}
\toprule
cut depth & locally admissible prefixes & share of the tree at or above it & best possible saving \\
\midrule
16 & $1.91\times10^{18}$ & $3.6\%$ & $28\times$ \\
17 & $5.00\times10^{18}$ & $11.4\%$ & $8.7\times$ \\
18 & $8.20\times10^{18}$ & $27.9\%$ & $3.6\times$ \\
19 & $7.41\times10^{18}$ & $52.0\%$ & $1.9\times$ \\
\bottomrule
\end{tabular}
\caption{The order-32 frontier, and the most any rule complete at a given depth could save if it cost
nothing.}
\label{p10-tab:frontier}
\end{table}

\noindent Solving the relaxation once per frontier prefix is out of the question. More fundamentally, even a
free, perfect replica of its verdict at depth 17 saves at most a factor of about nine. The depth of the transition was measured at orders 12 to 16, near the
peak of the Costas counts; that it stays near the same fraction of the order as the counts fall, which
they do exponentially fast in proportion to all permutations~\cite{p10-warnke2023density}, was checked at orders
18 and 20, beyond the peak, for binary displacement consistency (Section~\ref{p10-sec:late}): there too it
refutes almost nothing at half depth and about three quarters of prefixes at $0.55n$. That the fraction
stays put through the far sparser orders 24 to 32 remains an assumption.

\section{Triple consistency}
\label{p10-sec:triple}
The next level of the hierarchy adds, for every triple of columns, a distribution over distinct row triples
whose three projections equal the three pair distributions. An independent order-12 implementation, on
exactly equivalent prefix-reduced models, found it strictly stronger: it certified all 12 dead prefixes
sampled at depth 6, where the pair relaxation certifies about half. At depth 5 ($0.42n$) it certified 2 of 22,
far below the predeclared gate of 90\%. At depth 4 no instance was certified within the time limit. No live
control was cut. The transition moves by about one column, from near $0.58n$ to near $0.50n$ at order 12; at
order 32 that corresponds to depth 16, where the ceiling of Table~\ref{p10-tab:frontier} is $28\times$, for a
relaxation that could not be run at that order.

\section{Late mechanisms}
\label{p10-sec:late}
Three further mechanisms that work on exact displacement occupancy were measured, one by the project's
search and two by an independent implementation within the project, whose code and data are kept in
\texttt{mining/independent/}.

\paragraph{Binary displacement consistency.} At every node, each pair of unplaced columns must admit a pair
of rows that repeats no used displacement, either between the two new dots or against the placed dots,
and rows without such a partner are removed until nothing changes. On complete searches it is sound (every
order returns exactly the known arrays) and strong:

\begin{table}[H]
\centering
\small
\begin{tabular}{@{}rrrrr@{}}
\toprule
order & no lookahead & forward checking & binary consistency & reduction \\
\midrule
10 & 152{,}100 & 73{,}778 & 19{,}715 & $7.7\times$ \\
12 & 3{,}488{,}680 & 1{,}583{,}032 & 284{,}922 & $12.2\times$ \\
13 & 18{,}290{,}043 & 8{,}116{,}650 & 1{,}362{,}891 & $13.4\times$ \\
14 & 83{,}176{,}084 & 36{,}110{,}650 & 5{,}314{,}892 & $15.6\times$ \\
15 & 435{,}090{,}280 & 185{,}553{,}661 & 26{,}624{,}065 & $16.3\times$ \\
\bottomrule
\end{tabular}
\caption{Nodes of the left-to-right search with three levels of lookahead.}
\label{p10-tab:binac}
\end{table}

\noindent The reduction grows with the order, but more slowly at each step, which is consistent with a large
constant factor rather than a lower growth rate. Where it acts settles the question: at orders 13 and 14,
93\% and 97\% of its cuts fall between $0.54n$ and $0.64n$, almost none before $0.5n$, and the smallest
remaining domain at a cut is typically one or two rows. Binary consistency is therefore strong in the
same band as the pair-marginal relaxation, where the tree is widest, and not earlier. Recomputed from
scratch at every node, as here, it is about eight times slower than the plain search at order 15.

Whether it can pay for itself was then tested against a stronger baseline, the search with forward
checking, using one engine and three policies whose solution sets were checked to be identical. Against
forward checking, binary consistency saves only about four times the nodes at order 14, and each call
costs about forty forward-checking nodes. Triggering it only where it acts, from about $0.6n$ or when a
domain is down to one row, still left the search 1.3 to 1.5 times slower at orders 13 and 14. More
decisively, the time spent outside the consistency checks bounds what any implementation could achieve:
even checks that cost nothing would make the search at most 2.0 times faster at order 13 and 1.6 times at
order 14. That ceiling falls with the order. The tested implementation therefore fails a predeclared
production gate of a 1.25-fold speed-up, and an incremental one would have to be about four times cheaper
merely to reach that gate at order 14, against a shrinking margin; binary consistency was not retained in
the production solver. The comparison used only the vertical-flip symmetry, identically for all three
policies, and was run at orders 12 to 14. The two left-to-right policies did not commit forced columns out
of order, whereas the policy with dynamic column choice does so implicitly by always choosing the column
with fewest rows; that asymmetry favours dynamic choice, which nevertheless did not improve elapsed time.
Binary consistency thus remains informative about where pruning power lies, but not, on this evidence, a
route to a faster search.

The same transition holds beyond the peak of the Costas counts. The independent implementation applied
binary consistency to 24 random prefixes per depth at orders 18 and 20, which have 15{,}096 and 6{,}464
arrays. At half depth it refuted 1 prefix of 24 at each order; at depth 10 of 18 it refuted 18 of 24, and at
depth 11 of 20 it refuted 16 of the 21 that forward checking had not already closed. In both cases the
smallest remaining domain at that depth was typically two rows. Among prefixes that survive binary
consistency near half depth, the pair-marginal relaxation still proved some infeasible (2 of 7 at order 18
depth 9, 1 of 8 at order 20 depth 10), so on some prefixes it is strictly the stronger of the two.

\paragraph{Chain-filtered Hall analysis.} The independent implementation kept, for every pair of columns,
only the row pairs supported through every intermediate column, and then tested Hall's condition on each
lag's positions and available differences. Hall deficiencies appear in none of 20 dead prefixes at order 14
depth 6, in 4 at depth 7 and in 19 at depth 8 ($0.57n$); at order 13 depth 7 they appear in 10 of 20,
where the pair-marginal relaxation certifies about three quarters. Splitting on sets with one unit of slack
increased total work for every one of the 49 prefixes tested.

\paragraph{Where completion is decided.} The same implementation counted, at orders 12 to 14, which locally
valid prefixes extend to a Costas array. At depth 3, 96--98\% do; at the midpoint depth $\lceil0.53n\rceil$,
0.8\%, 0.5\% and 0.09\% do. No sound rule can prune much at depth 3 of these orders. The work between the two
depths is nearly uniform across depth-3 prefixes (Gini coefficient about 0.09), and only 30--38\% of it lies
above the midpoint level, as a branching factor near four implies, so compiling that stretch of the search
could save at most a factor of about 1.3. These orders sit near the peak of the Costas counts, which caps
what can be learned from them about early pruning at order 32.

\section{Aggregate relaxations}
\label{p10-sec:aggregate}
At lag $d$ the Costas condition asks for $n-d$ distinct values in a range of $2n-2$: injectivity at density
at most one half. Relaxations that keep only aggregates of the differences, such as sums, sums of squares
or residues, lose which values are occupied, which is where the constraint bites. A direct test: for each
lag, the sum of squares of the differences must lie between the least and greatest values achievable by
$n-d$ distinct allowed differences with the same sum. At order 32, $95.9\%$ of 20{,}000 random permutations,
none of them Costas, satisfy all 31 of these conditions, as do $N_{32}$ and the witness of Paper~7; at
order 12 the share is $62.3\%$, so the conditions weaken as the order grows. Correll and Swanson make the
corresponding observation for a single statistic: the number of positive first differences, which they
show lies within 2.5 of $(n-1)/2$ for every known Costas array, by itself yields only an $O(n^{-1/2})$
density~\cite{p10-correll2023difference}, against the exponential decay of the true
density~\cite{p10-warnke2023density}.

Classifying midpoint states by short-lag signatures falls to the same effect. With a signature recording,
for lags 1 to 4, how many differences are positive and how many are small, signature classes with at least
100 members and no completable state cover $9.35\%$ of the 987{,}664 locally valid order-12 midpoint states,
as the independent implementation found and a recomputation here reproduces exactly. With the completable
labels shuffled at random, such classes cover $8.3$--$9.0\%$: almost all of the apparent exclusion is what
sparsity alone produces.

\section{Implications}
\label{p10-sec:implications}
Three things should be distinguished.
\begin{itemize}
\item \emph{An obstruction language was found.} The pair-marginal relaxation proves dead prefixes infeasible
with linear certificates that no local rule and no resolution-based solver produces, and it does so
completely, in sample, beyond about half depth. That is new, and it is the first strong global obstruction
system for partial Costas permutations that we know of.
\item \emph{A cheap propagator was not.} Every reasoning tested below the transition, whether local
rules, clause learning, cutting planes or triple consistency, either fails to refute early or needs effort
comparable to the native search. The mechanisms that work on exact occupancy, namely the pair-marginal
relaxation, binary displacement consistency and chain-filtered Hall analysis, all become strong in the
same band near half depth, and aggregate statistics carry little beyond sparsity.
\item \emph{Order 32 is not made feasible.} The transition falls where the order-32 tree is widest, which caps
any rule acting there at about a nine-fold saving, or 28-fold for the triple level.
\end{itemize}
Two questions remain open. Why the transition sits near half depth is unexplained; the number of pairs
between placed and unplaced columns, $k(n-k)$, peaks at $k=n/2$, which suggests a mechanism but proves
nothing. And the certificates themselves, reconstructed exactly from dual multipliers, have not yet been
read for a common form.


\section{Status of the claims}\label{p10-sec:status}
The soundness of every relaxation is immediate: each constraint holds for every Costas array. All
other results are computational and sampled, and are stated with their sample sizes. The pair-marginal
transition was reproduced by an independent implementation, which also ran the triple-consistency test;
its archives are kept with the project's data in \texttt{mining/}. Linear certificates were not extracted
exactly, and the cutting-planes proofs were not logged or checked.

\paragraph{Relation to prior work.} The relaxation hierarchy is classical~\cite{p10-sherali1990}, and the solvers
are standard. Its application to partial Costas arrays, the transition near half depth, and the
measurements of obstruction effort are, as far as we know, new; a systematic literature search has not been
completed. Sections~\ref{p10-sec:late} and~\ref{p10-sec:aggregate} include results of an independent implementation within the project; one of them, the midpoint signature classes, is reinterpreted here by a shuffled-label control. The orders 18 and 20 measurements are from the same implementation; their data are in \texttt{mining/independent/orders18\_20/}.
\begin{paperrefs}
\bibitem{p10-elffers2018} J.~Elffers and J.~Nordstr\"om, \emph{Divide and conquer: towards faster pseudo-Boolean solving}, Proceedings of the 27th International Joint Conference on Artificial Intelligence (IJCAI), 2018.
\bibitem{p10-sherali1990} H.~D.~Sherali and W.~P.~Adams, \emph{A hierarchy of relaxations between the continuous and convex hull representations for zero-one programming problems}, SIAM Journal on Discrete Mathematics 3(3) (1990), 411--430.
\bibitem{p10-warnke2023density} L.~Warnke, B.~Correll~Jr. and C.~N.~Swanson, \emph{The density of Costas arrays decays exponentially}, IEEE Transactions on Information Theory 69(1) (2023), 575--581.
\bibitem{p10-correll2023difference} B.~Correll~Jr. and C.~N.~Swanson, \emph{Difference-based structural properties of Costas arrays}, Designs, Codes and Cryptography 91(3) (2023), 779--794.
\end{paperrefs}


\chapter*{Conclusions of the investigation}
\addcontentsline{toc}{chapter}{Conclusions of the investigation}

\paragraph{What has been excluded.} Any order-32 Costas array must be geometrically asymmetric
(known since 1993, and independently re-determined in Paper~8 both by SAT and by an exhaustive direct
search); must not arise from the unmodified Welch
or Golomb constructions, and none of the construction families examined in Paper~1 reaches order
32; must not be circular; must have no Costas deletion minor of order 28 or 29 (Paper~3, relative to
the published censuses); and must not be a two-point augmentation of any of the 664 generated
order-30 cores (Paper~2). These conditions are necessary, not sufficient, and together they do not
define a tractable residual enumeration.

\paragraph{What remains possible.} An order-32 Costas array with none of these structures. Nothing in
this volume bounds that class from above, and exhaustive enumeration of all order-32 permutations
is out of reach: Paper~9 measures it at about $1.9\times10^5$ core-years for the best search engineered here, about 18
years on $10^4$ cores, and
estimates $10^4$ for the fastest published method. Paper~6's conditional model suggests the class is very probably empty, but it is a
heuristic, and an equally good fit reverses it.

\paragraph{What is proved.} The deductive results are the soundness of the telescoping rule
(Paper~5), the diagonal parity floor (Paper~7), the classification of possible stabilisers, the
fixed-point partition and the exactness of the regional encoding (Paper~8), and the elementary
results of Paper~1. Everything else is computational: complete and, where stated, audited or
independently re-scored, but not certified. Only Paper~4's representatives are machine-checked.

\paragraph{Closed directions.} Several plausible routes were tested and closed, each for a stated reason:
the telescoping propagator as a speed-up (Paper~5), the dyadic filtration and, with an explicit witness,
unlabelled dyadic Fourier summaries including the M10 transform (Paper~7), SAT as a cheaper route to
the symmetric case (Paper~8), and the search for obstructions shorter than search, through local rules,
clause learning, cutting planes and a hierarchy of linear relaxations (Paper~10). The last found a genuine
obstruction language, the pair-marginal relaxation, but one that becomes complete only near half depth,
where the order-32 tree is widest.

\paragraph{What is worth continuing.} In order of value for cost:
\begin{enumerate}[nosep]
\item a systematic literature search for every strand, before publication, given that Paper~8's
result had been published in 1993;
\item reproducing the 2011 enumeration method and measuring its tree and rate at orders 20 to 24, which
would turn Paper~9's estimate for the general search, about $10^5$ KSU, into a calibrated figure;
\item complete bibliographic data, a persistent identifier for the archive, full hashes and a
volume-wide claim manifest;
\item an aggregate build of Paper~4's Lean development, a compact certificate for Paper~7's
union-find search, and proof certificates for representative regions of Paper~8.
\end{enumerate}
The investigation's broader result is a map: where structural reasoning reaches at order 32, where
it stops, and which apparently promising approaches do not provide a practical or logically
sufficient attack.

\end{document}
